% 本文件由 LaTeX 排版 Agent 根据有效 translation.json 自动生成。
% author=Aphrahat/Aphraates (c. 280-367); work_id=3701; title=论证集

\chapter{论证集}

% structure: p0001 source=h1 target=omitted reason=page-title-already-rendered-by-parent

\Emph{论证一} 论信德\newline{}\Emph{论证五} 论战争\newline{}\Emph{论证六} 论隐修士\newline{}\Emph{论证八} 论死者复活\newline{}\Emph{论证十} 论司牧\newline{}\Emph{论证十七} 论基督天主之子\newline{}\Emph{论证二十一} 论迫害\newline{}\Emph{论证二十二} 论死亡与末世\par

\SourceRecord{Translated by John Gwynn.；Nicene and Post-Nicene Fathers, Second Series；Vol. 13.；Edited by Philip Schaff and Henry Wace.；Buffalo, NY: Christian Literature Publishing Co.,；1890.；<http://www.newadvent.org/fathers/3701.htm>.}

% structure: p0001 source=h1 target=section reason=page-title

\section{论证一（论信德）}

% structure: p0002 source=h2 target=subsection reason=relative-heading-rank; base=h2

\subsection{求问者之信函}

1. 挚爱的，我向你寄送询问及问题，因为我不得不在许多事上向你寻求更进一步的教导。请不要拒绝听我。我的心神催迫我就许多题目警告你，好叫你为我揭示你心灵的属灵领悟，并将你从圣经\OriginalTerm{圣经}{Scripture}中所领悟的一切指示给我，这样，我的缺乏可由你补足，我的饥渴可由你的教导满足，你可从你教导的泉源解我的干渴。然而，我心中虽有许多事要问你，它们却仍都留在我这里，待我来到你那里时，你可在一切事上教导我。\par

2. 但在一切事之前，我渴望你写信教导我这件困扰我的事，即：关于我们的信德\OriginalTerm{信德}{faith}；它是怎样的，它的基础是什么，它建基于什么结构之上，它依靠什么，它的完满与成就以何种方式实现，以及为它所需要的是哪些善工。因为我自己坚信天主\OriginalTerm{天主}{God}是唯一的，祂起初创造了天地；祂以自己的工程装饰了世界；祂按照自己的肖像造了人；就是祂悦纳了\Person{亚伯尔}的祭献；祂因\Person{哈诺客}的卓越而将他提去；祂因\Person{诺厄}的义德而保全了他；祂因\Person{亚巴郎}的信德而拣选了他；祂因\Person{梅瑟}的温良而与他说话；就是祂在所有先知内说话，并且更进一步，祂派遣了祂的基督\OriginalTerm{基督}{Christ}进入世界。我的弟兄，既然我如此相信这些事是如此，因此，弟兄，我请求你写信指示我，为这我们的信德所需要的是哪些善工，好叫你使我安心。\par

% structure: p0005 source=h2 target=subsection reason=relative-heading-rank; base=h2

\subsection{论证开始}

1. 我收到了你的信，我的挚爱者，当我读它时，它大大欢悦了我，因为你将你的思想转向了这些探究。因为你们所请求于我的这事，将白白赐予，\BibleRef{玛 10:8} 因为这是白白领受的。凡有的，却愿意对那寻求者有所保留，他所保留的必被夺去。凡因白白的恩宠\OriginalTerm{恩宠}{grace}领受的，也该白白的施予。所以，我的挚爱者，关于你所询问于我的，就我微不足道所领悟的，我将写信给你。而且，凡你所没有寻求于我的，我呼求天主，也必向你解释。那么，我的挚爱者，请听，并向我开启你内心的耳，和你心灵的属灵知觉，来听我对你所说的话。\par

2. 信德由许多事物组合而成，并藉许多种类达至成全。因为它好像一座建筑，由许多工艺的构件建造起来，这样它的建筑物便升到顶端。我的挚爱者，你要知道，在建筑物的根基上安放石头，如此，整个建筑物靠着石头升起，直至完竣。同样，真磐石\OriginalTerm{磐石}{Stone}，我们的主耶稣基督，是我们一切信德的基础。而信德便建于祂，[这]磐石之上。又依靠信德，整个结构升起，直至完成。因为根基是整个建筑的开端。因为当任何一个人被带领接近信德时，它便为他奠定在那磐石，即我们的主耶稣基督之上。他的建筑不能被波浪摇撼，也不能被暴风伤害。在狂风骤雨中它不倒塌，因为它的结构是被建造在真磐石的岩石上。而在我称基督为磐石这事上，我不是出于我自己的思想，而是先知们预先称祂为磐石。这一点我必给你讲明。\par

3. 现在，请听关于建于磐石之上的信德，及关于那建在磐石上的结构。因为人先相信；当他相信时，他便爱慕；当他爱慕时，他便希望。当他希望时，他便成义\OriginalTerm{成义}{justification}。当他成义时，他便成全。当他成全时，他便完满。而当他的整个结构被建立起来，完满且成全时，他便成为房屋和殿宇，作\Person{基督}的居所，正如\Person{耶肋米亚}先知说过：\BibleQuote{\Emph{主的殿，主的殿，主的殿是你们，如果你们改善你们的行径和作为}}。\BibleRef{耶 7:4} 祂又藉先知说过：\BibleQuote{\Emph{我要住在他们中间，并在他们中间行走}}。\BibleRef{肋 26:12} 真福宗徒也这样说：\BibleQuote{\Emph{你们是天主的殿，且基督的神住在你们内}}。而我们的主又这样对祂的门徒说过：\BibleQuote{\Emph{你们在我内，我在你们内}}。\BibleRef{若 14:20}\par

4. 当房屋成为居所时，人便开始为那位居住在此建筑内的居住者所需的事物而焦虑。比如，如果一位国王或尊贵的人——被冠以君王之名——下榻于这房屋，那么就需要为这位国王预备一切君王应有的排场以及为国王的尊荣所需的一切侍奉。因为在一间缺乏一切美好事物的房屋中，国王是不会下榻的，也不会居住于其中；反之，房屋中一切最上等的都必须为国王预备，且不得有任何欠缺。如果国王下榻的房屋中有所缺欠，那房屋的看守人就要被处死，因为他没有为国王预备好侍奉。同样地，那成为基督的房屋、更确切地说，成为基督的居所的人，务要留心那些为居住在他内的基督所需用的侍奉，以及他应当用哪些事物来取悦基督。因为他首先将自己的建筑建造在那\OriginalTerm{磐石}{the Stone}上，那磐石就是基督。\OriginalTerm{信德}{faith}是建基在祂——这块磐石——之上，而所有的建筑都是在信德上被建造起来的。因为要作这房屋的居所，就需要纯洁的斋戒，而斋戒因着信德得以坚固。同时又需要纯洁的祈祷，祈祷因着信德而蒙悦纳。对它而言，爱德也是必须的，爱德与信德融合在一起。此外，施舍也是需要的，施舍因着信德而被施予。祂也要求温良，温良因着信德而得装饰。祂也挑选童贞之德，童贞因着信德而被爱慕。祂又联同圣洁，圣洁在信德中被栽植。祂也关心智慧，智慧因着信德而被获取。祂还盼望款待之德，款待因着信德而丰盛。祂又需要纯朴，纯朴与信德交融。祂也要求忍耐，忍耐因着信德而得以成全。祂又看重坚忍，坚忍因着信德而获得。祂也喜爱哀恸，哀恸因着信德而彰显。祂又寻求洁净，洁净因着信德而得以保存。这一切德性，都是那建基于真磐石——即基督——的岩石上的信德所要求的。这些善工都是为君王基督所要求的，祂居住在那些由这些善工所建造的人内。\par

5. 你或许会说：如果基督被立为基础，那么当建筑完成时，基督又怎样也居住于其中呢？因为真福宗徒提到了这两件事。他这样说：\BibleQuote{我作为明智的建筑师立了基础}。\BibleRef{格前 3:10} 在那里他定义了基础并说明了它，因为他这样说：\BibleQuote{任何人不能再奠立别的根基，除了那已奠立的，就是耶稣基督}。\BibleRef{格前 3:11} 而基督又居住在那建筑里，正是上面所写的话——就是\Person{耶肋米亚}称人为圣殿，并说天主居住在他们内的那段话。宗徒又说：\BibleQuote{基督的圣神住在你们内}。\BibleRef{格前 3:16} 我们的主也说：\BibleQuote{我与父原是一体}。\BibleRef{若 10:30} 因此那句话得以应验，即基督住在人内，就是住在那些信祂的人内，而祂就是那基础，整个建筑在其上被建造起来。\par

6. 但我必须回到我先前的陈述，即先知们称基督为磐石。因为古时\Person{达味}论到祂说：\BibleQuote{建筑者所弃的石头，已成了角石的基顶}。建筑者如何拒绝了这块磐石——也就是基督？\BibleRef{路 19:14} 无非是他们在\Person{比拉多}面前那样拒绝祂，说：\BibleQuote{这个人不可作我们的君王}。\BibleRef{若 19:15} 此外，在我们的主所说的比喻中，有个贵族去接受王权并回来统治他们；他们却在他后派去使者说：\BibleQuote{这个人不可作我们的君王}。\BibleRef{路 19:13} 藉着这些事，他们拒绝了那块磐石，即基督。而它怎样成了建筑的角石呢？无非是它被立在外邦人的建筑之上，他们所有的建筑都建基于其上。谁是建筑者呢？岂不是那些司祭和法利塞人吗？他们并没有建造坚固的建筑，反而推翻他正在建造的一切，正如\Person{厄则克耳}先知所记载的：\BibleQuote{他筑起一道墙壁，他们却粉刷它，使它倒塌}。\BibleRef{则 13:10} 经上又记载：\BibleQuote{我在他们中寻找一个重修围墙、在我面前为这地方站在缺口处的人，免得我破坏那地方，但我没有找到}。\BibleRef{则 22:30} 此外，\Person{依撒意亚}也关于这块石头说预言。他说：\BibleQuote{吾主上主这样说：看，我要在\Place{熙雍}安放一块精选的石头，一块宝贵的角石作基础，是那稳固的基础角石}。\BibleRef{依 28:16} 他又在那里说：\BibleQuote{凡信赖它的，必不恐惧}。\BibleRef{依 28:16} 并且\BibleQuote{谁若跌在这石头上，必被砸碎；这石头落在谁身上，必要把他压垮}。\BibleRef{玛 21:44} 因为以色列家的人民跌倒在祂上面，祂便成了他们永远的毁灭。又说\BibleQuote{这石头要打在大像上，把它压碎}。\BibleRef{达 2:34} 而外邦人却信赖了它，并不恐惧。\par

7. 关于那块石头，祂这样指示：它既是墙的角石，又是基础。但如果那块石头被立为基础，它又怎样成了墙的角石呢？无非是我们的主来临时，祂把自己的信德放在地上如同基础，而它升起高过诸天，如同墙的角石，整个建筑便由这些石头，从底部到顶部，完成了。至于我所说的祂把自己的信德放在地上的信德，\Person{达味}论到基督预先宣告了。因为他说：\BibleQuote{信德将从地上萌生}。而关于那在上面的，他又说：\BibleQuote{正义由天上俯视}。\par

8. 还有\Person{达尼尔}也论及这基石——即基督。因为他说道：— \BibleQuote{\Emph{那块石头，未经人手开凿，从山峦上滚下来，击中那像，然后整个大地都充满了它。}}\BibleRef{达 2:34} 这事他预先指明关于基督：整个大地将要充满祂。看啊！借着基督的信德，大地的四极都充满了，正如\Person{达味}所说的：— \Quote{\Emph{基督福音的声响，已传遍了整个大地。}} 还有，当祂派遣祂的宗徒们时，祂对他们这样说：— \BibleQuote{\Emph{你们去，使万民成为门徒，他们将会相信我。}}\BibleRef{玛 28:19} 还有，\Person{匝加利亚}先知也预言了关于那基石——即基督。因为他说道：— \Quote{\Emph{我看见一块首要的、合一的、爱的石头。}} 但他为什么说 \InnerQuote{\Emph{首要}} 呢？无疑是因为自太初祂就与祂的圣父同在。再者，他谈到爱，是因为当祂来到世上时，祂对祂的门徒们这样说：— \BibleQuote{\Emph{这是我的命令：你们要彼此相爱。}}\BibleRef{若 15:12} 祂又说：— \BibleQuote{\Emph{我称你们为我的朋友（所爱的）。}}\BibleRef{若 15:15} 而蒙福的宗徒这样说：— \Quote{\Emph{天主如在祂圣子的爱中爱了。}} 实在地，\BibleQuote{\Emph{基督爱了我们，并为我们牺牲了自己。}}\BibleRef{弗 5:2}\par

9. 祂肯定地论及这石头说：— \BibleQuote{\Emph{看！我要在这石头上开启七只眼睛。}}\BibleRef{匝 3:9} 那么，开启在石头上的七只眼睛是什么呢？显然是那居住在基督内的\OriginalTerm{天主的神}{Spirit of God}，有七种恩宠的运作，正如\Person{依撒意亚}先知所说的：— \BibleQuote{\Emph{天主的神将安息并停留在祂身上}，（一位灵）\Emph{智慧和聪敏的灵，超见和刚毅的灵，明达和敬畏主的灵。}}\BibleRef{依 11:1} 这就是开启在石头上的七只眼睛，而\BibleQuote{\Emph{这七只眼睛就是主的眼睛，遍察全地。}}\BibleRef{匝 4:10}\par

10. 同样，以下所述也是指着基督说的。他说祂被赐予，作为万邦的光明，正如\Person{依撒意亚}先知所说的：— \BibleQuote{\Emph{我已立祢作万邦的光明，为使祢成为我的救恩，直到大地的尽头。}}\BibleRef{依 49:6} 此外，\Person{达味}也说：— \Quote{\Emph{祢的言语是我脚步前的明灯，是我路途上的亮光。}} 而主的言语和\OriginalTerm{圣言}{Word}就是基督，正如我们救主的福音一开头所写：— \BibleQuote{\Emph{在起初已有圣言。}}\BibleRef{若 1:1} 关于光，他又作证说：— \BibleQuote{\Emph{光在黑暗中照耀，黑暗却未能胜过。}}\BibleRef{若 1:5} 那么，\Emph{光在黑暗中照耀，黑暗却未能胜过}是什么意思呢？显然是基督，祂的光在以色列家族百姓中间照耀，而以色列家族百姓未能领悟基督的光，因为他们不信祂，正如经上所载：— \BibleQuote{\Emph{祂来到自己的领域，自己的人却没有接受祂。}}\BibleRef{若 1:11} 此外，我们的主耶稣称他们为黑暗，因为祂对门徒们说：— \BibleQuote{\Emph{我在暗中告诉你们的，你们要在光中说出来，}}\BibleRef{玛 10:27} 也就是说，\BibleQuote{\Emph{你们的光当在外邦人中照耀；}}\BibleRef{玛 5:16} 因为他们接受了基督的光，祂是万邦的光明。祂又对宗徒们说：— \BibleQuote{\Emph{你们是世界的光。}}\BibleRef{玛 5:14} 祂又对他们说：— \BibleQuote{\Emph{你们的光当在人前照耀，好使他们看见你们的善行，而光荣你们在天之父。}}\BibleRef{玛 5:16} 祂又表明祂自己是光，因为祂对门徒们说：— \BibleQuote{\Emph{当光还在你们中时，你们要行走，免得黑暗笼罩你们。}}\BibleRef{若 12:35} 祂又对他们说：— \BibleQuote{\Emph{你们要信从光，好使你们成为光明之子。}}\BibleRef{若 12:36} 祂又说：— \BibleQuote{\Emph{我是世界的光。}}\BibleRef{若 8:12} 祂又说：— \Quote{\Emph{没有人点了灯，放在斗底下或床底下，或放在隐密处，而是放在灯台上，为使所有人看见灯光。}} 而那发光的灯就是基督，正如\Person{达味}所说的：— \Quote{\Emph{祢的言语是我脚步前的明灯，是我路途上的亮光。}}\par

11. 此外，\Person{欧瑟亚}先知也说过：——\BibleQuote{点起灯来寻求上主}。\BibleRef{欧 10:12} 我们的主耶稣基督说过：——\BibleQuote{哪有一个妇女，有十个\OriginalTerm{达玛}{drachmos}，若遗失了一个，哪有不点灯打扫房屋，细心寻找，直到找着的呢？}\BibleRef{路 15:8} 那么，这个妇女预表什么呢？显然预表以色列家的会众，就是领受了十诫的那会众。他们失落了第一条诫命——就是祂警告他们说：——\BibleQuote{我是上主你的天主，曾将你从埃及地领出来。}\BibleRef{出 20:2} 他们既然失落了这第一条诫命，随后的九条也就无法遵守，因为九条都倚赖第一条。因为当他们崇拜\OriginalTerm{巴耳}{Baal}时，就绝不可能遵守那九条诫命。因为他们失落了第一条诫命，正像那妇女失落了十个达玛中的一个。于是先知向他们呼喊：——\BibleQuote{点起灯来寻求上主}。\BibleRef{欧 10:12} 此外，\Person{依撒意亚}先知也说过：——\BibleQuote{趁上主可找到的时候，你们应寻找祂；趁祂在近处的时候，你们应呼求祂。罪人应离开自己的行径，恶人该抛弃自己的思念。}\BibleRef{依 55:6-7} 那灯已经照耀，他们却没有借着它寻求上主他们的天主。光在黑暗中照耀，黑暗却没有领悟它。灯被放在灯台上，屋里的人却没有看见它的光。那么，灯被放在灯台上，这又意味着什么呢？显然是指祂被高举在十字架上。因此，整个房屋为他们变成了黑暗。因为他们钉死祂的时候，光就离开他们而变暗，却照耀在\OriginalTerm{外邦人}{Gentiles}中；因为从他们钉死祂的第六时辰（日间）直到第九时辰，以色列全地都变为黑暗。太阳在正午落下，大地在白日光辉里变为黑暗，正如\Person{匝加利亚}先知所记载的：——\BibleQuote{上主说：到那一天，我要使太阳在正午落下，使大地在白日光辉里变为黑暗。}\BibleRef{匝 14:6-7}\par

12. 现在我必须回到我先前的主题——信德，因为一切善工的建筑都建立在它之上。再者，关于建筑，我所说的并非怪异之谈，而是真福宗徒在\BibleBook{格林多}{前书}中所写的，他说：——\BibleQuote{我好像一个精明的建筑师，奠定了根基，别人在上面建筑。有的用金、银、宝石建筑；有的用芦苇、禾秸、残茬建筑。在最后的日子，那建筑要被火试验；因为金、银、宝石的工程在火中得以保存，因为是坚固的建筑。至于芦苇、禾秸、残茬的工程，火必对它有权势，它必被烧毁。}\BibleRef{格前 3:10-15} 那么，那建立建筑的金、银、宝石是什么呢？显然是指信德的善工，它们必在火中得以保存；因为基督居住在那坚固的建筑内，祂就是保护它免受火烧的那位。让我们从天主在先前计划中所赐的榜样来思考并领悟，因为那计划的许诺必为我们恒久确定。那么，让我们从那三个被投入火中却没有被烧毁的义人身上领悟，就是\Person{哈纳尼雅}、\Person{阿匝黎雅}和\Person{米沙耳}，火对他们毫无权势，因为他们建造了坚固的建筑，拒绝了\Person{拿步高}王的命令，没有朝拜他所造的像。至于那些违背天主诫命的人，火立刻就胜过他们，烧毁了他們，毫无怜悯地烧了他们。因为\OriginalTerm{索多玛人}{Sodomites}就像禾秸、芦苇、残茬一样被烧毁。再者，\Person{纳达布}和\Person{阿彼胡}因违背天主的诫命被烧死。还有那献香的二百五十人被烧死。再者，那两个官长和随同他们的一百人，因为走近\Person{厄里亚}所坐的山，被火烧死；厄里亚后来乘火马车升了天。那些诽谤者因给义人挖了陷阱，也被烧死了。因此，亲爱的诸位，义人必如金、银、宝石般被火试验，而恶人必如禾秸、芦苇、残茬般在火中被烧，火必对他们有权势，他们必被烧毁；正如\Person{依撒意亚}先知说：——\BibleQuote{上主要以火施行审判，以火考验一切血肉。}\BibleRef{依 66:16} 他又说：——\BibleQuote{你们要出去观看那些违抗我的人尸骸；他们的虫不死，他们的火不灭；他们为一切血肉成为惊骇。}\BibleRef{依 66:24}\par

13. 宗徒又为我们诠释了这建筑和这根基；因为他这样说：——\BibleQuote{除已奠定的根基，即耶稣基督外，谁也不能再奠立别的根基。}\BibleRef{格前 3:11} 宗徒又说，信德与望德及爱德相联，因为他这样说：——\BibleQuote{现今存在的，有信德、望德、爱德这三样。}\BibleRef{格前 13:13} 他也表明了，信德首先被奠立在坚固的根基上。\par

14. 论及\Person{亚伯尔}，因他的信德，他的祭献蒙了悦纳。\Person{哈诺客}，因他藉信德中悦天主，被免于死亡。\Person{诺厄}，因他相信，从洪水中得以保全。\Person{亚巴郎}，藉他的信德，获得了祝福，这就算为他的正义。\Person{依撒格}，因他相信，得蒙喜爱。\Person{雅各伯}，因他的信德，得以保全。\Person{若瑟}，因他的信德，在争议之水中受了试探，并从他的试炼中蒙了拯救，他的主在他内立了见证，如\Person{达味}所说：——\BibleQuote{他在若瑟内立了见证}。\Person{梅瑟}也藉他的信德行了许多异能奇事。他藉信德以十灾毁灭了\Place{埃及}人。又因信德分开大海，使自己的百姓渡过，而将埃及人淹没海中。他藉信德将木头投入苦水，水就变甜了。他藉信德降下\OriginalTerm{玛纳}{manna}，使他的百姓饱足。他藉信德伸开双手，战胜了\Person{阿玛肋克}，正如经上所载：——\BibleQuote{他的双手在信德中恒举不落，直到日落}。又因信德他上了\Place{西乃山}，在那里两次禁食四十天。再因信德他战胜了\Person{阿摩黎}人的王\Person{息红}和\Person{敖格}。\par

15. 亲爱的诸位，这是奇妙的，\Person{梅瑟}在\Place{红海}所行的伟大奇迹，当众水因信德分开，涌起如高山或峻峭的悬崖。它们受命止住，僵立不动；它们如被关闭在器皿中，在高处亦如深处被紧紧束缚。它们的流动性并未溢出界限，反而改变了其受造的本性。无灵之物变得顺从。巨浪凝固，等待着报复，等待着百姓过去。波浪如何停立，静候命令与报复，这真是奇妙。自世世代以来隐藏的根基显露了，那自起初为液体的，骤然干涸。\BibleQuote{城门昂起头来，古老的门户高抬}。火柱进入，照亮整个营幕。百姓因信德而过。正义的审判遂临于\Person{法郎}、他的军旅和战车之上。\par

16. 同样，\Person{农}的儿子\Person{若苏厄}因信德分开了\Place{约旦河}，\OriginalTerm{以色列子民}{children of Israel}过了河，如同在\Person{梅瑟}的日子。但亲爱的诸位，要知道，这约旦河的通道曾三次因分开而显露。第一次藉\Person{农}的儿子\Person{若苏厄}，第二次藉\Person{厄里亚}，然后藉\Person{厄里叟}。正如\WorkTitle{书}上的话所表明的，对着\Place{耶里哥}的这通道处，厄里亚在那里被提升天；因为当厄里叟从他跟随转身，分开约旦河过去之后，\Place{耶里哥}的先知弟子们出来迎接厄里叟说：——\BibleQuote{厄里亚的神魂临到厄里叟身上了}。\BibleRef{列下2:8} 此外，当百姓在\Person{农}的儿子\Person{若苏厄}的时代过河时（正是在那里），因为经上这样记载：——\BibleQuote{百姓正对着耶里哥过了河}。\BibleRef{苏3:17} \Person{农}的儿子\Person{若苏厄}也因信德使\Place{耶里哥}的城墙倒塌，毫不费力就倒了。又因信德他击败了三十一个君王，使以色列子民继承了土地。再者，他藉信德向天伸开双手，使太阳停在\Place{基贝红}，月亮停在\Place{阿雅隆}山谷。\BibleRef{苏10:13} 它们停住了，从它们的轨道上静止不动。但够了！所有义人，我们的祖先，在他们所做的一切事上，都是凭信德而获胜，正如\OriginalTerm{真福宗徒}{the blessed Apostle}论及他们众人所证实的：——\BibleQuote{他们凭信德得了胜利}。\BibleRef{希11:33} \Person{撒罗满}也说：——\BibleQuote{仁慈的人虽多，但忠信的人谁能找到}？\BibleRef{箴20:6} \Person{约伯}也这样说：——\BibleQuote{我的正义我决不离弃，我要坚持我的清白}。\BibleRef{约27:5}\par

17. 同样，我们的救主习惯对每个前来求助祂医治的人说：\BibleQuote{照你的信德，给你成就吧！}当那盲人走近祂时，祂对他说：\BibleQuote{你信我能治好你吗？}那盲人对祂说：\BibleQuote{主，我信。}\BibleRef{玛 9:28}他的信德使他眼睛开了。对那个儿子患病的人，祂说：\BibleQuote{你信吧，你的儿子必活。}他对祂说：\BibleQuote{我信，主，请补助我的薄弱信德。}因他的信德，他的儿子便痊愈了。又有一个王臣来到祂跟前，因他的信德，他的孩子得了医治，当时他对我们的主说：\BibleQuote{只要你说一句话，我的仆人就会痊愈。}\BibleRef{玛 8:8}我们的主惊奇他的信德，照他的信德，给他成就了。又有一个会堂长，为他的女儿来求祂，祂这样对他说：\BibleQuote{只要坚信，你的女儿必活。}\BibleRef{谷 5:23}他就信了，他的女儿便活了，起来了。当拉匝禄死去时，我们的主对玛尔大说：\BibleQuote{如果你信，你的兄弟必要复活。}玛尔大对祂说：\BibleQuote{主，我信。}\BibleRef{若 11:23}祂便在四天后复活了他。还有西满，因他的信德被称为刻法，即磐石。再者，当我们的主将圣洗圣事交给宗徒们时，祂这样对他们说：\BibleQuote{信而受洗的必要得救；不信的必被判罪。}祂又对宗徒们说：\BibleQuote{如果你们有信德，不疑惑，你们没有不能做的事。}\BibleRef{玛 21:22}当我们的主在海面上行走时，西满也凭信德在水面上与祂同行；但当他因信德动摇而疑惑，开始下沉时，我们的主称他为\BibleQuote{小信德的人}。\BibleRef{玛 14:31}当宗徒们请求我们的主时，他们只向祂求了这件事，对祂说：\BibleQuote{增加我们的信德吧！}祂对他们说：\BibleQuote{如果你们有信德，就连一座山也会从你们面前移开。}祂又对他们说：\BibleQuote{不要怀疑，免得你们沉沦于世界，如同西满怀疑时便开始沉于海中。}祂又这样说：\BibleQuote{这些奇迹将伴随信的人：他们要说新语言，驱除魔鬼，按手在病人身上，病人就会痊愈。}\BibleRef{谷 16:17}\par

18. 所以，我亲爱的，让我们亲近信德吧，因为它的能力如此浩大。因为信德将\Person{哈诺客}提升到天上，并战胜了洪水。它使不生育的生育，从刀剑下解救，从深坑中提升，使贫穷的富足，释放了被掳的，拯救了受迫害的，使火从天降下，分开海洋，劈开磐石，给口渴的人水喝，饱饫了饥饿的，复活了死人，将他们从\OriginalTerm{阴府}{Sheol}中领出，平息了波浪，治愈了病人，征服了敌军，推倒了城墙，堵住了狮子的口，熄灭了烈火的火焰，羞辱了骄傲者，举扬了谦卑的人。这一切大能的作为，都由信德完成。\par

19. 信德是这样的：即人相信天主，万有的主，祂创造了天地海洋及其中所有的一切；祂以自己的肖像造了\Person{亚当}；祂将法律赐给了\Person{梅瑟}；祂将祂的圣神派遣到先知们身上；祂更派遣了祂的基督到世界上来。再者，人应信死者的复活；并且还应信圣洗圣事。这就是天主的教会的信德。并且（必须）人要远避遵守时辰、安息日、月朔和节期，远避占卜、巫术、\OriginalTerm{迦勒底法术}{Chaldæan arts}和邪术，远避淫乱和宴乐，远避虚妄的道理——这些都是恶者的工具——远避甜言蜜语的谄媚，远避亵渎和奸淫。人不可作假见证，也不可说一口两舌的话。这些就是信德的行为，这信德建立在真磐石——基督之上，整个建筑都立在祂上面。\par

20. 此外，我亲爱的，\WorkTitle{圣经}中还有许多关于信德的内容。但我从这些丰富的内容中，写下这少许，为唤起你们的爱心，叫你们能认识并传扬，且能相信，也使人相信。当你们读了并学习了信德的行为，你们便可能会像那被耕好的田地，好种子落在其上，结出百倍的、六十倍的、三十倍的果实。当你们来到主面前时，祂必称你们为善良忠信的仆人，这仆人因他丰厚的信德，得以进入他主人的国。\par

\SourceRecord{Translated by John Gwynn.；Nicene and Post-Nicene Fathers, Second Series；Vol. 13.；Edited by Philip Schaff and Henry Wace.；Buffalo, NY: Christian Literature Publishing Co.,；1890.；<http://www.newadvent.org/fathers/370101.htm>.}

% structure: p0001 source=h1 target=section reason=page-title

\section{证道五（论战争）}

\Emph{1.} 当前关于此时将发生的震动，以及那为刀剑而聚集的军队，我心中浮想联翩。时间是由天主预先安排的。和平的时日完成于善良与义人的日子；诸多灾祸的时日完成于邪恶与罪人的日子。因为经上这样记载说：\BibleQuote{\Emph{善必须发生，使之实现的人是有福的；}}并且\BibleQuote{\Emph{恶必须发生，但使之实现的人是有祸的。}}善已经临于天主的子民，而有福分等待着那带来善的人。至于那邪恶骄傲的人，恶被煽动起来，因着那夸耀的邪恶傲者所集结的军队；也有祸患为他存留，因他煽动了恶。然而，我亲爱的，不要责备那将恶加于多人的恶人；因为时间早已被预先安排，它们应验的时辰已经到了。\par

\Emph{2.} 因此因为这是\OriginalTerm{那恶者}{the Evil One}的时刻，请你们在奥迹中聆听我为你们所写的。因为经上这样记载说：\BibleQuote{\Emph{凡在人前自高的，在天主前是可憎的。}}\BibleRef{路 16:15} 又记载说：\BibleQuote{\Emph{凡高举自己的，必被贬抑；凡贬抑自己的，必被高举。}}\BibleRef{路 14:11} 耶肋米亚也说过：\BibleQuote{\Emph{勇士不可夸耀自己的勇力，富人不可夸耀自己的财富。}}\BibleRef{耶 9:23} 再有蒙福的宗徒说过：\BibleQuote{\Emph{凡要夸耀的，应当因主而夸耀。}}\BibleRef{格后 10:17} 达味曾说：\BibleQuote{\Emph{我看见恶人骄横高耸，有如黎巴嫩的香柏；但我再经过时，他已不在；我寻找他，却未能寻见。}}\par

\Emph{3.} 因为凡夸耀的，必被贬抑。加音在对自己的弟弟亚伯尔夸耀后，将他杀害。于是他受诅咒，成为流离飘荡于地上的人。索多玛人对罗特夸耀，于是有火从天降下，烧毁了他们，他们的城也倾覆于他们身上。厄撒乌对雅各伯夸耀，并迫害他，而雅各伯却得到了厄撒乌的长子名分和祝福。雅各伯的子孙对若瑟夸耀，后来却在埃及俯伏下拜于他。法郎对梅瑟和他的百姓夸耀，于是法郎及其军队被淹没在海中。达堂和阿彼兰对梅瑟夸耀，他们便活活地陷落\OriginalTerm{阴府}{Sheol}。哥肋雅威胁达味，却倒在他面前，被彻底击败。撒乌耳再次迫害达味，他却倒在培肋舍特人的刀下。阿贝沙隆对他（达味）自高，约阿布便在战场上将他杀死。本哈达得对阿哈布夸耀，却被交在以色列手中。散乃黑黎布亵渎了希则克雅和他的天主，当一位\OriginalTerm{守卫者}{Watcher}出去，因希则克雅的祈祷和先知中最光荣的先知依撒意亚的祈祷，在营中击杀了十八万五千人时，他的军队便成了死尸。阿哈布对米该亚自高，他上去后却倒在了辣摩特基肋阿得。依则贝耳对厄里亚夸耀，狗便在依次勒耳的田地里吞食了她。哈曼对摩尔德开夸耀，他的邪恶却归到了自己头上。巴比伦人对达尼尔夸耀，将他投入狮子圈，达尼尔却得胜而出，他们反被投入圈中代替了他。巴比伦人又夸耀而控告阿纳尼雅和他的同伴，他们被投入烈火窑中；他们却得胜而出，火焰反吞灭了控告者。\par

\Emph{4.} 拿步高曾说过：\BibleQuote{\Emph{我要升到天上，将我的宝座高举于天主的星辰之上，坐在极北边界的高山上。}}\BibleRef{依 14:13} 依撒意亚论到他说：\BibleQuote{\Emph{因为你的心如此高举了你，所以你必被带到\OriginalTerm{阴府}{Sheol}，所有看见你的人必因你而惊愕。}}\BibleRef{依 14:15} 散乃黑黎布也曾这样说：\BibleQuote{\Emph{我要上升山巅，直到黎巴嫩的山脊。}} \BibleQuote{\Emph{我要掘井饮水，用我的马蹄踏干所有深河。}}因为他如此自高，依撒意亚又论到他说：\BibleQuote{\Emph{斧头岂可对用它的砍伐者自夸？锯子岂可对用它的锯木者自大？棍棒岂可举起自己攻击那挥动它的？}}\BibleRef{依 10:15} 因为你散乃黑黎布，是砍伐者手中的斧头，是锯木者手中的锯子，是挥动你为惩罚的棍棒，是责打的杖。你被派遣攻击反复无常的百姓，又被派定攻击顽固的百姓，好叫你掳掠俘虏，夺取掠物；你使他们如街上的淤泥，面对众人及所有外邦人都如此。当你做了这一切，为何对那持着你的自高，为何对那用你锯木的自夸，为何辱骂圣城？并对耶路撒冷的子孙说：\BibleQuote{\Emph{难道你的天主能从我的手中救出你们吗？}}\BibleRef{列下 18:35} 你竟敢说：\BibleQuote{\Emph{谁是那能救你们脱离我手的主？}}因此，听上主的话说：\BibleQuote{\Emph{我要在我的土地上击碎亚述人，在我的山岭上将他们践踏。}}\BibleRef{依 14:25} 当他被击碎践踏后，\BibleQuote{\Emph{童贞女，熙雍女儿，將鄙视他，耶路撒冷女儿将摇头说：\InnerQuote{你辱骂了谁，亵渎了谁？你向谁扬声？你举目向天攻击以色列的圣者，又藉你使者的手辱骂了主。如今你看，你的鼻孔已被钩子钩住，你的口唇被嚼环勒住，你满心骄傲而来，却满心破碎而归。}}}他的被杀是在他所爱之人手中；在他寄望的家中，他被倾覆，倒在他自己的神前。我亲爱的人们啊，确实理当如此：他的身体竟成为祭品与奉献，献给他所依靠的那神前，在他的庙宇里，作为他偶像的记念。\par

5. 那公绵羊又举高自大，用角向西、北、南抵触，压制了许多走兽。它们都不能在他面前站立，直到那公山羊从西方来，抵触公绵羊，打断他的双角，使公绵羊全然俯伏。但公绵羊就是\Place{玛待}和\Place{波斯}王，即\Person{达理阿}；公山羊就是\Place{马其顿}人\Person{斐理伯}的儿子\Person{亚历山大}。因为\Person{达尼尔}看见那公绵羊时，他在东方，在\Place{厄蓝}省境内，\Place{乌来}河畔的\Place{稣撒}堡垒门前。他正在向西、北、南抵触。没有走兽能在他面前站立。那公山羊从\Place{希腊}地区上来，自高自大攻击公绵羊，并且打断了他的双角，大角和小角。为什么他说打断了他的双角呢？很明显是因为他压制了公绵羊所统治的两个王国，较小的\Place{玛待}王国和较大的\Place{波斯}王国。但是，当希腊人\Person{亚历山大}的到来，他杀死了\Place{玛待}和\Place{波斯}王\Person{达理阿}。因为天使向\Person{达尼尔}解释异象时，这样对他说：—— \BibleQuote{你所看见的那只公绵羊，就是\Place{玛待}和\Place{波斯}王；那只公山羊，就是希腊王。}\BibleRef{达 8:20} 如今，从公绵羊的双角被折断时，直到现在，已有六百四十八年了。\par

6. 因此，论到那公绵羊，它的角已被折断。虽然它的角已被折断，看！它却向那第四兽自高自大，那兽是\BibleQuote{强大有力，有铁牙铜蹄，吞吃、嚼碎，把剩下的用脚践踏}\BibleRef{达 7:19}。角被折断的公绵羊啊，你要离开那兽，不要激怒它，免得它吞吃你、把你嚼成粉末。公绵羊尚且不能在那公山羊面前站立，又怎能在那可怕的兽面前站立呢？它的口\BibleQuote{说话傲慢}\BibleRef{达 7:8}，又向它所见的一切，像狮子向猎物一样俯伏。凡激怒狮子的，就成了它的掠物；凡招惹那兽的，它必吞吃他。当那兽用脚践踏人时，谁能逃脱它的蹄下呢？因为那兽必不被杀戮，直到\OriginalTerm{万古常存者}{Ancient of Days}坐在宝座上，\OriginalTerm{人子}{Son of Man}来到他面前，权柄赐给了他。那时那兽要被杀戮，它的尸首必败坏。人子的国度要建立，一个永远的国度，他的权柄世世代代不断。\par

7. 你这自高自大的人，要安静！不要自夸！因为即使你的财富使你心高气傲，也不及\Person{希则克雅}的财富丰厚；他曾进去在\Place{巴比伦}人面前夸耀它，但所有的财富都被掳去，到了\Place{巴比伦}。如果你以你的儿女为荣，他们也要被掳到你面对那兽，如同\Person{希则克雅}王的儿女被掳去，在\Place{巴比伦}王的宫中成为太监一样。如果你以自己的智慧夸口，你也比不上\Person{提洛的君主}，\Person{厄则克耳}曾谴责他说：—— \BibleQuote{你比\Person{达尼尔}更聪明吗？或是凭自己的智慧看见了隐藏的事？}\BibleRef{则 28:3} 如果你因自己的年岁众多而骄傲，你的年岁也不比\Person{提洛的君主}多，他在位的时候，犹大家中有二十二位国王，即四百四十年。这位\Place{提洛}王年岁虽多，却始终心中说：\BibleQuote{我是天主，坐在海中天主的宝座上。}\BibleRef{则 28:2} 但\Person{厄则克耳}对他说：\Quote{你是人，而不是天主。} 因为当\Person{提洛的君主} \Emph{在\OriginalTerm{火石}{stones of fire}中行走，毫无过失}时，他蒙受了怜悯。但当他的心高傲时，\BibleQuote{那遮掩的革鲁宾就消灭了他。}\BibleRef{则 28:14}\par

8. 那么，\OriginalTerm{火石}{stones of fire}是什么呢？不就是\Place{熙雍}的子女和\Place{耶路撒冷}的子女吗？因为在古时，在\Person{达味}和他的儿子\Person{撒罗满}的日子，\Person{希兰}是以色列家之友。但当他们被掳离本土时，他却幸灾乐祸，用脚踢他们，不记念\Person{达味}家的友谊。至于我所说犹大的子孙称为\OriginalTerm{火石}{stones of fire}，这不是出于我自己的意思，而是\Person{耶肋米亚}先知论到他们时说的；因为他在\BibleBook{}{哀歌}中为他们流泪呼唤时，曾说：—— \BibleQuote{熙雍的子女，比精金更贵。}\BibleRef{哀 4:2} 他又说：—— \BibleQuote{圣所的石头，怎么倒在每个街头上？}\BibleRef{哀 4:1} 上主又藉先知说：—— \BibleQuote{那些被丢弃在他土地上的石头原是圣的。}\BibleRef{匝 9:16} 至于这些石头，火在其中燃烧，正如\Person{耶肋米亚}说的：—— \BibleQuote{上主的话在我心中像燃烧的火，在我的骨髓里燃烧。}\BibleRef{耶 20:9} 上主又对\Person{耶肋米亚}说：—— \BibleQuote{看，我将我的话放在你口中如烈火，使这百姓成为木柴。}\BibleRef{耶 5:14} 祂又说：—— \BibleQuote{我的话要像火一般发出，像击碎磐石的铁鎚。}\BibleRef{耶 23:29} 为此，那些先知，就是\Person{提洛的君主} \Person{希兰}曾与之同行的，被称为\Emph{火石}。\par

9. 天主又对他说：\BibleQuote{你曾与那受傅且遮蔽的革鲁宾同在。}\BibleRef{则 28:14}因为那受圣油傅抹的君王，就被称为一个革鲁宾\OriginalTerm{革鲁宾}{Cherub}。他遮蔽了他所有的子民，正如\Person{耶肋米亚}所说的：\BibleQuote{主的受傅者是我们鼻孔的气息，我们曾论及他说：我们在他的荫影下，得以生活于外邦人中间。}\BibleRef{哀 4:20}因为他们原坐在君王的荫影下，当他站在他们为首时。当他们的头顶上的冠冕倒下时，他们便失去了荫庇。\BibleQuote{我们有祸了，因为我们的头顶的冠冕已经倒下了！}\BibleRef{哀 5:16}但基督并未倒下，因为祂第三日复活了。因为君王从犹大家堕落了，他们的国度从此再未建立。祂又说：\BibleQuote{我要毁灭那遮蔽的革鲁宾。}\BibleRef{则 28:16}因为祂所要毁灭的革鲁宾就是\Person{拿步高}，正如经上记载：\BibleQuote{他在\Place{提洛}作工，\Place{提洛}却没有给他军队的酬劳；为补偿\Place{提洛}的工价，将\Place{埃及}地赐给了他。}\BibleRef{则 29:18}\Place{提洛}为何不给\Person{拿步高}酬劳呢？显然因为它的财富已被投入海中，以致\Person{拿步高}没有得到它。在那个时候，祂毁灭了\BibleQuote{那遮蔽的革鲁宾}，就是\Person{拿步高}。因为有两个革鲁宾：一个是受傅且遮蔽的，另一个是遮蔽却不受傅的。因为祂在上文说：\BibleQuote{你曾与那受傅且遮蔽的革鲁宾同在。}\BibleRef{则 28:14}而在下面祂说：\BibleQuote{我要毁灭你，那遮蔽的革鲁宾}；\BibleRef{则 28:16}却没有说\Quote{受傅者}。因为\Person{拿步高}未曾受傅；但\Person{达味}和\Person{撒罗满}，以及他们以后起来的那其他君王，则是受傅的。\Person{拿步高}如何被称为\Emph{遮蔽}的呢？显然是由于那树的异象，当时他看见一棵树立于大地中央，其下住着田野的一切走兽，枝上栖着天空的一切飞鸟，一切血肉都由它得以饱饫。当\Person{达尼尔}为他解释他的梦时，\Person{达尼尔}对他说：\BibleQuote{你就是那棵树，就是你在大地中央所见的树，万民都住在你之下。}因此，他就是那\BibleQuote{遮蔽的革鲁宾}；他毁灭了\Place{提洛}的王子，因为他幸灾乐祸于\Person{以色列}民被掳离开本土，又因他的心高气傲。这个\Place{提洛}也曾荒凉七十年，正如\Place{耶路撒冷}荒凉七十年一样。因为\Person{依撒意亚}论到\Place{提洛}说：\BibleQuote{\Place{提洛}将漂泊七十年，如同一王的岁月；她将与地上万国行淫。}\par

10. 你们这些被高举、被提升的人啊，不要让你们心中的狂傲误导你们，不要说：\Quote{我要上去攻击那富饶之地和那强壮的兽。}因为那兽不会被公绵羊杀死，因为公绵羊的角已被折断。因为公山羊折断了公绵羊的角。\BibleRef{达 7:7}如今，公山羊已成了那强壮的兽。因为当\Person{耶斐特}的子孙执掌王权时，他们便杀了\Place{波斯}王\Person{达理阿}。如今第四兽吞灭了第三兽。这第三兽由\Person{耶斐特}的子孙组成，第四兽由\Person{闪}的子孙组成，因为他们就是\Person{厄撒乌}的子孙。因为当\Person{达尼尔}见到四兽的异象时，他首先见到\Person{含}的子孙，即\Person{尼默洛得}的后裔，也就是\Place{巴比伦}人；其次，\Place{波斯}人和\Place{玛待}人，他们乃是\Person{耶斐特}的子孙；第三，\Place{希腊}人，他们是\Place{玛待}人的弟兄；第四，\Person{闪}的子孙，他们就是\Person{厄撒乌}的子孙。因为\Person{耶斐特}的子孙和\Person{闪}的子孙之间曾结盟。随后，政权从年幼的\Person{耶斐特}子孙手中夺去，交给了年长的\Person{闪}；这事延续至今，并将永远存续。然而，当\Person{闪}的子孙统治终结的时候来到时，那位出自\Person{犹大}子孙的统治者，将在祂第二次降临时，承受王权。\par

11. 因为在\Person{拿步高}的异象中，就是\Person{达尼尔}向他指示和说明的那个，当他看到那立在面前的巨像时，\BibleQuote{这像的头是精金的，胸和臂是银的，腹和股是铜的，腿和脚是铁和陶土混合的。}\BibleRef{达 2:31}\Person{达尼尔}对\Person{拿步高}说：\BibleQuote{你就是那金头。}他为何被称为金头呢？岂不是因为\Person{耶肋米亚}的话应验在他身上吗？因为\Person{耶肋米亚}说：\BibleQuote{巴比伦是主手中的金杯，使全地都喝它的酒。}\BibleRef{耶 51:7}而且\Place{巴比伦}也被称为万国之首，正如经上记载：\BibleQuote{\Place{巴比伦}是\Person{尼默洛得}之首。}\BibleRef{创 10:10}\par

12. 他说\BibleQuote{像的胸和臂是银的。}这指的是一个次于它的王国；即\Place{玛待}人\Person{达理阿}。因为天主把那王国放在天平上。\Person{尼默洛得}家的王国被过秤，被发现欠缺。既然欠缺，\Person{达理阿}就承受了它。因此他说\BibleQuote{他的王国是次等的。}\BibleRef{达 2:39}因为它次等，\Place{玛待}的子孙并未统治全地。然后\BibleQuote{像的腹和股是铜的}，他又说：\BibleQuote{第三王国将统治全地。}\BibleRef{达 2:39}那就是\Person{雅汪}子孙的王国，他们是\Person{耶斐特}的子孙。因为\Person{雅汪}的子孙来攻击他们弟兄的王国。因为\Person{玛待}和\Person{雅汪}是\Person{耶斐特}的儿子。\BibleRef{创 10:2}但\Person{玛待}愚昧，不能掌管王国，直到他的弟兄\Person{雅汪}来，他聪明狡猾，要来毁灭这王国。因为\Person{斐理伯}的儿子\Person{亚历山大}统治了全地。\par

13. \BibleQuote{\Emph{像的腿和脚是铁的}}这是闪的子孙，即厄撒乌的子孙的王国，它坚硬如铁。他说：——\BibleQuote{\Emph{正如铁击碎并制伏一切，同样这第四个王国也将击碎并压碎一切}}\BibleRef{达 2:40}他又解释脚和脚趾，一部分是铁，一部分是陶工泥土。因为他说：——\BibleQuote{\Emph{这样，他们必将与人种混合，但彼此不能粘合，正如铁不能与泥混合}}\BibleRef{达 2:43}这是指着第四个王国。因为在厄撒乌子孙的王国里，没有一位王，王的儿子，被立来治理王国；但当厄撒乌的子孙聚集在一座强大的城时，他们便设立了一个元老院。然后他们常常从那里立一位智者作那城的首领，以治理王国，免得他们王国的主宰衡量他们时，发现他们有所欠缺，而王国就从他们手中被夺去，如同傲慢的\Person{尼默洛得}子孙的王国被夺去，给了愚蠢的\Person{玛代}子孙一样。那被立起的王，总是被前代王的后裔毁掉；他们彼此不能粘合。至于那与人种相比的泥土，意思是：当王被选立治理王国时，他便与铁国的根源相混合。\par

14. 他又指明\BibleQuote{\Emph{在那些将兴起于王国的君王时代，天上的\OriginalTerm{天主}{God}将建立一个永不毁灭、永远长存的王国}}\BibleRef{达 2:44}这就是默西亚君王的王国，它将使第四个王国消逝。他先前又说：——\BibleQuote{\Emph{你看见一块未经人手开凿的石头，打在像的铁与泥土的脚上，将它们击碎}}\BibleRef{达 2:34}他并没有说那石头打像的头，没有打它的胸和手臂，也没有打它的腹和腿，而是打它的脚；因为在整座像中，当那石头来时，将只会遇见脚。在下一节中他又说：——\BibleQuote{\Emph{铁、铜、银和金都一同被粉碎}}\BibleRef{达 2:35}因为在那之后，当默西亚君王为王时，祂必贬抑第四个王国，并击碎整座像；因为整座像指的是世界。它的头是\Person{拿步高}；它的胸和手臂是\Place{玛待}和\Place{波斯}的王；它的腹和腿是\Place{希腊}的王；它的腿和脚是\Person{厄撒乌}子孙的王国；那击碎像并使整个大地充满的石头，是默西亚君王的王国，祂将摧毁今世的王国，并永远为王。\par

15. 且再听\Person{达尼尔}所见的、从海中上来、彼此各异的四兽的异像：它们的外表是这样的：——\BibleQuote{\Emph{第一个相似狮子，具有鹰的翅膀；我见它的翅膀被拔去，它像人一样双脚站立，人的心给了它}}\BibleRef{达 7:4}第二个兽相似熊，\BibleQuote{\Emph{它从一边挺身，口内牙齿间衔着三根肋骨}}第三个兽\BibleQuote{\Emph{相似豹，有四个翅膀和四个头}}第四个兽\BibleQuote{\Emph{极其可怕、强壮有力，有大牙；吞噬并粉碎一切，剩下的用脚践踏}}达尼尔所见的\BibleRef{达 7:2}那大海，就是世界；这四兽就是上文所指的四个王国。\par

16. 关于第一个兽，他论到它说，它相似狮子，具有鹰的翅膀。因为第一个兽就是\Place{巴比伦}王国，它如同狮子。因为\Person{耶肋米亚}这样写道：——\BibleQuote{\Emph{\Place{以色列}是只迷路的羊，狮子们使它们流离。先是\Place{亚述}王吞噬了他，然后那最后的，\Place{巴比伦}王\Person{拿步高}，比他更强}}\BibleRef{耶 50:17}所以\Person{耶肋米亚}称他为狮子。他又说：——\Quote{\Emph{他有鹰的翅膀}}因为经上这样写：当\Person{拿步高}与兽一同去旷野时，他的头发长得如同鹰的羽毛。他又说：——\BibleQuote{\Emph{我见它的翅膀被拔去，它像人一样双脚直立，人的心给了它}}\BibleRef{达 4:30}因为首先，在像的异像中，他被比作金子，比世上一切使用的都宝贵。所以在兽的异像中，他被比作狮子，力胜一切兽。又再被比作鹰，超越一切飞鸟。凡经上关于他所写的，都应验在他身上。因为主论到他说：——\BibleQuote{\Emph{我把铁轭放在万民的颈上，他们要服事\Place{巴比伦}王七十年；连旷野的走兽和天空的飞鸟，我也给了他，好服事他}}既然王如同金头，人便服事他如同王。当他去到旷野，兽便服事他如同狮子。当他的头发如同鹰的羽毛，天空的飞鸟便服事他如同鹰。但当他的心骄傲起来，不知道那能力是从天上赐给他的，铁轭便从人的颈上折断，他就与兽一同出走；王的心被换走，给了他狮子的心。当他凌驾兽类之上，狮子的心也被取去，给了他飞鸟的心。当他的翅膀长出如同鹰的翅膀，他便高居飞鸟之上。然后他的翅膀也被拔去，给了他一颗谦卑的心。当他明了至高者在人的国中掌权，要将国赐与谁就赐与谁时，他便如同人一样赞颂了祂。\par

17. 至于第二个兽，先知论到它说：\BibleQuote{它有如熊，挺立一边}。因为\Place{玛待}和\Place{波斯}王国兴起时，是在东方兴起。\BibleQuote{它口中衔着三根肋骨}。因为那只公绵羊朝西、朝北、朝南抵触，向着天上的三方风。它掌握了这三方风，并加以抵触，就如熊口中的三根肋骨；直到那只公山羊从西方出来，击打公绵羊，并取出它口中的肋骨。\par

18. 关于第三个兽，先知论到它说：\BibleQuote{有如豹子，背上生有四只鸟翼，这兽又有四个头}。这第三个兽就是\Person{马其顿人亚历山大}。因为他勇猛如豹。至于那兽有四翼四头，是因为他在前来杀死\Person{达理阿}并取代他作王之后，将王国交给他四个朋友治理。\par

19. 至于第四个兽，先知论到它说：极其可怖，强壮有力，吞吃、压碎，并用脚践踏所剩下的。这就是\Person{厄撒乌}子孙的王国。因为\Person{马其顿人亚历山大}成为国王之后，希腊人的王国就建立了，因为亚历山大也是希腊人之一。但第三个兽的异象应验在他身上，因为第三个和第四个原本是一个。亚历山大在位十二年。继亚历山大之后兴起的希腊诸王，共有十七王，从\Person{色娄苛尼加诺尔}到\Person{托勒密}，共二百六十九年。而从\Person{奥古斯都}到\Person{斐理伯凯撒}的凯撒们，共有十七王，他们的年数是二百九十三年；再加上\Person{塞维鲁}的十八年。\par

20. 因为\Person{达尼尔}说：\BibleQuote{我注意观看兽头上的那些角，因为那些角就是十个国王}\BibleRef{达 7:8}，他们相继而起，直到\Person{安提约古}。他又说：\BibleQuote{从那些角中又生出一个较小的角，先前的角中有三个竟被它拔除。}\BibleRef{达 7:8}因为当\Person{安提约古}兴起为王时，他就击败了三个王，高抬自己，反对\OriginalTerm{至高者}{Most High}的圣民和\Place{耶路撒冷}。他玷污了圣所\BibleRef{加下 6:2}。他使祭祀和供物停止了一周和半周，即十年半。他将淫乱者引入主的殿，并废除了法律的规条\BibleRef{加下 5:26}。他屠杀义人，将他们交给天空的飞鸟和地上的走兽作为食物。因为在他的日子里，应验了\Person{达味}所说的话：\BibleQuote{天主，外邦人侵入了你的产业，玷污了你的圣殿，使耶路撒冷倾覆；将你仆人的尸首，给天空的飞鸟做食物；将你圣徒的肉，给地上的走兽做食粮。在耶路撒冷四周血如水流，但出来埋葬的人一个也没有。}因为这事在那时应验了：当可敬的老年\Person{厄肋阿匝尔}被杀，以及那七位蒙福的\Person{撒慕纳}的儿子们被杀，还有\Person{犹大·玛加伯}和他的兄弟们为他们的百姓奋斗，住在藏身处之时\BibleRef{加下 5:27}。在那时，\BibleQuote{那角与圣民交战}\BibleRef{达 7:21}，并且他们的势力得胜。那邪恶的\Person{安提约古}\BibleQuote{说亵渎至高者的话，并试图改变节期和法律}\BibleRef{达 7:25}。他废除了\Person{亚巴郎}的盟约，废止了安息的安息日\BibleRef{加下 6:10}。因为他命令犹太人不得行割损。所以，（先知）论到他说：\BibleQuote{他必企图改变节期、法律和盟约，他们将交付他手中，直到一段时期、两段时期和半段时期。}\BibleRef{达 7:25}这时期和半时期，就是一周半，即十年半。他又说：\BibleQuote{审判者已开庭，他们夺去了他的统治权，要将他铲除，彻底消灭，直到结局。}\BibleRef{达 7:26}因为审判临到了\Person{安提约古}，来自天上的审判\BibleRef{加下 9:5}；他患了严重而恶毒的病，由于他腐烂的臭味，没有人能接近他，因为蛆虫从他身上爬出掉落，吃他的肉，因他压迫了\BibleQuote{雅各伯虫豸}\BibleRef{依 41:14}。他活着的时候肉体腐烂，因为他使耶路撒冷子孙的尸体腐烂，不得埋葬。他在自己眼中成了不洁的，因为他玷污了天主的圣所。他祈祷而不得蒙允，因为他不听从他所杀义人的哀声。他曾写信给犹太人，称他们为\Quote{我的朋友，}但天主没有怜悯他，他就在痛苦中死去。\par

21. 先知又说：\BibleQuote{至高者的圣民要承受王国}\BibleRef{达 7:27}。对此我们该说什么呢？以色列子民是否接受了至高者的王国？绝对没有。或者那民是否乘着天上的云彩而来？这已远离了他们。因为\Person{耶肋米亚}论到他们说：\BibleQuote{称他们为被弃的银，因为主已弃绝了他们}\BibleRef{耶 6:30}。他又说：\BibleQuote{他不再看顾他们}\BibleRef{哀 4:16}。\Person{依撒意亚}论到他们说：\BibleQuote{走过去！走过去！不要接近不洁的人}\BibleRef{依 52:11}。而关于至高者的圣民，（\Person{达尼尔}）这样说：\BibleQuote{他们将继承王国直到永远}\BibleRef{达 7:27}。因为这些人\BibleQuote{从君王和首领的重压下稍得苏息}\BibleRef{欧 8:10}，就是从\Person{安提约古}死后直到六十二个七满。\OriginalTerm{人子}{Son of Man}来释放并聚集他们，但他们并没有接受他。因为他来向他们收取果实，他们却未给他。因为他们的葡萄树是\BibleQuote{属索多玛的葡萄树，来自哈摩辣的枝条，一个葡萄园}\BibleRef{申 32:32}，其中长出荆棘，并\BibleQuote{结了野葡萄}\BibleRef{依 5:2}。他们的葡萄树是苦的，他们的果子是酸的。荆棘无法变软，苦味无法变成酒的本性，酸果也无法变成甘甜的本性。\par

22. 因为依撒意亚首先立犹大人作他们的审判者，并在他们中间栽种了一棵新的、可爱的树。但这些人正是那些审判者，\Emph{将坐在十二宝座上，审判以色列十二支派}。\BibleRef{玛 19:28} 祂又这样对审判者说：\BibleRef{依 5:1}——\BibleQuote{\Emph{你们应审断我和我的葡萄园：审判者们，我为我的葡萄园所做的，还有什么没有做到呢？看，我栽种了精选的葡萄枝}，它们却变成了异样的葡萄树。\Emph{我四周围上了篱笆}，由天上的\OriginalTerm{守卫者}{Watchers}守护；\Emph{我建造了其中的塔楼}，即圣殿。\Emph{我掘了榨酒池}，即司铎的洗礼。我降下雨水浇灌它，即我众先知的话语。我修剪它，整理它，除掉阿摩黎人的作为。\Emph{我原指望它结出好葡萄}，即正义，\Emph{它反倒结了野葡萄}，即邪恶与罪恶。\Emph{我指望公平，却见压迫；指望公义，却闻冤声。审判者们，请听，我要怎样对待我的葡萄园：我要拆毁它的篱笆，任它被践踏。我要推倒它的塔楼，任它被掠夺。我要使它变成荒野，因为它结了野葡萄。它将不再被修剪，不再被整理；荆棘和杂草将在其中蔓生。我还要命令云彩不再向它降下雨水。}} 因为天上的守卫者离开了葡萄园的篱笆；他们所倚靠的坚固塔楼被拆毁了。榨酒池，即清除他们罪过之处，也被倾覆了。\BibleQuote{\Emph{当蔓枝完好时，尚且不中用；如今被火烧毁，被焚烧殆尽，还有什么用呢？看，火已吞灭它的两根枝子，内部也已烧毁。}} \BibleRef{则 15:4} 因为它的两根枝子就是两个王国，被烧毁的内部就是\Place{耶路撒冷}。葡萄园的主人派了许多仆人给他们，他们却杀了他们，不把果实交给葡萄园的主人。众仆人之后，又派了爱子去，要从他们收取果实，带回给派遣祂的那位。他们却抓住祂，把祂推出葡萄园外；他们从园中的荆棘上砍下尖刺，钉在祂手上。祂饿了，向他们求食；他们拿葡萄园的果实调成苦胆给祂。祂渴了，向他们求饮；他们给祂醋，祂不肯喝。他们用园中长出的荆棘编了茨冠，戴在葡萄园主人的儿子的头上。因为自从葡萄园建成以来，它所结的就是这样的果实。因此，它的主人把它连根拔起，丢进火里；又在园中栽种了结好果子的葡萄树，使园丁喜悦的葡萄树。因为基督就是葡萄园，祂的父是园丁；凡喝祂爵杯的，就是葡萄枝。因此，一个葡萄园取代了另一个葡萄园。此外，在祂来临之时，祂把王国交给了罗马人，即所谓\Person{厄撒乌}的子孙。这些\Person{厄撒乌}的子孙将为赐予者保守这王国。\par

23. 而圣民继承了永远的王国；这圣民是代替那选民而被拣选的。因为 \BibleQuote{\Emph{祂以那不是子民的民族激起了他们的妒火，以一个愚昧的民族惹动了他们的怒气}。} \BibleRef{申 32:21} 祂释放了圣民。看，天主的每一个盟约，都 \BibleQuote{\Emph{摆脱了君王和首领的重担}。} \BibleRef{欧 8:10} 因为即使一个人曾服侍外邦人，一旦他接近天主的盟约，他就获得自由。但犹太人却在外邦人中劳苦为奴。关于圣徒，祂曾这样说：\BibleQuote{\Emph{他们将继承普天下的王国。}} \BibleRef{达 7:27} 然而，如果这是指着他们（犹太人）说的，那么他们为什么还在外邦人中服役呢？如果他们声称这事尚未发生，那么（我们要问）将要赐给人子的王国，是天上的还是地上的呢？看，王国的子民已受印记，他们已经从这世界得到了解放。因为既然这世界现今存在，它就不会愿意臣服于那位将要来取回祂王国的君王权下。但它会恭敬地保守祂的抵押，这样，当祂来摧毁这王国时，便不会在怒气中临到他们。因为当那位 \BibleQuote{\Emph{王国是属于祂的}} 的人 \BibleRef{创 49:10} 在第二次来临时，祂将收回祂所赐予的一切。祂自己将永远为王。祂的王国永不消逝，因为这是永远的王国。\par

24. 首先，祂将王国赐予雅各伯的子孙，并使厄撒乌的子孙臣服于他们；如依撒格对厄撒乌所说：—— \BibleQuote{你应服事你的弟弟雅各伯}。\BibleRef{创 27:40} 当他们在王国中不再亨通时，祂又从雅各伯子孙手中夺去王国，赐予厄撒乌的子孙，\BibleQuote{直到那应得权柄者来到}。\BibleRef{创 49:10} 他们必将这寄托物交还给它的赐予者，不会以欺诈对待它。王国的守护者服从于那位万物都服从的主。因此，厄撒乌子孙的王国必不交在那些聚集起来、意图上来攻击它的军队手中；因为这王国正为它的赐予者保全，而祂自己将守护它。至于我所写给你们的，亲爱的，即厄撒乌子孙的王国正为它的赐予者保全一事，你们不要怀疑，这王国必不被征服。因为有一位名叫耶稣的强大战士，将以大能来临，并身披王国的全部能力作为铠甲。你们查看便知，祂也曾藉那\OriginalTerm{登记}{poll-tax}被列在他们当中。\BibleRef{路 2:1} 祂既藉那\OriginalTerm{登记}{poll-tax}被列在他们当中，也必援助他们。祂的旗帜在那地方飘扬，他们穿戴祂的铠甲，在战争中必不缺乏。若你们对我说：—— \Quote{在这些之前的列王时代，他们为何能征服并制伏那兽？}那是因为当时在厄撒乌子孙王国中兴起的长官与君王，不愿带领那位与他们一同在\OriginalTerm{登记}{poll-tax}中被列的人同赴战场。因此兽稍被制伏，但未被杀死。\par

25. 但是关于我写给你们的这些事，我亲爱的，就是关于\BibleBook{达尼尔}{书}所载的事，我并未将它们完结，只是到了末了之前（便停止了）。若有人对此争论，就这样对他说：这些话并没有完结，因为天主的\OriginalTerm{圣言}{word of God}是无穷尽的，也不会被完结。因为愚昧人说：\Quote{话就说到这里为止。}另又记着：\BibleQuote{你们不可加添，也不可删减}。\BibleRef{申 4:2} 因为天主的富藏，无法计算，也无从限定。譬如你们从海中取水，几乎看不出减少。又如从海滨移去沙，它的数量不见减损。又如点数天上的星辰，你们不能得其总数。又如从燃烧的火中引火，火丝毫不因之减弱。又如你们领受基督的圣神，基督丝毫不因之减损。又如基督居于你们内，祂却不会在你们内被完成。又如太阳透过窗牖进入你们的房屋，整个太阳却不会全到你们这里。凡此种种我为你们列举的事物，都是由天主的圣言造成的。因此要知道，就天主的圣言而言，没有人达到过、也没有人将会达到它的尽头。所以，你们不要为此争论，只说：—— \Quote{这些事就是这样。这就够了。}但你们当听我这话，也要向我们信仰内的弟兄、我们信仰的子女，究问这些事。但凡讥诮自己弟兄之言的，即使他说：\Quote{我的（话）是智慧的，}你们也不要听他的话。至于我写给你们、关于这些被鼓动去作战的军队，并不是因为有什么启示赐给了我、使我将这些事告知你们，而是要留意信首的话：—— \Emph{凡高举自己的，必被贬抑。}因为即使那些军队上前得胜，你们仍要知道，这是天主的惩罚；而且他们虽胜，也必在公义的审判中被定罪。但你们要确信这事：那兽必在预定的时刻被杀死。然而，我的弟兄，此时要恳切祈求仁慈，好使天主的子民得享平安。\par

\SourceRecord{Translated by John Gwynn.；Nicene and Post-Nicene Fathers, Second Series；Vol. 13.；Edited by Philip Schaff and Henry Wace.；Buffalo, NY: Christian Literature Publishing Co.,；1890.；<http://www.newadvent.org/fathers/370105.htm>.}

% structure: p0001 source=h1 target=section reason=page-title

\section{第六篇论证（论隐修士）}

1. 我所说的话是合宜的，且值得接受：——\Emph{让我们从睡眠中醒来}，\BibleRef{罗 13:11}并举起我们的心和手，朝向天上的天主；免得家主突然来到，\Emph{祂来的时候，见我们正在醒寤}。\BibleRef{路 12:37}让我们留意那荣耀新郎指定的时刻，\BibleRef{玛 25:4}好使我们能与祂一同进入祂的婚室。让我们为灯预备油，好能欢欢喜喜地出去迎接祂。让我们为永居之所准备好必需品，因为那道路是窄狭的。并且让我们除去并丢掉一切污秽，穿上婚宴的礼服。让我们用所领受的银钱去交易，\BibleRef{玛 25:21}好能被称为忠信的仆人。让我们恒心祈祷，好能通过那恐惧所居之处。让我们洗净心中的邪恶，好能瞻仰至高者的荣耀。让我们慈悲为怀，如经上所载，好使天主怜悯我们。\BibleRef{玛 5:7}让我们彼此和平相处，好能被称为基督的弟兄。让我们渴慕义德，好能饱饫，\BibleRef{玛 5:6}从祂王国的筵席上。让我们成为真理的盐，以免成为蛇的食物。让我们清理种子中的荆棘，好能结出百倍的果实。让我们把房屋建在磐石上，\BibleRef{玛 7:24}好使它不被风吹浪打所动摇。让我们成为贵重的器皿，\BibleRef{弟后 2:21}好让主用以供祂使用。让我们变卖所有的一切，为自己买下那颗珍珠，\BibleRef{玛 13:46}好使我们成为富足。让我们积攒财宝于天上，\BibleRef{玛 6:20}好使我们在到来时能打开它们，享受其中的欢乐。让我们在病人身上探望我们的主，\BibleRef{玛 25:33}好使祂能邀请我们站在祂的右边。让我们憎恶自己而爱基督，如同祂爱了我们，并为我们舍弃了自己。让我们尊敬基督的灵，好使我们从祂领受\OriginalTerm{恩宠}{grace}。让我们对世界成为陌生人，\BibleRef{若 17:14}正如基督不属于世界一样。让我们谦卑温和，好能继承生命之地。让我们在事奉祂上孜孜不倦，好使祂能让我们在圣者的居所事奉。让我们纯洁地祈祷祂的祷文，好使它能达于威严的主前。让我们分担祂的苦难，好使我们也能在祂的复活中起来。\BibleRef{弟后 2:11}让我们在身上带着祂的印记，好使我们能从将来的忿怒中得救。因为祂来临的日子是可畏的，谁能承受得住呢？\BibleRef{岳 2:11}祂的忿怒猛烈如火，必将毁灭一切恶人。让我们戴上救赎的头盔，\BibleRef{弗 6:14}好使我们在战斗中不受伤害而死亡。让我们以真理束腰，好使我们在争战中不被发现无能。让我们起来唤醒基督，好使祂平息我们周围的风浪。让我们拿起我们救赎主的福音的准备，作为抵挡那恶者的盾牌。让我们从主领受能力，践踏蛇蝎。\BibleRef{路 10:19}让我们放下忿怒，以及一切愤恨和恶毒。不要让辱骂的话从我们口中说出，我们就是用这口向天主祈祷的。让我们不要咒骂，好使我们从法律的诅咒中得救。让我们作勤奋的工人，好使我们与古时的人同得赏赐。让我们承担整日的重负，好能寻求更丰厚的赏报。让我们不要作闲散的工人，看哪！我们的主已雇用我们进入祂的葡萄园。\BibleRef{玛 20:1}让我们被栽种为祂葡萄园中的葡萄树，因为那是真实的葡萄园。让我们作多结果子的葡萄树，好使我们不被从祂的葡萄园中拔除。让我们成为馨香之气，使我们的香气散发到周围的人身上。让我们在世上一无所有，却以我们主的道理使许多人富足。让我们不要在地上称任何人为父，\BibleRef{玛 23:9}好使我们成为在天之父的子女。虽然一无所有，却拥有一切。\BibleRef{格后 6:9}虽然无人认识我们，但认识我们的人却有许多。让我们时刻在\OriginalTerm{望德}{hope}中喜乐，\BibleRef{罗 12:12}好使那作我们希望和救赎主的因我们而喜乐。让我们公正地审判自己并自责，免得我们在那些将坐在宝座上审判支派的审判者面前，低下头来。\BibleRef{玛 19:28}让我们拿起福音的准备，作为争战的盔甲，\BibleRef{弗 6:16}让我们叩响天堂的门，\BibleRef{玛 7:7}好使它为我们开启，我们得以由此进入。\par

让我们勤恳地祈求仁慈，好能获得我们所需的一切。让我们寻求祂的国和祂的义德，\BibleRef{玛 6:33} 好使我们在那土地上获得增长。让我们思念天上的事，\BibleRef{哥 3:1} 思念天上的事物，并默想它们，在那里基督被高举并受到举扬。但要让我们舍弃那不属于我们的世界，好能抵达我们受邀前往的地方。让我们举目向高天，好能看见那将要彰显的光辉。让我们如鹰展翅上腾，好能在那里看见那身体。让我们为君王预备悦纳的果实——斋戒和祈祷——作为祭献。让我们纯洁地守护祂的质押，好使祂能把全部宝库信托给我们。因为凡对祂的质押行事虚假的，他们就不许他进入宝库。让我们小心守护\OriginalTerm{基督的身体}{body of Christ}，好使我们的身体在号角声中复活。让我们听从新郎的声音，好能同祂一起进入新房。让我们为祂的婚期预备结婚礼物，并欢欢喜喜地出去迎接祂。让我们穿上圣洁的礼服，好能在被选者的首席中坐席。凡不穿婚宴礼服的，\BibleRef{玛 22:13} 他们便把他扔到外面的黑暗里。凡推辞婚宴的，必不得品尝宴席。\BibleRef{路 14:18} 凡喜爱田产和买卖的，必被关在圣徒的城外。凡在葡萄园里不结果实的，必被连根拔除，扔去受折磨。凡从主人那里领了银钱的，应当连本带利归还给那赐予者。\BibleRef{玛 25:16} 凡想成为商人的，就当为自己买下那块田地和其中的宝藏。\BibleRef{玛 13:44} 凡领受好种子的，就当从自己的土地上清除荆棘。\BibleRef{玛 13:7} 凡想成为渔夫的，就当时常撒网。凡为竞赛而训练的，就当使自己远离世界。凡想获得冠冕的，就当像得胜者那样在赛跑中奔跑。凡想下场角力的，就当学会与对手搏斗。凡想奔赴战场的，就当拿起作战的武器，并时刻洁净自己。凡效法天使模样的，就当与世人为陌路。凡背负圣徒之轭的，就当远离获取和花费。凡想获得自己的，就当摒弃世间的利益。凡喜爱天上居所的，就不要为那必将倒塌的泥土建筑劳碌。凡期待被提到云中的，就不要为自己制造华丽的战车。凡期待新郎婚宴的，就不要贪爱现世的宴乐。凡想在那备妥的宴席中享乐的，就当戒绝醉酒。凡为晚餐预备自己的，就不要推辞，\BibleRef{路 14:18} 也不要做商人。凡有好种子落在他身上的，就不要让\OriginalTerm{恶者}{Evil One}在他里面撒下莠子。凡开始建造塔楼的，就当先计算全部花费。\BibleRef{路 14:29} 凡建造的应当完成，免得成为过路人的笑柄。凡把房屋建在磐石上的，就当深挖根基，免得被波浪冲倒。凡想逃离黑暗的，就当趁着有光的时候行走。\BibleRef{若 12:35} 凡害怕在冬天飞行的，\BibleRef{玛 24:20} 就当在夏日预备自己。凡盼望进入安息的，\BibleRef{希 4:11} 就当为\OriginalTerm{安息日}{Sabbath}预备所需。凡恳求主人宽恕的，也当宽恕自己的债户。\BibleRef{玛 18:24} 凡不追讨一百银币的，他的主人必宽免他一万塔冷通。凡把主人的银钱投在钱庄的桌案上，\BibleRef{玛 25:27} 必不被称为无用的仆人。凡喜爱谦卑的，必在生命之地成为承继者。凡愿意缔造和平的，必称为天主的子女。\BibleRef{玛 5:9} 凡知道主人旨意的，就当按那旨意行事，免得受多重的责打。\BibleRef{路 12:47} 凡洁净自己心灵、除去诡诈的，\BibleQuote{他的眼目必瞻仰君王的荣美}。\BibleRef{依 33:17} 凡领受\OriginalTerm{基督的圣神}{Spirit of Christ}的，就当装饰他的内在之人。凡被称为\OriginalTerm{天主的圣殿}{temple of God}的，\BibleRef{格前 3:16} 就当洁净自己的身体，远离一切不洁。凡使基督的圣神忧伤的，\BibleRef{弗 4:30} 必不能从忧伤中抬起头来。凡领受基督的身体的，就当保守自己的身体远避一切不洁。凡脱去\OriginalTerm{旧人}{old man}的，\BibleRef{弗 4:22} 就不要再回到从前的行为中去。凡穿上\OriginalTerm{新人}{new man}的，就当保守自己远避一切污秽。凡从\OriginalTerm{圣洗之水}{water of baptism}穿上盔甲的，就不要脱下盔甲，免得被定罪。凡拿起盾牌\BibleRef{弗 6:16} 抵御恶者的，就当保守自己免遭他射来的箭矢。凡\Emph{退缩}的，他的主就\Emph{不喜悦他}。\BibleRef{希 10:38} 凡思念\OriginalTerm{他主的法律}{Law of his Lord}的，必不被这世界的思虑所困扰。凡默想他主法律的，就像一棵栽在水边的树。凡再信任他主人的，就像一棵栽在河边的树。凡信赖人的，必承受\Person{耶肋米亚}的诅咒。凡受邀赴新郎宴的，就当预备自己。凡点亮自己灯的，就不要让它熄灭。凡期待婚礼呼声的，就当在器皿里备油。\BibleRef{玛 25:6} 凡守门的，就当警醒等候主人。凡喜爱\OriginalTerm{童贞}{virginity}的，就当像\Person{厄里亚}一样。凡背负圣徒之轭的，就当静坐缄默。凡喜爱和平的，就当仰望他的主人，视之为生命的希望。\par

2. 我亲爱的，我们的\OriginalTerm{对敌}{adversary}是诡诈的。与我们对战的那一位是狡猾的。他预备自己攻击勇敢和有名望的人，好使这些人衰弱。因为软弱者本是他的所有物，他不同已被他掳去的俘虏交战。有翅膀的人从他逃离，他投出的飞镖碰不到这人。属神的人在他进攻时看见他，他的\OriginalTerm{全副武装}{panoply}在他们的身上无能为力。所有\OriginalTerm{光明之子}{children of light}都不畏惧他，因为黑暗在光明前奔逃。\OriginalTerm{善者}{the Good}的子女不惧怕\OriginalTerm{邪恶者}{the Evil}，因为天主已将他交在他们脚下，任他们践踏。当他向他们伪装成黑暗时，他们却成为光明。当他像蛇一样爬近他们时，他们却成为盐，是他无法吞吃的。如果他向他们装成毒蛇，他们却变成婴儿。如果他借着贪食欲来攻击他们，他们就像我们的\OriginalTerm{救赎主}{Redeemer}一样，以禁食战胜他。若他企图用眼目的情欲与他们争战，他们却举目仰望高天。若他想用诱惑胜过他们，他们却不给予他聆听的机会。若他公开与他们对战，看，他们就穿上\OriginalTerm{全副武装}{panoply}，站起来抵抗他。若他想借着睡眠攻击他们，他们却警醒戒备，歌唱圣咏并祈祷。若他用财物引诱他们，他们却将财物分施给穷人。若他以甘甜的形式来到他们面前，他们却不品尝，知道他是苦的。若他以\Person{厄娃}的欲望点燃他们，他们就独居，不与厄娃的女儿们同住。\par

3. 因为正是通过\Person{厄娃}，他进入了\Person{亚当}内，\BibleRef{创 3:6} 而亚当因缺乏经验受了诱惑。他又借着主人之妻攻击\Person{若瑟}，但若瑟洞悉他的狡猾，不给他聆听的机会。他借一个女人与\Person{三松}争战，直到夺去他的\OriginalTerm{纳齐尔愿}{Nazariteship}。\Person{勒乌本}是众兄弟中的长子，却因他父亲的妻子，\BibleRef{创 35:22}（那对敌）给他留下了污点。\Person{亚郎}是以色列家的\OriginalTerm{大司祭}{high-priest}，却因他的姐妹\Person{米黎盎}而嫉妒\Person{梅瑟}。梅瑟被派遣去解救人民离开埃及，却带了对他说出可耻行径的女人同行，于是主在路上遇见梅瑟，想要杀死他，直到他让妻子返回\Place{米德杨}。\Person{达味}在一切战事中都得胜，却因一个\Person{厄娃}的女儿而在自己身上发现了污点。\Person{阿默农}容貌俊美，却（那对敌）因对他姐妹的贪欲将他掳去，\Person{阿贝沙隆}因\Person{塔玛尔}受辱而杀了阿默农。\Person{撒罗满}超过世上一切君王，但在他晚年时，他的妻妾们勾引了他的心。\BibleRef{列上 11:1} 借\Person{厄特巴耳}的女儿\Person{依则贝耳}，\Person{阿哈布}的邪恶加重了，他全然成了一个异教人。此外，对敌通过\Person{约伯}的儿女和财产试探他，当无法胜过他时，他就去搬来自己的铠甲，并带来一个厄娃的女儿——那曾使亚当堕落的女人，并通过她的口对正直的丈夫约伯说：\BibleQuote{\Emph{诅咒天主}}。\BibleRef{约 2:9} 但约伯拒绝了她的建议。\Person{阿撒}王也战胜了那当受诅咒的生命，当他想借着阿撒的母亲来攻击他时。\BibleRef{列上 15:13} 因为阿撒知道他的狡猾，就将自己的母亲从高位上移去，并将她的偶像砍碎扔掉。\Person{若翰}大于所有先知，但\Person{黑落德}却因一个厄娃女儿的舞蹈而杀了他。\Person{哈曼}富足，荣誉仅次于君王第三位，但他的妻子却鼓动他灭尽犹太人。\BibleRef{艾 6:13} \Person{齐默黎}是\Person{西默盎}支派的首领，但\Place{米德杨}酋长的女儿\Person{科兹彼}倾覆了他，并且因一个女人，以色列人中一天之内倒毙了两万四千人。\BibleRef{户 25:6}\par

4. 所以，弟兄们，若有人是隐修士或圣人，喜爱\OriginalTerm{独居生活}{solitary life}，却又渴望一位像他一样许过\OriginalTerm{修道誓愿}{monastic vow}的女子与他同居，那么为他而言，倒不如公开地（为妻）娶一个女子，总好过因情欲而放荡。同样，那女子若没有与独居者分开，她倒不如公开出嫁。因此，女人应与女人同住，男人应与男人同住。而且，无论谁渴望在\OriginalTerm{圣洁}{holiness}中持续下去，便不要让他的配偶与他同住，免得他转回从前的状况，而被视为\OriginalTerm{奸夫}{adulterer}。因此，这个劝告是相称、正直且美善的，我给予我自己，也给予你们——我亲爱的独居者们，即那些不娶妻的，以及那些不出嫁的贞女，以及那些爱慕圣洁的人。这是公正、正直且相称的：即使一个人会感到痛苦，他也应当独自生活。正如先知\Person{耶肋米亚}所记载的，他这样居住是合宜的：\BibleQuote{\Emph{人在少年时负起你的\OriginalTerm{轭}{yoke}，独自静坐，乃是好的，因为他已肩负你的轭}}。\BibleRef{哀 3:27} 因为，我亲爱的，凡负起\Person{基督}之轭的人，理应以纯洁保守他的轭。\par

5. 我亲爱的，因为关于\Person{梅瑟}是这样记载的：自从\OriginalTerm{圣者}{Holy One}向他显现的那一刻起，他也爱上了\OriginalTerm{圣洁}{holiness}。而从他受\OriginalTerm{圣化}{sanctified}的时候起，他的妻子就不再服事他。但经上是这样记载的：——\BibleQuote{\Person{农的儿子若苏厄}从幼年就是\Person{梅瑟}的侍从。} \BibleRef{出33:11} 关于若苏厄，经上又这样记载他说：\BibleQuote{他不离开会幕。} \BibleRef{出33:11} 而\OriginalTerm{暂世的会幕}{temporal tabernacle}不由女人服事，因为\OriginalTerm{法律}{the Law}不准女人进入暂世的会幕，甚至她们来祈祷时，也只是在暂世的会幕门口祈祷，然后回去。此外，他命令\OriginalTerm{司祭}{Priests}，在他们供职的时期应当保持圣洁，不可亲近自己的妻子。关于\Person{厄里亚}，经上也这样记载：他有时住在\Place{加尔默耳山}，有时住在\Place{革黎特溪}旁，并由他的门徒服事；因为他的心在天上，天上的飞鸟就给他带来食物；因为他取了天上天使的样式，当他逃避\Person{依则贝耳}时，那些天使亲自给他带来饼和水。因为他将全部心思放在天上，他就被火马车接升天，\BibleRef{列下2:11} 在那里他的居所永远确立。\Person{厄里叟}也步其师的后尘。他常住在\Person{叔能妇人}的楼上，并由他的门徒服事。因为叔能妇人这样说：——\BibleQuote{他是天主的一位圣先知，常从我们这里经过；按他的圣洁，我们理当为他盖一间楼，并为他预备其中所需用的服事。} \BibleRef{列下4:8} 那么，厄里叟楼上所需用的服事是什么呢？显然，仅需床、桌、凳和灯台。但对于\Person{若望}，我们该说什么呢？他也曾住在人群中，却尊严地保守了\OriginalTerm{童贞}{virginity}，并领受了\OriginalTerm{圣神}{Spirit of God}。此外，\OriginalTerm{真福宗徒}{blessed Apostle}论及自己和\Person{巴尔纳伯}时说：——\BibleQuote{难道我们没有权柄吃喝，并带着妻子同行吗？但这不合宜，也不正当。} \BibleRef{格前9:4}\par

6. 因此，弟兄们，因为我们知道也看见，从起初就是借着女人，仇敌得以接近人，并且直到终末他也要借着她来完成这事——因为她是撒殚的武器，并且借着她，他攻击勇士。她时时刻刻为他奏乐，因为从第一天起她就如同为他所用的竖琴。因为她，法律的诅咒得以确立，并且因为她，死亡的许诺得以制定。因为她，她以痛苦生产儿女，并将他们交给死亡。因为她，大地受了诅咒，生出荆棘和蒺藜。因此，借着真福玛利亚的后裔的来临，荆棘被拔除，汗流被抹去，无花果树受诅咒\BibleRef{玛 21:19}，灰尘成为盐\BibleRef{玛 5:13}，诅咒被钉在十字架上\BibleRef{哥 2:14}，刀剑的锋刃从生命树前移开，并赐给信徒作食物，乐园应许给蒙福者、童贞女和圣人。所以生命树的果实被赐给信徒和童贞女作食物，并且对于那些承行天主旨意的人，门已敞开，道路已铺平。泉源涌流，给口渴的人水喝。桌子摆设，晚餐预备。肥牛宰杀，救赎之杯调和。宴席预备，新郎已近，即将入席。宗徒发出邀请，蒙召的人极多。你们蒙拣选的人啊，预备自己吧。光明已照耀，明亮而美好，非人手所制的衣裳已经备妥。婚筵的呼声已近。坟墓将开启，宝藏将显露。死者将复活，生者将飞去迎接君王。宴席已设，号角激励，号筒催促。天上的守望者急忙前行，宝座为审判者设立。劳苦的人将欢欣，无益的人将惧怕。作恶的人不得接近审判者。右边的人将狂喜，左边的人将哀哭切齿。在光明中的人将受光荣，在黑暗中的人将呻吟，以湿润自己的舌头。恩宠已过，公义为王。在那里没有悔改。冬天已近；夏天已过。安息的安息日已来临；劳苦已止息。黑夜已过；光明为王。至于死亡，它的刺已折断，已被生命吞灭\BibleRef{格前 15:54}。那些归于阴府的人将哀哭切齿，而那些进入天国的人将欢欣狂喜，跳舞歌唱赞美。因为那些不娶妻的人将由天上的守望者服侍。那些保守贞洁的人将在至高者的圣所安息。出自父怀的独生子将使所有独修士欢欣。在那里没有男女之分，也没有为奴或自主之别\BibleRef{迦 3:28}，但他们全都是至高者的子女。所有许配给基督的纯真童贞女要点燃自己的灯\BibleRef{玛 25:10}，并与新郎一同进入洞房。所有许配给基督的人都远离了法律的诅咒，并从厄娃女儿的刑罚中被救赎；因为她们没有嫁给人，以致承受诅咒和痛苦。她们不顾念死亡，因为她们没有将孩子交给死亡。代替会死的丈夫，她们许配给了基督。并且，\Emph{因为她们不生育孩子，就赐给她们比子女更好的名号}\BibleRef{依 56:5}。并且代替厄娃女儿的哀叹，她们唱出新郎的歌曲。厄娃女儿的婚宴只持续七日；但对于这些（童贞女）而言，新郎永不离去。厄娃女儿的妆饰是会朽坏磨灭的羊毛，但这些人的衣裳永不会朽坏。衰老使厄娃女儿的美貌凋残，但这些人的美貌将在复活时更新。\par

7. 你们将自己许配给基督的童贞女啊，若有隐修士对你们中的一位说\Quote{我要与你同住，你来服侍我}，你们要这样对他说——\Quote{我已经许配给一位王室的丈夫，我服侍的是祂；如果我离开对祂的服侍而来服侍你，我的新郎必对我发怒，给我写休书，把我逐出祂的家；你虽想从我受尊敬，我也想从你得尊敬，却要小心，免得伤害临到你我。不可怀中藏火\BibleRef{箴 6:27}，免得烧着你的衣裳；但你要独居在尊贵中，我也要独居在我的尊贵中。至于新郎为祂永世婚宴所预备的这些事物，你要备好婚筵的礼物，预备自己迎接祂。至于我，我要预备油，好与明智的童贞女一同进去，而不致被关在门外，与愚拙的童贞女一起。}\par

8. 所以，我亲爱的，请听我写给你们的，就是论到\OriginalTerm{独修士}{solitary}、\OriginalTerm{隐修士}{monk}、\OriginalTerm{贞女}{virgins}、\OriginalTerm{圣人}{saints}当行的事。首先，负轭的人最应当的是：信德坚定，正如我在前一封书信中写给你们的；热心守斋与祈祷；热切于基督的爱；并应谦逊、温良、明智。他的言语应和平悦人，他的思念应对众人真诚。他说话时要仔细掂量，为口设置屏障，免出有害之言，并要远离轻率的笑。不要喜爱服饰的华丽，也不宜蓄长发并装饰，或用香膏涂抹。不可赴宴席，也不宜穿着华丽的衣服。不可放胆饮酒过量。要远离骄傲的思想。不宜注视华丽的衣服，或穿戴精美的服饰。要除去诡诈的口舌；要驱逐嫉妒和忿怒，抛弃诡诈的嘴唇。当人不在场时谈论他的话，不要听也不要接受，免得犯罪，直到查明真相。嘲笑是可憎的过失，将之放在心上是不对的。不要借贷取利，不要喜爱贪婪。要忍受伤害，而不伤害人。此外，要远离混乱，不出口戏谑之言。不要轻看任何悔改的人，不要嘲笑守斋的弟兄，也不要羞辱不能守斋的人。在接纳他的地方，他应规劝人；在不接纳他的地方，他应懂得自重。在适当的时机说话，否则便缄默。不要为肚腹的缘故，因乞讨而招致轻视。要向敬畏天主的人透露自己的秘密；却要远离恶人。不要与恶人谄媚交谈，也不要与仇敌交谈。当这样竭力，使自己毫无仇敌。当人因他的善行而嫉妒他时，他应更增添善行，不要因嫉妒而受损。若拥有，就施舍穷人，应当欢喜；若没有，也不要忧愁。不与恶人交谈，不与傲慢者说话，免得自己陷入傲慢。不与亵渎者争论，免得因他的缘故使主受到亵渎。要远离诽谤者，不要让任何人用花言巧语取悦他人。这些事适合于负起天上之轭，成为基督门徒的独修士。因为基督的门徒理应如此，相似他们的师傅基督。\par

9. 我亲爱的，让我们效法我们的\OriginalTerm{救主}{Saviour}，祂本是富有的，却成了贫穷的；\BibleRef{格后 8:9} 祂本是崇高的，却谦抑了自己的尊威；祂的居所本在天上，却没有枕首之地；\BibleRef{玛 8:20} 祂虽然将要在云彩上降临，\BibleRef{玛 26:64} 却骑着一匹驴驹进入耶路撒冷；\BibleRef{玛 21:2} 祂虽是\OriginalTerm{天主}{God}及\OriginalTerm{天主子}{Son of God}，却取了仆人的形体；\BibleRef{斐 2:6} 祂本是万民的安息，免去一切劳苦，却自己因行路困乏；祂本是解渴的泉源，却自己渴了，向人求水；\BibleRef{若 4:6} 祂本是丰盛，饱饫我们的饥饿，却自己在旷野受试探时饿了；\BibleRef{玛 4:2} 祂是那不眠的守望者，却在海中船上沉睡；\BibleRef{玛 8:24} 祂虽在圣父的帐幕中受服事，却让人手来服事自己；祂是一切病人的医治者，却有铁钉穿透祂的双手；祂的口曾发出美善的言语，人却拿苦胆给祂吃；祂从未伤害任何一人，没有损害过谁，却受人鞭笞、忍受羞耻；祂虽是所有凡人的救主，却将自己交付于\OriginalTerm{十字架}{cross}的死亡。\par

10. 我们救主亲自向我们显示了这一切谦卑。所以，我亲爱的，让我们也谦卑自己。当我们的主超出祂的性体时，祂行在我们的性体之中。让我们安居于我们的性体，好使在审判之日祂让我们分享祂的性体。我们的主离开时，从我们取了质；祂升天时，给我们留下了祂自己的质。那毫无需要的，因我们的需要而设此方便。我们所有的，从太初就是祂的；但祂所有的，谁会给我们呢？然而，我们的主所应许我们的是真实的：——\BibleQuote{我在哪里，你们也要在那里}。\BibleRef{若 14:3} 因为祂从我们所取的，在祂那里备受尊荣，如同王冠系在祂的头上。同样，我们从祂所领受的，我们也应当加以尊荣。我们所有的，在非我们性体的祂那里备受尊荣：让我们尊荣那在祂本性中属于祂的。如果我们尊荣祂，我们就会去到那取了我们的性体而升天的祂那里。但如果我们轻视祂，祂就要从我们收回祂所赐给我们的。如果我们欺诈祂的质，祂就在那里收回祂的，并剥夺祂所应许我们的一切。让我们光荣地颂扬那与我们同在的王子，因为我们中的一人被取为祂的人质。凡尊荣王子的人，必从君王获得许多恩赐。那属于我们而偕同祂的，已经尊荣地坐下了，王冠系在祂头上，祂与君王同坐。至于我们这些贫乏的人，我们当怎样对待那与我们同在的王子呢？祂不需要我们什么，只求我们为祂装饰我们的圣殿；好使时候一到，祂回到圣父那里时，能因我们而感谢圣父，因为我们尊荣了祂。当祂来我们这里时，祂没有我们任何东西，我们也没有任何祂的东西，尽管这两个性体都是祂和圣父的。因为当加俾额尔向那生育祂的真福玛利亚预报时，从高天发出的\OriginalTerm{圣言}{Word}出发而来，并且\BibleQuote{圣言成了血肉，寓居在我们中间}。\BibleRef{若 1:14} 当祂回到派遣祂的那位那里时，祂离去时带走了祂未曾带来之物，正如\OriginalTerm{宗徒}{Apostle}所说：——\BibleQuote{祂使我们与祂一同复活，并在天上与祂同坐}。\BibleRef{弗 2:6} 当祂往圣父那里去时，祂给我们派遣了祂的\OriginalTerm{圣神}{Holy Spirit}，并对我们说：\BibleQuote{我与你们天天在一起，直到今世的终结}。\BibleRef{玛 28:20} 因为基督\Emph{坐在圣父的右边}，而基督\Emph{寓居在人间}。藉着圣父的智慧，祂在天上地下都丰裕充足。祂虽然是一，却寓居在众人内；祂逐一荫庇所有信友，从自己而出，却毫不减损，正如经上所载：\BibleQuote{我要使祂与许多人分受}。\BibleRef{依 53:12} 虽然祂被分施于众人，祂却仍坐在圣父的右边。祂在我们内，我们在祂内，正如祂所说的：\BibleQuote{你们在我内，我也在你们内}。\BibleRef{若 14:20} 祂又在另一处说道：\BibleQuote{我与圣父原是一体}。\BibleRef{若 10:30}\par

11. 如果有人\OriginalTerm{良心}{conscience}缺乏知识，对这事争辩说：——\Quote{基督既是一，祂的圣父也是一，那么基督如何寓居，祂的圣父又如何寓居在信友内呢？义人如何成为天主的圣殿，使得天主寓居在他们内呢？如果这样，每一位信友都有各别的基督，以及在基督内的天主——若是如此，对他们便有许多神明和无数的基督了。} 但是，我亲爱的，请听关于这论点的辩护。让这样说的人从可见之物获得教导。因为人人都知道，太阳固定在天上，然而它的光线播散在地上，它的光透过许多门户和窗子进入房屋；凡日光照射之处，虽仅如手掌大小，也称之为太阳。虽然它照在许多地方，仍如此称呼，但真实的太阳本体在天上。既然如此，难道\Emph{他们就有}许多太阳吗？再者，海水浩瀚，当你从中取一杯时，那也被称为水。纵然你将之分开倒入一千个器皿，它仍按其名被称为水。又如，当你用火点燃多处的火时，那取火的地方在你点燃时不至于匮乏，而火仍被称为同一个名字。你虽将火分布到许多地方，它并不因此就拥有许多名字。当你从地上取尘土，撒在许多地方，它丝毫不会减少，你也不能用许多名字称呼它。同样，天主和祂的基督，虽然是一，却寓居在众多的人内。他们本体在天上，寓居在众人内时丝毫不减损；正如太阳的光力倾注于地上时，太阳在天上也丝毫不减损。天主的德能何其更大，因为太阳本身也是赖天主的德能而存在。\par

12. 我亲爱的，我还要提醒你们经上所记载的。因为经上这样记载：当\Person{梅瑟}独自带领营队成为难以承受的重担时，主对他说：\BibleQuote{\Emph{看！我要将你身上的神能取去，赐给七十位\Person{以色列}的长老。}}\BibleRef{户 11:17} 但当祂从梅瑟的\OriginalTerm{圣神}{Spirit}中取去一些，七十位长老被充满时，\Person{梅瑟}毫无缺乏，也无人能知晓他的圣神被取去了什么。再者，那蒙福的宗徒也说：\Quote{\Emph{天主将基督的圣神分施，差遣到先知们里面。}} 基督毫无损伤，因为\BibleQuote{\Emph{父将圣神赐给祂，是没有限量的。}}\BibleRef{若 3:34} 借着这个反思，你们能领悟基督居住在有信德的人里面；然而基督虽分散在许多人中间，祂却毫无损失。因为先知们领受了基督的圣神，各人按自己所能承受的。而如今，基督的圣神再次倾注在一切有血肉的人身上，\BibleRef{岳 2:28} 儿女们要说预言，老年人要见异梦，青年人要见异象，仆婢们也要说预言。基督的一些在我们内，而基督却在天上，在祂父的右边。基督领受了圣神不是按限量的，而是父爱了祂，将一切交在祂手中，将祂所有产业的权柄赐给了祂。因为\Person{若望}说：\BibleQuote{\Emph{父赐圣神给子，不是按限量的，而是爱了祂，将一切交在祂手中。}}\BibleRef{若 3:34} 我们的主也说：\BibleQuote{\Emph{一切由我父交给了我的。}}\BibleRef{玛 11:27} 祂又说：\BibleQuote{\Emph{父不审判任何人，而是将一切审判都交给了子。}}\BibleRef{若 5:22} 宗徒也说：\BibleQuote{\Emph{万物都将归属于基督，除了那使万物归属于祂的父。当万物通过父归属于子时，子自己也要归属于那使万物归属于祂的天主父，这样天主就是万有中的万有，并在每个人之内。}}\BibleRef{格前 15:27}\par

13. 我们的主为\Person{若翰}作证，说他是先知中最大的。然而他领受圣神是有限量的，因为\Person{厄里亚}怎样领受圣神，\Person{若翰}也怎样获得了它。如同\Person{厄里亚}惯居旷野，圣神也引领\Person{若翰}进入旷野，他惯居于山岭与洞穴中。飞鸟供食于\Person{厄里亚}，而\Person{若翰}惯食飞蝗。\Person{厄里亚}腰束皮带，\Person{若翰}也腰束皮带。\Person{依则贝耳}迫害\Person{厄里亚}，\Person{黑落狄雅}迫害\Person{若翰}。\Person{厄里亚}谴责\Person{阿哈布}，\Person{若翰}谴责\Person{黑落德}。\Person{厄里亚}分开\Place{约旦河}，\Person{若翰}开启了洗礼。\Person{厄里亚}的\OriginalTerm{精神}{spirit}加于\Person{厄里叟}身上多一倍，因此\Person{若翰}按手于我们的救主，祂便领受了无限量的圣神。\Person{厄里亚}开了天门，升了天；\Person{若翰}见天开了，圣神降下，临于我们的救主。\Person{厄里叟}领受了双倍于\Person{厄里亚}的精神；而我们的救主领受了\Person{若翰}的精神与来自天上的圣神。\Person{厄里叟}得了\Person{厄里亚}的外衣，而我们的救主领受了司铎的覆手。\Person{厄里叟}使水变油，而我们的救主使水变酒。\Person{厄里叟}仅以少许面饼饱饫了一百人；而我们的救主用些许饼饱饫了五千人，还不计算孩子及妇女。\Person{厄里叟}洁净了麻风病人\Person{纳阿曼}，而我们的救主洁净了十个（麻风病人）。\Person{厄里叟}咒骂了孩童，他们便被熊撕裂；而我们的救主却祝福孩童。孩童辱骂\Person{厄里叟}，而孩童却以贺三纳光荣我们的救主。\Person{厄里叟}咒骂了他的仆人\Person{革哈齐}；而我们的救主咒骂了祂的门徒\Person{犹达斯}，并祝福了所有（其余的）门徒。\Person{厄里叟}只复活了一个死人；而我们的救主复活了三个人。在\Person{厄里叟}的骨骸上，一个死人复生；但当我们的救主降入阴府时，祂复苏了多人并使他们复活。并且有许多奇迹，是圣神所行的，就是先知们从祂所领受的。\par

14. 因此，我亲爱的诸位，我们也领受了\OriginalTerm{基督之神}{Spirit of Christ}，并且基督\Emph{居住在我们内}，正如经上所载，圣神藉先知的口说过：\BibleQuote{\Emph{我要在他们内居住，在他们中往来}}。\BibleRef{肋 21:12} 所以，我们当预备我们的殿宇为基督之神，不要使他忧伤，以免他离开我们。记得宗徒给我们的劝告：\Quote{\Emph{你们不要使圣神忧郁，因为你们是在他内受了印证，以待得救的日子。}} 因为我们从圣洗领受了基督之神。因为在司铎呼求圣神的那时刻，天开了，他降下，\BibleQuote{\Emph{运行在水面上}}。\BibleRef{创 1:2} 受洗的人就穿上他；因为圣神远离一切由肉身所生的人，直到他们因水重获新生，那时才领受圣神。因为在第一次诞生时，他们带着一个被造在人内、且此后不死不灭的\OriginalTerm{生魂}{animal soul}出生，正如经上说：\BibleQuote{\Emph{亚当成了一个有灵的生物}}。\BibleRef{创 2:7} 但在第二次诞生——即藉着圣洗——他们从\OriginalTerm{天主性}{Godhead}的一个微粒领受了圣神，这圣神就不再屈服于死亡。因为人死的时候，\OriginalTerm{属生魂的灵}{animal spirit}与身体一同埋葬，其感觉也被夺走；然而他们所领受的\OriginalTerm{属天的灵}{heavenly spirit}，按着他的本性，往基督那里去。宗徒把这两者都说明了，因为他说：\BibleQuote{\Emph{播种的是属生灵的身体，复活的是属神的身体}}。\BibleRef{格前 15:44} 圣神按着他的本性，再回到基督那里，因为宗徒又说：\BibleQuote{\Emph{我们出离肉身，便与主同住}}。\BibleRef{格后 5:8} 因为那属神的人所领受的基督之神，是往主那里去的。而那属生魂的灵按着他的本性被埋葬，感觉也被夺走。凡以纯洁保守基督之神的人，当他回到基督那里时，他就这样对他说：\Quote{我所进入的，那从圣洗的水中使我穿上的身体，以圣洁保守了我。} 圣神将为那以纯洁保守他的身体的复活而向基督恳切代求，且要求再次与他结合，为使那身体能在光荣中复活。任何人若从（圣洗的）水领受了圣神，却使他忧伤，他便离开他，直到他死去，回到基督那里，控告那人使他忧伤。当末后的终结时刻来临，复活的时候临近，那曾被纯洁保守的圣神，从其本性获得大能，来到基督前，站在那些以纯洁保守他的人们所埋葬的坟墓门口，等待（复活的）呼声。当\OriginalTerm{守卫者}{Watchers}在君王面前开了天门，\BibleRef{得前 4:16} 那时号角召唤，号筒齐鸣，那等待（复活的）呼声的圣神一听见，就立刻打开坟墓，使身体和其中所埋葬的一切起来，并穿上随之而来的光荣。并且（圣神）在身体内带来复活，光荣在外面装饰身体。属生魂的灵将被属天的灵所吞灭，整个的人成为属神的，因为他的身体被他（圣神）所占有。死亡被生命吞灭，\BibleRef{格后 5:4} 身体被圣神吞灭。靠着圣神的大能，那人将飞升迎接君王，祂要欢欢喜喜地接待他，基督将为那曾在纯洁中保守祂圣神的身体而感谢。\par

15. 我亲爱的诸位，这就是先知们所领受的圣神，我们也照样领受了。他并非时刻常与领受他的人同在，有时他回到派遣他的那位那里，有时又去到领受他的人那里。你们倾听我主所说的话：\BibleQuote{\Emph{你们小心，不要轻视这些小子中的一个，因为我告诉你们：他们的天使在天上，常见我在天之父的面}}。\BibleRef{玛 18:10} 因此，这圣神时常前去，侍立在天主面前，瞻仰他的圣容；凡伤害他居住其中的殿堂的人，他要在天主面前控告他。\par

16. 我要指教你们所记载的事，就是圣神并非时刻常与领受他的人同在。关于撒乌耳，经上这样记载：当他受傅时领受了圣神，但他使圣神忧伤，圣神就\Emph{离开了他}，天主遂打发一个恶神来扰乱他。每逢他被恶神搅扰时，达味就弹琴奏乐，而达味受傅时所领受的圣神就降临，那搅扰撒乌耳的恶神便从他面前逃跑了。可见达味所领受的圣神也并非时刻常与他同在。只要他还在弹琴，圣神才来。因为如果圣神时常与他同在，就不会任凭他与乌黎雅的妻子犯罪了。因为在他为自己的罪祈祷、向天主认罪时，他这样说：\BibleQuote{\Emph{不要从我身上收回你的圣神}}。至于厄里叟，经上也这样记载：\BibleQuote{\Emph{当他弹琴的时候，主的手降在他身上，他便说预言说：\InnerQuote{\Emph{主这样说：你们不见风，也不见雨，这谷中却要充满水沟}}}}\BibleRef{列下 3:15} 还有，当叔能妇人因她死去的儿子来见他时，他对她说：\BibleQuote{\Emph{主向我隐瞒了，没有告诉我}}。\BibleRef{列上 4:27} 然而，当以色列王打发人来捉拿他、要杀他的时候，圣神在使者到来之前预先通知了他，他便说：\BibleQuote{\Emph{看，这个凶手之子，打发人来斩我的头}}。\BibleRef{列下 6:32} 他又预告了次日撒玛黎雅将出现的丰裕。还有一次，当革哈齐偷取了银钱并藏起来的时候，圣神也告诉了他。\par

所以，我亲爱的，当圣神离开一个曾领受祂的人，直到祂返回并来到他那里时，撒殚就接近那人，使他犯罪，好让圣神完全离开他。因为只要圣神与人同在，撒殚就害怕接近他。你们要留意，我亲爱的，我们的主，也是由圣神所生，直到他在圣洗中领受了自天而来的圣神后，才受撒殚的试探。然后圣神引领他去受撒殚的试探。因此，在人身上也是如此；当他发觉自己内心不再热切于圣神，心灵开始偏向世俗的思念时，他就该知道圣神已不在他内，就该起来祈祷、守夜不懈，好使天主的圣神临于他，免得被那仇敌所胜。盗贼不会挖墙进入房屋，直到他看见屋主离开。同样，撒殚也不能接近那房屋，即我们的身体，除非基督的圣神从其中离开。你们要确信，我亲爱的，盗贼并不确实知道屋主是否在家，但他先\Emph{侧耳细听}，并观察。如果他听见屋主在屋内说：\Quote{我要出门远行，}当他探知并看见屋主已出发去办理事务后，盗贼便来，挖墙进入房屋，进行偷窃。但如果他听见屋主的声音在告诫并吩咐家人要警醒看守他的房屋，并说：\Quote{我也在屋内，}那么盗贼就会害怕逃跑，免得被抓住。同样，撒殚也没有预先的知识，不知道或看见圣神何时离开，好前来掠夺人；但他也侧耳细听、观察，然后发动攻击。如果他听见有基督内住的人说可耻的话，或发怒，或争吵，或争斗，那么撒殚就知道基督不与他同在，便前来在他身上成就自己的意愿。因为基督居住在爱好和平和温良的人内，寄居在敬畏他言语的人内，正如他藉先知所说：\BibleQuote{我要垂顾谁呢？我将在谁那里居住呢？岂不是在爱好和平、温良且敬畏我言语的人那里吗？}\BibleRef{依 66:2} 我们的主也说：\BibleQuote{凡遵行我的诫命，并持守我爱的，我们要到他那里去，并要在他那里作我们的住所。}\BibleRef{若 14:23} 但是，如果他听见一个人警醒祈祷，昼夜默想他主的律法，他便转身离开他，因为他知道基督与他同在。如果你们说：\Quote{撒殚是多么众多啊！看，他与许多人作战，}那么，请听我前面关于基督所证明的，并从中学习：无论他怎样被分配在许多人中，他也毫未减少。因为，如同房屋，只要有一线阳光从窗户射入，整个房屋就被照亮；这样，只要有一点撒殚进入一个人，他就完全变黑。请听宗徒所说的：\BibleQuote{撒殚既能使自己化装成光明的天使，那么他的仆役化装成正义的仆役，也不足为奇了。}\BibleRef{格后 11:14} 我们的主又对门徒们说：\BibleQuote{看，我已赐予你们权柄，可以践踏仇敌的势力。}\BibleRef{路 10:19} 圣经也表明，他拥有权柄和仆役。约伯论到他时说：\BibleQuote{天主造了他，叫他去作战。}\BibleRef{约 40:14} 所以，这些他所拥有的仆役，他使他们奔走在世界上，发动战争。但你们要确信，他不会公开作战；因为自从我们的救主来临，（天主）已赐予权柄胜过他。但他肯定会掠夺和偷窃。\par

18. 但是，我所爱的，我要向你们解释宗徒所说的那句话，藉此可以衡量那些作为恶者工具的道理和欺骗的道理。因为宗徒说过：\BibleQuote{有属生灵的身体，也有属神的身体，正如经上记载的：第一个人亚当成了生灵，第二个人亚当成了使人生活的神。}\BibleRef{格前 15:44}所以他们就说，将有两个亚当。但他却说：\BibleQuote{我们怎样带了那由土而出的亚当的肖像，也要怎样带那由天而出的亚当的肖像。}\BibleRef{格前 15:49}因为那由土而出的亚当是犯了罪的那一位，而那由天而出的亚当是我们的救主，我们的主耶稣基督。所以，凡领受了基督的圣神的人，就变得与那属天的亚当相似；\Emph{祂}就是我们的救主，我们的主耶稣基督。因为属生灵的将要被属神的吞灭，正如我先前给你们写的。而那使基督的圣神忧伤的人，在他的复活中将是属生灵的；因为属天的圣神不与他同在，使那属生灵的不被吞灭于其中。但他复活起来时，仍将处于他本性的状态，缺少圣神而赤身露体。因为他从自己身上脱去了基督的圣神，他就被交付于完全的赤身。凡尊敬圣神，并以纯洁守护圣神在自己内的人，在那一天，圣神必要护佑他，他将完全成为属神的，不至于被发现在赤身之中；正如宗徒所说：\BibleQuote{我们若穿上衣服，就不至于被发现赤身了。}\BibleRef{格后 5:3}他又说：\BibleQuote{我们众人都要安眠，但在复活时，我们不都改变。}他又说：\BibleQuote{这可死的将穿上那不死的，这可朽坏的将穿上那不朽坏的；几时这可死的穿上了那不死的，这可朽坏的穿上了那不朽坏的，那时就要应验经上所记载的话：死亡已被胜利吞灭了。}\BibleRef{格前 15:53}他又说：\BibleQuote{在一刹那，眨眼之间，死人必要复活，成为不朽的，我们也要改变。}\BibleRef{格前 15:52}那些将要改变的人，将穿上那属天的亚当的形体，成为属神的。而那些不改变的人，将继续处于亚当受造的本性，即出于尘土的本性，且将继续处于他们在地下的本性中。然后，那些属天的将被提升到天上，他们所穿上的圣神将使他们飞翔，并将承受那从起初就为他们预备了的国度。而那些属生灵的，将因他们身体的重量停留在下地，返回\OriginalTerm{阴府}{Sheol}；在那里要有哀号和切齿。\par

19. 我写下这些，既提醒了自己，也提醒了你们，我所爱的；所以，你们要爱慕童贞，这天上的分，这天上守望者的团体。因为没有甚么能与之相比。在那些如此生活的人内，有基督居住。\BibleQuote{夏季近了，无花果树已发芽，长出叶子了}\BibleRef{玛 24:32}——我们的救主所给的这些征兆已开始应验。因为他说过：\BibleQuote{民族要起来攻击民族，国家攻击国家；并且要有饥荒、瘟疫和从天上来的恐怖。}\BibleRef{路 21:10}看！这一切都在我们的时代应验了。\par

20. 所以，你们，弟兄们，以及爱慕童贞的隐修士们，要研读我写给你们的这一切。并且要提防嘲笑者。因为凡嘲笑和讥讽自己弟兄的，福音上所记载的话正适用于他；即当我们的主愿意和贪婪者及法利塞人算账时。因为经上记载：\BibleQuote{因为他们贪爱钱财，便嗤笑祂。}\BibleRef{路 16:14}同样，如今那些不同意这些事的人，也以同样的方式嘲笑。你们要阅读，要学习。要热心于阅读和实行。要让天主的法律时常成为你们的默想。当你们读了这封信，我恳求你们，我所爱的，起来祈祷吧，并在你们的祈祷中记念我这罪人。\par

\SourceRecord{Translated by John Gwynn.；Nicene and Post-Nicene Fathers, Second Series；Vol. 13.；Edited by Philip Schaff and Henry Wace.；Buffalo, NY: Christian Literature Publishing Co.,；1890.；<http://www.newadvent.org/fathers/370106.htm>.}

% structure: p0001 source=h1 target=section reason=page-title

\section{第八篇论证（论死者的复活）}

1. 对于此事，历来都有争论，\BibleQuote{\Emph{死人将如何复活，他们将带着什么样的身体而来呢}？}\BibleRef{格前 15:35} 看啊！身体逐渐朽坏；骨骸亦然，无疑随着时日渐长，亦将腐朽而不可辨识。当你进入一座埋葬了百具死尸的坟墓，你在那里连一把尘土也找不到。于是那些思虑此事的人这样说：——\Quote{我们当然知道死人将要复活；但他们将穿上一个\OriginalTerm{天上的身体}{heavenly body}和属神的形态。若非如此，这同葬一墓的百具死尸，时日既久，其中已全然无存，当死人得苏生，被穿戴身体而复活时，若非穿上天上的身体，他们的身体将从何而来？你看！坟墓中什么也没有。}\par

2. 如此思虑者，实属愚昧无知。当死者被抬入时，他们曾是某种存在；当他们在那里日久，便化为乌有。待到死者复活之时，那乌有将按其原先本性重成存在，其本性且将添一变化。你这不智之人，如此思虑者，请听那有福的宗徒在教导像你这样的愚人时所说的话；因他说：——\BibleQuote{\Emph{糊涂人哪，你所播的种子若不先死去，就不能复活；并且你所播种的，并不是那将来长出青苗的形体，而是一粒赤裸的麦子，或大麦，或其他某种种子。每类种子，天主都赋予其自有形体。但天主按自己的旨意，给你的种子穿上形体。}}\BibleRef{格前 15:36}\par

3. 因此，愚人哪，你当由此学知：每种种子都穿上了它自己的形体。你从不播下麦种却收获大麦，也从不栽植葡萄却结出无花果；万物都依其本性生长。这样，那安葬在地里的身体，也就是将要复起的身体。至于身体朽坏消散之事，你当由种子的比喻学知：正如种子被撒在地里，腐烂朽坏，又从朽坏中萌发生长，开花结果。那被耕犁过的土地，若不播种，就不能结实，即使那地饮尽了雨水。同样，那未埋葬死者的坟墓，在死人复活时，即或有号筒的巨响在其中吹响，也不会有人从其中出来。而假如，如他们所说，义人的灵魂将上升天堂，并穿上一个\OriginalTerm{天上的身体}{heavenly body}，那么他们已在天堂。那使死者复活者就居住在天堂。那么，当我们的救主来临时，祂将从地中使谁复活呢？祂又为何为我们写下：——\BibleQuote{\Emph{时候要到，且现在就是，死者将要听见人子的声音，凡听见的就要生活，并从坟墓中出来}}？\BibleRef{若 5:28-29} 因为天上的身体不会前来进入坟墓，然后又从坟墓中出去。\par

4. 那些顽固愚昧的人这样说：——宗徒为什么说——\BibleQuote{\Emph{在天上的形体与在地上的形体是不同的}}？\BibleRef{格前 15:40} 但听这话的人，也当听听宗徒所说的另一件事：——\BibleQuote{\Emph{有一个属生灵的身体，也有一个属神的身体}}\BibleRef{格前 15:44}。他又说：——\BibleQuote{\Emph{我们众人不是都要睡眠，我们却都要改变}}\BibleRef{格前 15:51}。他又说：——\BibleQuote{\Emph{这可死的，必须穿上那不可死的；这可朽坏的，必须穿上那不可朽坏的}}\BibleRef{格前 15:53}。他又说：——\BibleQuote{\Emph{我们众人都必须出现在基督的审判台前，为使各人藉他身体所行的，或善或恶，领取相应的报应。}}\BibleRef{格后 5:10} 他又说：——\BibleQuote{\Emph{那些为死者受洗的人，将怎样行呢？如果死者总不复活，为什么还为他们受洗呢？}}\BibleRef{格前 15:29} 他又说：——\BibleQuote{\Emph{如果死人没有复活，基督也就没有复活；如果基督没有复活，你们的信仰便是假的，我们的宣讲也是假的。若这样，我们在为天主作证时，就是假证人，因为我们证实天主复活了基督，而祂并没有复活祂。}}\BibleRef{格前 15:14-15} 所以，如果死人不得复活，也就没有审判。如果没有审判，那就\BibleQuote{\Emph{让我们吃喝吧，因为明天我们就要死了。}}\BibleRef{格前 15:32} \BibleQuote{\Emph{你们不要受迷惑；滥交败坏善行。}}\BibleRef{格前 15:33} 现在，关于宗徒所说的：——\BibleQuote{\Emph{在天上的形体与在地上的形体是不同的}}，你当这样理解这话。当义人的身体复活并改变时，就被称为属天的。而那不改变的，就按它属地的本性，被称为属地的。\par

5. 然而，我亲爱的，请听宗徒所说的另一段类似的话。因他曾说：——\BibleQuote{\Emph{属神的人能审断一切，却不受任何人的审断}}\BibleRef{格前 2:15}。他又说：——\BibleQuote{\Emph{凡随从圣神的人，都思念圣神的事；凡随从肉身的人，都思念肉身的事}}\BibleRef{罗 8:5}。他又说：——\BibleQuote{\Emph{当我们还在肉身时，罪恶的邪欲在我们的肢体内运作，为使我们结出死亡的果实}}\BibleRef{罗 7:5}。他又说：——\BibleQuote{\Emph{如果天主的圣神住在你们内，你们就不属于肉身，而是属于圣神}}\BibleRef{罗 8:9}。宗徒说这一切话时，虽仍穿着肉身，却行圣神的工作。因此，在死者的复活中，义人将被改变，属地的形体将被属天的吞没，并被称为\OriginalTerm{天上的身体}{heavenly body}。而那不改变的，则将被称为属地的。\par

6. 我亲爱的，关于这死者的复活，我要按我的能力指示你们。因为从一开始，天主创造了\Person{亚当}；用泥土形成了他，并使他兴起。因为如果\Person{亚当}不存在时，祂从虚无中造了他，现在祂使他复活岂不更容易？看！他像种子一样被撒在地里。如果天主只做那些对我们来说容易的事，祂的工程在我们看来就不会显为伟大。看！人间有些工匠制造奇妙的东西，而那些不是这些工程的工匠的人，站在那里惊奇它们是如何做成的；在他们看来，同类的工程是困难的。天主的工程岂不更当视为奇迹！但对天主来说，使死人复活并非大事。在种子被撒在地里之前，土地已生出未曾撒入之物。在它未曾怀孕之前，它已在童贞中生产了。那么，土地使曾撒入其中的再生出来，在怀孕后生产，这有何难呢？看！她的产痛临近了；正如\Person{依撒意亚}所说：\BibleQuote{谁曾见过这样的事，谁曾听过这样的事？土地竟在一天之内产痛，一个民族竟在一时之间诞生}？\BibleRef{依 66:8} 因为\Person{亚当}未经撒种就生长；未经怀孕就出生。但看！如今他的后裔被撒种，等候雨水，必将生长。看！土地充满了众多的人，她生产的时候近了。\par

7. 因为我们所有的先祖，都怀着对死者复活和生命的希望，期盼着并急切向往；正如真福宗徒所说：\BibleQuote{如果义人曾经向往亚巴郎所离开的那座城，他们就有机会再转回去；但他们表明，他们向往一个更好的城，就是在天上的那一个}。\BibleRef{希 9:15} 他们期盼着被释放，快快到那里去。从我写给你们的这封信中，你们可以明白并观察到，他们是在期盼复活。因为我们的先祖\Person{雅各伯}临终时，曾用誓言捆绑他的儿子\Person{若瑟}，对他说：\BibleQuote{将我葬在我祖先的坟墓里，与亚巴郎、撒辣、依撒格和黎贝加在一起}。我亲爱的，\Person{雅各伯}为什么不愿意葬在\Place{埃及}，而要葬在他祖先那里呢？他预先表明，他期盼死者的复活；以便当复活的呼喊发出，号角声响起时，他能在祖先近旁复活，而在复活之时不致与那些将回到\OriginalTerm{阴府}{Sheol}并受惩罚的恶人混在一起。\par

8. 同样，\Person{若瑟}也用誓言捆绑了他的兄弟们，\BibleRef{创 50:24} 对他们说：\BibleQuote{当天主眷顾你们时，你们要把我的骨骸从这里带上去}。他的兄弟们照\Person{若瑟}的话做了，并守了这誓言一百二十五年。当上主的军队从\Place{埃及}地出来时，那时\Person{梅瑟}在离开时带上了\Person{若瑟}的骨骸。\BibleRef{出 13:19} 在\Person{梅瑟}看来，这义人的骨骸比以色列子民从\Place{埃及}夺取的金银更宝贵、更好。\Person{若瑟}的骨骸在旷野中四十载；在\Person{梅瑟}安眠之时，他将它们交给\Person{农}的儿子\Person{若苏厄}作为产业。他父亲\Person{若瑟}的骨骸在他眼中比他所征服那地的一切战利品更好。\Person{梅瑟}为什么将\Person{若瑟}的骨骸交给\Person{若苏厄}呢？显然，因为\Person{若苏厄}属于\Person{若瑟}的儿子\Person{厄弗辣因}支派。他把骨骸埋葬在预许之地，为使那地有一珍宝，就是埋葬在那里的\Person{若瑟}的骨骸。还有，在\Person{雅各伯}临终时，他祝福了他的支派，并向他们指明在\Emph{末日}会临到他们的事，他对\Person{勒乌本}说：\BibleQuote{勒乌本，你是我的长子，我的力量和我精力的初果。你过于放荡，如同水，你不能再居首位，因为你上了你父亲的床，你玷污了我的榻！} \BibleRef{创 49:3} 从\Person{雅各伯}安眠到\Person{梅瑟}安眠，历时二百三十三年。于是\Person{梅瑟}愿以他\OriginalTerm{司祭的权柄}{priestly power}，赦免\Person{勒乌本}因与父亲的妾\Person{彼耳哈}同寝所犯的过犯罪恶，好使他在兄弟们复活时，不致从他们的数目中被剪除。因此他在祝福的开端说：\BibleQuote{勒乌本要生存，不至死亡，并且要列在数目中}。\BibleRef{申 33:6}\par

9. 当\Person{梅瑟}应与他的祖先同眠的时刻到来时，他忧伤愁闷，向他的主恳求，希望能进入预许之地。亲爱的，为什么义人梅瑟因未能进入预许之地而忧伤呢？显然，因为他愿意去与他的祖先同葬，而不愿被葬在仇敌之地，即\Place{摩阿布}地。因为\Place{摩阿布}人曾雇佣\Person{贝敖尔}的儿子\Person{巴郎}来诅咒\Place{以色列}。因此梅瑟不愿被葬在那片土地上，免得摩阿布人前来报复他，掘出那义人的骨骸，抛掷于外。主向梅瑟施行了恩宠：祂带他上到\Place{尼波山}，将全地指给他看，使一切地方在眼前经过。当梅瑟观望全地，并看到耶步斯人的山——那里将是圣幕居住之处——时，他看到\Place{赫贝龙}的坟墓，他的祖先\Person{亚巴郎}、\Person{依撒格}和\Person{雅各伯}都葬在那里，他却不能与他们同葬，他的骨骸也不能叠加在他们的骨骸上，好能在复活时与他们一同起来；想到这里，他忧伤哭泣。但是，当他看完全地后，他的主安慰他说：\Quote{我要亲自埋葬你，将你藏起，无人知道你的坟墓。} \BibleQuote{\Emph{于是，梅瑟按主口中所说的话死了；主将他葬在\Place{摩阿布}地，面对\Place{贝特培敖尔}的一个山谷内，就是在\Place{以色列}犯罪的地方；直到今日没有人知道他的坟墓。}}\BibleRef{申 34:5} 他的主不为以色列子民指明他的坟墓，为梅瑟成就了两件美事：一、仇敌不会知道他，从他坟墓中掘出他的骨骸；二、他本族的人不会知道他，将他的坟墓变为朝拜之所，因为他在他本族人的眼中被视如天主。亲爱的，请从此事明白：当他离开他们上山时，他们说——\BibleQuote{\Emph{至于这个从\Place{埃及}地领我们上来的梅瑟，我们不知道他遭遇了什么事。}}\BibleRef{出 32:1} 于是他们为自己造了一只金牛犊并朝拜它，他们记不起天主曾借着\Person{梅瑟}、以\BibleQuote{\Emph{大能的手和伸展的手臂}}\BibleRef{申 5:15} 领他们出\Place{埃及}。为此，天主眷顾了梅瑟，没有指明他的坟墓；免得，倘若指明了他的坟墓，他本族的人会误入歧途，为自己制造偶像，朝拜它，向它献祭，从而因他们的罪惊扰那义人的骨骸。\par

10. \Person{梅瑟}又清楚地宣讲了死者的复活，因为他以他的天主的口吻说：\BibleQuote{\Emph{是我使人死，也是我使人活。}}\BibleRef{申 32:39} 此外，\Person{亚纳}在她的祈祷中这样说：\BibleQuote{\Emph{主使人死，也使人活；祂使人下到\OriginalTerm{阴府}{Sheol}，也使人从那里上来。}}\BibleRef{撒上 2:6} \Person{依撒意亚}先知也这样说：\BibleQuote{\Emph{主，你的亡者将活，他们的尸体将起；睡在尘埃中的人们将醒起歌唱你。}}\BibleRef{依 26:19} \Person{达味}也宣讲说：\BibleQuote{\Emph{看！你为死者行了奇事，勇士们将起来颂扬你，坟墓中的人将述说你的恩宠。}} 而他们如何\Emph{在坟墓中} \Emph{述说天主的恩宠}呢？显然，当他们听到号角召唤的声音，听到从高天发出的角声，以及将要发生的地震和坟墓将裂开时，勇士们将荣耀地起来，在坟墓中彼此述说：\Quote{向我们施行的恩宠何其伟大！我们的希望曾被断绝，但另一希望已为我们升起。我们曾被囚禁在黑暗中，现已出来进入光明。我们曾在可朽坏中被播种，现已复活在荣耀中。我们曾是本性被埋葬，现已属神地复活。我们曾在软弱中被播种，现已在能力中复活。} 这就是他们在坟墓中将要述说的恩宠。\par

11. 亲爱的，天主不仅用言语说：\Quote{我使死者复活，}更借着许多见证以行动显示给我们；使我们对此毫不怀疑。祂事先已显明；因为借着\Person{厄里亚}显示了一个奇迹，为证明死者将活，睡在尘埃中的人将起来。因为当寡妇的儿子死去时，厄里亚使他复活，交还给他的母亲。他的门徒\Person{厄里叟}又使叔能妇人的儿子复活；这样，两个见证就为我们建立并稳固了。又，当以色列子民将一个死人抛在厄里叟的骨骸上时，那死人便复活站了起来。三个见证是确实的。\par

12. 也借着\Person{厄则克耳}先知，死者的复活被清楚地显示出来：当天主领他到平原，向他显示许多骨骸，并使他围绕它们经过，对他说：\BibleQuote{\Emph{人子，这些骨头可以复生吗？}} 厄则克耳回答祂：\BibleQuote{\Emph{你知道，万主之主！}}\BibleRef{则 37:1} 主对他说：\BibleQuote{\Emph{人子，你向这些骨头讲预言，对枯骨说：听万主之主的话！}} 当他使它们听到这些话，就有震动和响声，骨头聚拢起来，甚至那些碎成片断的也都聚拢了。先知见了，大为惊奇，因为骨头从四面而来，每个骨头接受其同伴，每个关节接近其同关节，互相连接排列。它们的干枯变得湿润，关节由韧带连接，血在动脉中变暖，皮肤覆盖在肉上，毛发按其性质生长起来。但它们仍俯卧着，没有气息。于是祂又命令先知，对他说：\BibleQuote{\Emph{你向气息讲预言，说：气息，从四方来，吹在这些被杀者身上，使他们复活。}} 当他使它们听到这第二句话，\BibleQuote{\Emph{气息进入他们内，他们便复活了，站立起来，成了一支极庞大的军队。}}\par

13. 但是，我心爱的，为何那些死人并未藉着厄则克耳所说的那句话而复活？为何他们的复活——无论是骨还是神魂——未藉着那句话完成呢？看哪！藉着一句话，骨便接合起来；再藉另一句话，神魂便到来。这是为使完全的完美留待我们的主耶稣基督，祂在末日要用一句话和一言片语，使每个人的身体复活过来。因为不是那话本身不足，而是传达者稍逊。关于这点，你们要明白并注意：当厄里亚及其门徒厄里叟复活死人时，也不是仅用一句话使他们复活的，而是先祈祷、转求，延误了相当长的时间，然后他们才起来。\par

14. 我们的主自己，在他首次降临时，复活了三个死人，为使三个人的见证得以确定。他每复活一个人都用了两句话。当他复活寡妇的儿子时，他唤了他两次，对他说：\BibleQuote{青年，青年，起来罢！}\BibleRef{路 7:14}他便苏醒过来，就起来了。还有一次，他两次呼唤会堂长的女儿，对她说：\BibleQuote{女孩，女孩，起来罢！}\BibleRef{谷 5:41}她的神魂便返回，她就起来了。拉匝禄死后，当他来到墓地时，他恳切祈祷，大声喊说：\BibleQuote{拉匝禄，出来罢！}\BibleRef{若 11:43}他便苏醒过来，从坟墓中出来了。\par

15. 对于我所解释的这一切——那些死人是各用两句话复活的，这是因为对他们而言，有两个复活发生：先前那个，和将要来临的第二个复活。因为在万民都将起来的那个复活中，没有人会再跌倒；藉着由基督所发出的天主的一句话，所有死人都将在眨眼之间、迅速地复活。因为那使之成就的那一位并不软弱，也不乏能。因为凭一句召唤的话，他将使天涯海角都听见，所有躺在坟墓中的都将跃出起来；而话语不会空空地回到派遣它们的那位那里，而是就如先知依撒意亚上记载的，\BibleRef{依 55:10}他将话语比作雨和雪；因为他这样说：——\BibleQuote{正如雨和雪从天降下，不再返回原处，而是灌溉土地，使之生长，给播种者种子，给食用者食粮，同样，从我口中发出的言语也不会空空地回到我这里，它必完成我所喜悦的，必成就我派遣它所办的事。}因为雨和雪并不回到天上，却在地上完成派遣它们者的旨意。所以，那他将藉着自己的基督所发出的言语，那位基督本身就是\OriginalTerm{圣言}{Word}和讯息，必将带着大能回到他那里。因为当他来临并带来它时，他将如雨和雪一般降下，并藉着他，所有已播种的都将发芽并结出正义的果实，而话语将回到派遣它的那位那里；但他的去并不是徒劳的，他将在派遣他的那位面前这样说：——\BibleQuote{看，我和上主赐给我的孩子们。}\BibleRef{依 8:18}这就是那使死人得生的声音。关于这声音，我们的救赎主作证说：——\BibleQuote{时候要到，那时死人将听见人子的声音，凡听见的就必活下去。}\BibleRef{若 5:25}就如经上所载：\BibleQuote{在起初已有声音，即圣言}\BibleRef{若 1:1}；他又说：\BibleQuote{圣言成了血肉，寄居在我们中间}\BibleRef{若 1:14}。这就是那将从至高之处发出，并使所有死人复活的，天主的声音。\par

16. 又一次，我们的主向撒杜塞人解释关于死人复活的事，那时他们向他提出一个关于那嫁给七个丈夫的妇人的譬喻，并对他说：——\Quote{看！这妇人曾是他们的妻子；在死人复活时，她究竟是哪一位的妻子呢？}\BibleRef{玛 22:28}于是我们的主对他们说：——\Quote{你们大错了，你们既不明了圣经，也不明了天主的能力。因为在那复活后，那堪当来世及从死人中复活的人，男人不娶，女人不嫁，因为他们不能再死，因为他们好像天主的天使，又是复活之子。至于死人复活，你们没有念过圣经上天主对梅瑟所说的话吗？他说：\InnerQuote{我是亚巴郎的天主，依撒格的天主和雅各伯的天主。}看！祂不是死人的天主，而是活人的天主：因为所有人，为祂都是活着的。}\BibleRef{玛 22:29}\par

17. 还有一些人，即使活着，在天主前却是死的。因为天主给亚当颁布了一条诫命，对他说：——\BibleQuote{你吃那树上的果子的那天，必定要死。}\BibleRef{创 2:17}在他违犯了诫命，吃了之后，他活了九百三十岁；但因他的罪，在天主前已算是死的了。但为使你们明白，罪人即使在活着的时候也被称为死人，我要给你们说明。因为在厄则克耳先知书上这样记载说：——\BibleQuote{我凭我的永生起誓——吾主上主的断语——我决不喜欢死去的罪人死亡。}\par

18. 此外，我们的主对那向他说：\BibleQuote{请让我回去埋葬我的父亲，然后我再来跟随你。}\BibleRef{路 9:59} 我们的主又对他说：\BibleQuote{任凭死人去埋葬自己的死人吧！至于你，你要去宣扬天主的国。}然而，我亲爱的，你们如何理解这句话呢？你们曾见过死人埋葬死人吗？一个死人怎能起来埋葬另一个死人呢？但请听我的解释：一个罪人，虽然活着，在天主面前却是死的；而一个义人，虽然死了，在天主面前却是活的。因为这样的死就是睡眠，正如达味说：\BibleQuote{我躺下安睡，我又醒来。}依撒意亚又说：\BibleQuote{睡在尘埃中的，都要醒来。}\BibleRef{依 26:19} 我们的主论及会堂长的女儿说：\BibleQuote{这女孩没有死，只是睡着了。}\BibleRef{玛 9:26} 至于拉匝禄，他对他的门徒说：\BibleQuote{我们的朋友拉匝禄睡着了；我要去叫醒他。}\BibleRef{若 11:11} 宗徒说：\BibleQuote{我们众人都要安眠，但不全要改变。}\BibleRef{格前 15:51} 他又说：\BibleQuote{关于长眠的人，你们不要忧伤。}\BibleRef{得前 4:13}\par

19. 然而，我们理应害怕第二次死亡，那充满哀哭切齿、呻吟悲苦，处于外面黑暗中的死亡。但忠信的人和义人，在那复活中必蒙福，因为他们期望被唤醒，并领受许予他们的美好恩许。至于不信实的恶人，在复活时他们有祸了，因那为他们预备的惩罚！就他们所怀有的信德而言，他们若不复活倒好。因为那主人正为他预备鞭打捆锁的仆人，在睡觉时不愿醒来，因为他知道黎明一到他醒来，主人就要鞭打他、捆绑他。但那善仆，主人已应许给他赏报，就切望黎明之时来到，好从主人那里领受恩赐。即使他正熟睡，在梦中他也见到好似主人即将要给他的，就是主人所应许他的一切，就在梦中欢喜雀跃，满心高兴。至于恶人，他的睡眠并不舒适，因为他想象看哪！黎明已为他来到，就在梦中胆战心惊。但义人安睡，他们的睡眠在白天黑夜都舒适，不理会那漫长的黑夜，在他们眼中如同一小时。然后在黎明更次，他们欢然醒来。至于恶人，他们的睡眠沉重地压着他们，他们如同被大而深的发热所压倒的人，在床上辗转反侧，整夜恐惧，黑夜为他漫长，且害怕那他的主人要惩罚他的黎明。\par

20. 我们的信德这样教导：当人安眠时，他们沉睡不知善恶。义人并不期待他们的恩许，恶人也不期待他们的惩罚判词，直到审判者来临，把那些站在他右边的人与那些站在他左边的人分开。要知道经上所记载的：当审判者就座，案卷在他面前展开，宣读善行恶行时，那行善的，必从善者领受善的赏报；那作恶的，必从公义的审判者遭受恶的惩罚。因为对善人，他并不改变他的本性；他显明自己是公义的，因为他公义地惩罚多人。但对恶人，在那一世界里，恩宠在公义中失落，他就改变他的本性；他显明自己对所有人都是公义的。对于他们，恩宠不会与公义相结合。正如恩宠不能补救损害，公义也不能辅助恩宠。因为恩宠远离审判者，而公义催逼审判者。恩宠若临近某人，他就应当转向它，不要把自己交在公义的手中，免得公义惩罚他，为他的缺失，从他索取惩罚。恩宠若远离某人，公义就会带他到审判台前，他因此被定罪，离去进入刑罚。\par

21. 我亲爱的，你们要听这证据，证明报应要在末期才实现。因为当牧人分开他的羊群，把一些放在他的右边，一些放在他的左边，直到他嘉许善人的服事，那时他才会使他们承受国度；直到他斥责恶人，他们被定罪，那时他才会送他们去受刑。至于那些打发使者跟在君王后边说：\BibleQuote{这个人不该作我们的君王}的人，当他得了国回来时，他的仇敌就要在他面前被杀戮。那些在葡萄园里急忙劳苦的工人，不到工作停止，不会得到工钱。那些领了钱的商人，当钱的主人来到时，他就要索要利润。那些等候新郎的童女，因新郎迟延，都打盹睡着了，当她们听见呼喊时，就要醒来整理灯；明智的进去，糊涂的却被关在门外。那比我们先进入信德的人，若没有我们，\BibleQuote{便不能达到完美。}\BibleRef{希 11:40}\par

22. 我亲爱的，从这一切事上，你们要明白，正如已经向你们确定的，直到现在，还没有人得到他的赏报。因为义人尚未承受国度，恶人也未进入刑罚。牧人尚未分开他的羊群。看！工人进入葡萄园，至今尚未领受工钱。看！商人们正用钱做买卖，他们的主人还没有来算账。君王已经去接受国度，但至今尚未第二次回来。那些等候新郎的童女至今还在睡觉，等候着唤醒他们的呼喊。那些在先的人，在信德上劳苦，要等到最后的人来到，才能达到完美。\par

23. 但是那些在悟性上像婴孩的人说：\Quote{如果没有人得到赏报，宗徒为什么说：\BibleQuote{我们离开身体后，就与主同在}？}\BibleRef{格后 5:8} 但是，亲爱的，请回想我在\WorkTitle{论独修士}中关于这件事教导你们的：义人所领受的\OriginalTerm{神魂}{spirit}，按其属天的本性，要到我们的主那里去，直到复活的时候，那时它要来穿上它曾居住的身体。它时时刻刻在天主台前记得这事，并热切期待那它曾居住的身体的复活，正如\Person{依撒意亚}先知论到外邦人的教会时所说的：\BibleQuote{那记念你的，必忠信侍立在上主面前，你也不会让他们安歇}。但是至于恶人，他们没有人能在上主面前记念他们，因为圣神远离他们，因为他们是\OriginalTerm{属血气的}{animal}，并且像禽兽一样被埋葬。\par

24. 再者，那些随从成为\OriginalTerm{恶者}{Evil One}工具的教义的人，被我们的主所说的话绊倒：\BibleQuote{没有人升过天，除了那自天降下的人子，祂原在天上。}\BibleRef{若 3:13} 他们说：\Quote{看！我们的主证实了，没有任何属地的身体升了天。}他们因无知而不能理解这话的力度。因为当我们的主教导\Person{尼苛德摩}时，他也不理解这话的力度。于是我们的主对他说：— \Emph{没有人升过天，以便下来告诉你们那里的一切。}因为\BibleQuote{如果我对你们说地上的事，你们尚且不信，如果我说天上的事，你们怎么会信呢？}\BibleRef{若 3:12} 因为看！除我以外，没有别的见证者从那上面下来，为证明天上的事，好使你们相信。因为\Person{厄里亚}升到那里去了，但他没有与我一同下来作见证，好使两个见证人的见证可靠。\par

25. 但是，亲爱的，对于死人的复活，你们不要怀疑。因为生活的（天主的）口作证说：\BibleQuote{我使人死，也使人活。}\BibleRef{申 32:39} 这两者都出自同一张口。正如我们确信祂使人死，并且我们看见这事；同样，祂使人活也是确实可信的。从我向你们解释的一切，要接纳并相信：在复活之日，你们的身体将完整地复活，你们将从我们的主那里领受信德的赏报，并且在你们所相信的一切事上，你们必将欢欣鼓舞。\par

\SourceRecord{Translated by John Gwynn.；Nicene and Post-Nicene Fathers, Second Series；Vol. 13.；Edited by Philip Schaff and Henry Wace.；Buffalo, NY: Christian Literature Publishing Co.,；1890.；<http://www.newadvent.org/fathers/370108.htm>.}

% structure: p0001 source=h1 target=section reason=page-title

\section{示范第十篇（论牧者）}

牧者被设立在羊群之上，并将生命的食粮给羊。凡警醒并为自己的羊劳苦的，关心自己的羊群，便是我们的善牧的门徒；善牧为祂的羊舍了自己\BibleRef{若 10:11}。凡不细心领回自己羊群的，便如同那不顾念羊的佣工。牧者们啊，要效法古时那些正义的牧者。\Person{雅各伯}牧放了\Person{拉班}的羊，守护它们，劳苦警醒，因此得了赏报。因为\Person{雅各伯}对\Person{拉班}说：\BibleQuote{看！我这二十年来在你这里：你的母羊和母山羊我从未盗取过，你羊群中的公羊我从未吃过。被野兽撕裂的，没有给你带回，反由我自己赔偿；不论白日黑夜被窃的，你都要我赔偿。白日热风将我吞噬，夜里寒霜将我侵袭。}\BibleRef{创 31:38}\BibleQuote{我的两眼无法入睡。}牧者们，请看那位牧者怎样关心他的羊群。他夜间守更以防卫，时时警醒；白日劳苦以喂养。正如\Person{雅各伯}是牧人，\Person{若瑟}也是牧人，他的兄弟也是牧人。\Person{梅瑟}是牧人，\Person{达味}也是牧人。\Person{亚毛斯}也是牧人。所有这些都是牧人，喂养了羊群并带领得好。\par

现在，我心爱的，这些牧者为何先喂养羊群，而后被选作人群的牧者呢？显然，为要学习一位牧者怎样关心他的羊，怎样警醒并为自己的羊劳苦。当他们学会了牧人的方式，便被选担任牧职。\Person{雅各伯}牧放了\Person{拉班}的羊，劳苦警醒，带领得好；之后，他又牧养教导自己的儿子们，传授他们牧灵工作的样式。\Person{若瑟}曾与他的兄弟们一起牧放羊群；在\Place{埃及}，他成了众多百姓的向导，如同一位好牧人带领他的羊群，将他们领回。\Person{梅瑟}牧放他岳父\Person{耶特洛}的羊，他从牧养羊群中被选去牧养自己的百姓，并作为一位好牧人引导了他们。\Person{梅瑟}将牧杖扛在肩上，走在他所带领的百姓前头，牧养了他们四十年；他为自己的羊群警醒劳苦，是一位勤勉而良善的牧者。当他的主因他们的罪想要毁灭他们——就是他们拜牛犊之时，\Person{梅瑟}祈祷恳求他的主说：\BibleQuote{求你赦免这百姓的罪吧，不然，就从你所写的书上涂抹我吧。}\BibleRef{出 32:31}这是一位极其勤勉的牧者，为自己的羊舍了自己。这是一位卓越的领袖，为他的羊舍了自己。这是一位仁慈的父亲，疼爱并养育了他的子女。伟大而明智的牧人\Person{梅瑟}，懂得如何引领羊群，他教导了\Person{农}的儿子\Person{若苏厄}——一个充满圣神的人，后来他引领羊群，即\Place{以色列}全军。他摧毁诸王，征服土地，将地赐给他们作牧场，为他的羊分赐安息之所和羊栈。此外，\Person{达味}喂养他父亲的羊，从羊群中被提去牧养自己的百姓。于是他按照内心的正直牧养了他们，用他手中的巧妙引导了他们。当\Person{达味}清点羊群数目时，义怒临到他们身上，他们开始被毁灭。于是\Person{达味}为自己的羊舍了自己，祈祷说：\BibleQuote{主啊，天主！我犯了罪，在于我清点了\Place{以色列}。愿你的手攻击我和我父家，这些无辜的羊群，他们犯了什么罪呢？}\BibleRef{撒下 24:17}因此，所有勤勉的牧者都这样为自己的羊舍了自己。\par

然而，那些不顾念羊群的牧者，那些只喂养自己的佣工。为此，先知指责他们说：\BibleQuote{你们这些毁灭和驱散我牧场羊群的牧者，请听主的话。主这样说：看！在旋风之日，我必追讨我的羊，如同牧者追讨他的羊群，我必向你们索要我的羊。无知的牧者啊，你们用羊毛给自己穿衣，吃肥羊的肉，却不牧养羊群。那有病的，你们没有医治；那受伤的，你们没有包扎。瘦弱的，你们没有扶壮；那迷失和被驱散的，你们没有领回。那强壮的肥羊，你们看守了，却以严厉辖制它们。你们自己啃食好牧草，剩下的用脚践踏。你们自己喝清水，剩下的用脚搅浑。我的羊吃你们脚践踏的草，喝你们脚搅浑的水。}这些就是贪婪卑鄙的牧者和佣工，不喂养羊群，也不好好引领它们，也不从豺狼中救出它们。但是当那位伟大的牧者，牧者之首来临时，他必呼唤察视自己的羊，并认识自己的羊群。他必将那些牧者带来，与他们清算，按他们的行为定他们的罪。而那些牧养好羊群的，牧者之首必使他们欢欣，承受生命与安息。\BibleQuote{无明而愚昧的牧人啊！至我怎样将我的羊交付你的右手和右眼，因为你对羊说：该死的就死吧，该灭亡的就灭亡吧，剩下的就吃彼此的肉吧——看！我必使你的右眼失明，使你的右臂枯萎。你那看贿赂的眼必昏盲，你那不按正义行事的手必枯槁。}\BibleRef{匝 11:9}\BibleQuote{而你们，我的羊，我牧场的羊，乃是人；我是主，你们的天主。}\BibleRef{则 34:31}\BibleQuote{看！今后我必在美好丰饶的牧场牧放你们。}\BibleRef{则 34:14}\par

4. \BibleQuote{\Emph{善牧为羊舍掉自己的性命。}}\BibleRef{若 10:11} 祂又说：— \BibleQuote{\Emph{我还有别的羊，不属于这一栈，我也该把他们引来，他们也要听我的声音，这样，将只有一个羊群，一个牧人。父爱我，因为我舍掉我的性命，为再取回它来。}}\BibleRef{若 10:16} 祂又说— \BibleQuote{\Emph{我就是羊的门；凡从我进来的，必得安全；可以进，可以出，可以找着草场。}}\BibleRef{若 10:9} 噢，\OriginalTerm{牧者}{pastors}，你们要效法那位勤奋的\OriginalTerm{牧者}{pastor}，整个羊群的\OriginalTerm{总牧}{chief}，祂如此关心祂的羊群。祂把远处的羊领近，把迷路的羊带回，探望病弱的羊，坚固软弱的羊，包扎受伤的羊，守护肥壮的羊。祂为羊舍命。祂拣选并训导了杰出的领袖，将羊群托付在他们手中，并赐给他们管理整个羊群的权柄。因为祂曾对\Person{西满·刻法}说：— \BibleQuote{\Emph{你喂养我的小羊，你牧放我的绵羊，你喂养我的母羊。}}\BibleRef{若 21:15} 于是西满喂养了祂的羊；他完成了自己的时代，将羊群交付给你们，然后离去了。你们也要好好喂养并引导它们。因为关心羊群的\OriginalTerm{牧者}{pastor}不会同时从事其他事务。他不建葡萄园，不植果园，也不陷入此世的烦恼中。我们从未见过一个牧者将他的羊留在旷野，自己去做商人；也未见一个牧者离开羊群去漫游，而成为农夫。但如果他丢弃羊群而做这些事，就是把羊群交给豺狼了。\par

5. 亲爱的诸位，请记住，我曾写信告诉你们关于我们古代的祖先：他们先学会了牧放羊群的方法，并在其中经历了谨慎的考验，然后才被选为领导者的职务，以便他们能学习并观察牧者何等关心他的羊群，并像他们曾经小心引导羊群那样，也能在这种引导的职务上达于完善。这样，\Person{若瑟}从羊群中被拣选，为在困苦时期引导\Place{埃及}人。\Person{梅瑟}从羊群中被拣选，为引导他的子民并牧养他们。\Person{达味}从追随羊群中被取来，成为治理\Place{以色列}的君王。上主从追随羊群中取来\Person{亚毛斯}，使他成为自己子民的先知。\Person{厄里叟}也是这样，从耕牛后面被取来，成为\Place{以色列}的先知。梅瑟没有回到他的羊那里，也没有离开托付给他的羊群。达味没有回到他父亲的羊那里，而是以正直的心引导自己的子民。亚毛斯没有回头去喂养自己的羊，或去采收树木的果实，而是引导他们并履行先知的职务。厄里叟没有回头去套他的耕牛，而是服事\Person{厄里亚}并接替他的位置。至于那个为他如同牧者的人，因为喜爱田地、贸易、葡萄园、橄榄园和耕种，不愿意成为他的门徒；因此，他没有把羊群托付在他手中。\par

6. 我恳求你们，\OriginalTerm{牧者}{pastors}，不要立那些愚昧蠢笨、贪婪且爱财的人作羊群的领袖。凡喂养羊群的，要吃羊的奶。\BibleRef{格前 9:7} 凡牵引耕牛的，要从自己的劳作中得到供给。司铎有权分享祭台，肋未人应收取十一之物。凡吃奶的，要将心放在羊群上；凡从耕牛的劳作中得供给的，要留心自己的耕种。分享祭台的司铎，要恭敬地服务祭台。至于收取十一之物的肋未人，他们在以色列中没有产业。噢，\OriginalTerm{牧者}{pastors}，我们伟大\OriginalTerm{牧者}{Pastor}的门徒，不要像\OriginalTerm{佣工}{hirelings}；因为佣工不关心羊。要像我们\OriginalTerm{甘饴的牧者}{Sweet Pastor}，祂的生命对祂来说并不比祂的羊更宝贵。养育青年，抚养少女；喜爱羊羔，让它们在你们怀中抚养；这样，当你们来到\OriginalTerm{总牧}{Chief Pastor}面前时，可以把全部羊群完整地献给祂，祂便会赐给你们所应许的：\BibleQuote{\Emph{我在那里，你们也在那里。}}\BibleRef{若 12:26} 这些事虽然简短，对好的牧者和领袖们已足够了。\par

7. 亲爱的诸位，以上我写了，为提醒你们关于全羊群应有的品格。在这篇讲论中，我给你们写了关于牧者，即羊群的引导者的事。亲爱的，这些提醒我是按你们在亲切的来信中向我所请求的，写给了你们。\par

8. 那管家领我进入君王的宝库，向我展示了许多珍宝；我一见，心神就被那巨大的宝库所吸引。我注视它时，它眩惑了我的眼目，掳掠了我的思想，使我的思绪四处游荡。凡从中领受的，自己富足，也能使他人富足。它向一切寻求者敞开，无人看守；虽有许多人从中取用，却不见缺损；当他们将自己所领受的施予出去时，自己的一份反而大大增多。那些白白领受的人，就当白白施予，一如他们所领受的\BibleRef{玛 10:8}。因为这宝库不能以任何代价买卖，因为没有任何事物能与它等价。再者，这宝库永不枯竭；领受它的人不会餍足。他们喝了，仍渴望再喝；吃了，仍觉饥饿。不渴的人，找不到可喝的；不饿的人，找不到可吃的。对它的饥饿能饱足多人，对它的渴望能涌出泉源。因为那走近敬畏天主的人，恰如口渴的人走近泉源、饮水止渴，而那泉源丝毫不会减损。需要饮水的土地，从泉源饮水，泉水却不会枯竭。当土地饮水后，仍需再饮，而泉源并不因涌流而减少。对天主的认识也是如此。即使所有人都从其中领受，其中也不会有任何缺乏，也不能被血肉之子所局限。从它取用的，不能取尽；当他施予时，他也毫无缺乏。你从火焰中用蜡烛取火，即使你点燃许多蜡烛，火焰也不会因你取火而减少，蜡烛也不会因点燃许多而耗损。一个人无法领受君王全部的宝藏；口渴的人从泉源饮水，泉源的水也并不因此枯竭。人站在高山上，他的眼不能同样看清近处和远处；他站着数算天上的星辰时，也无法限定天象的数目。同样，当人走近对天主的敬畏时，他也不能完全达致它；当他领受了许多珍贵事物时，它并不显得减少；当他将他所领受的施予出去时，它也不会耗尽，对他而言仍没有终结。我亲爱的，请记念我在第一篇讲论中关于信德给你们所写的：凡白白领受的，就当白白施予，正如他所领受的，一如我们的主所说：\BibleQuote{你们白白得来的，也要白白分施。}\BibleRef{玛 10:8}。因为凡将他所领受的任何一部分扣留不交的，连他已有的也要被夺去。\BibleRef{玛 25:29}。所以，我亲爱的，正如我现今能从这不枯竭的宝库中得到一些，我就从中寄送给你们。然而，我虽寄送给你们，一切仍与我同在。因为这宝库永不枯竭，因为它是天主的智慧；而管家就是我们的主耶稣基督，正如祂作证说：\BibleQuote{一切都已经由我父交给我了。}\BibleRef{玛 11:27}。再者，祂是那智慧的管家，又如宗徒所说：\BibleQuote{基督是天主的德能和天主的智慧。}\BibleRef{格前 1:24}。这智慧被分赐给许多人，却毫无缺乏，正如我前面给你们解说的；先知们领受了基督的灵，而基督却丝毫没有减损。\par

9. 我亲爱的，我已给你们写了十篇讲论。你们问我的一切，我都给你们解说了，并未从你们收取什么。你们没有问我的，我也给了你们。我问了你们的名字，并写给你们。我自己问了你们的问题，并按我所能回答，为使你们信服。凡我写给你们的，要随时默想这些事；并要努力研读那些在天主的教会中诵读的书籍。我给你们写的这十本小书，它们互相借用，彼此依存。不要把它们彼此分开。我从\Emph{Olaph}到\Emph{Yud}为你们写了，每个字母依次相连。你们和弟兄们、隐修士们，以及那些远离讥诮的信友们，都要阅读、学习；就如我前面写给你们的那样。并要记念我向你们指出的：我没有将这些事讨论到终点，而是离终点尚远。这些事也还不充足；但你们要毫不争辩地听从我这些事，并要与那些善于劝导的弟兄们一同探讨。凡你们听到确实能建树的事，就要领受；凡建立怪异道理的事，就要推倒，彻底拆毁。因为争辩不能建树。而我，我亲爱的，像石匠一样，为建筑运来石头，让聪明的建筑师将它们雕琢并安置在建筑中；所有在建筑中劳苦的工人，都将从家主领受报酬。\par

\SourceRecord{Translated by John Gwynn.；Nicene and Post-Nicene Fathers, Second Series；Vol. 13.；Edited by Philip Schaff and Henry Wace.；Buffalo, NY: Christian Literature Publishing Co.,；1890.；<http://www.newadvent.org/fathers/370110.htm>.}

% structure: p0001 source=h1 target=section reason=page-title

\section{证明十七（论基督，天主之子）}

1.（这是）对犹太人的一个答复，他们亵渎那从外邦人中聚集的子民；因为他们这样说，\Quote{你们敬拜并事奉一个被生的人，一个被钉十字架的人子，并称呼一个人子为天主。而天主没有儿子，你们却论到这位被钉的耶稣，说祂是天主子。}他们提出一个论据，天主曾说：—— \BibleQuote{\Emph{我是天主，在我以外并没有别神。}}\BibleRef{申 32:39} 祂又说：—— \BibleQuote{\Emph{你不可敬拜别的神。}}\BibleRef{出 34:14} 因此，（他们说）你们称一个人为天主，就是反对天主。\par

2. 关于这些事，我亲爱的，以我卑微的能力所能领悟的，我将教导你们，我们虽然向他们承认祂是人，同时我们尊敬祂并称祂为天主和主，但这并非以任何新奇的方式如此称呼祂，也不是我们加给祂一个他们所不用的新奇名号。然而我们确信，我们的主耶稣是天主、天主子、君王、君王之子、出自光明的光明、创造者、谋士、向导、道路、救赎主、牧者、聚集者、门、珍珠、灯；祂被冠以许多（这样的）名号。但我们将放下所有（其余）的名号，并证明论到祂，那从天主而来的，是天主子，且（是）天主。\par

3. 因为那尊崇的\OriginalTerm{天主性}{Godhead}名号也曾加给义人，他们被认为配得如此称呼。那些天主所喜悦的人，祂称他们为我的儿子们，我的朋友们。当祂拣选梅瑟作祂的朋友和所爱的，立他作百姓的首领、师傅和司祭时，就称他为天主。因为祂对他说：—— \BibleQuote{\Emph{我立你作法郎的天主。}} \BibleRef{出 6:1} 并且祂将祂的司祭赐给他作先知，\BibleQuote{\Emph{你的哥哥亚郎要向法郎为你说话；你在法郎面前要如同天主，他就要作你的代言人。}} \BibleRef{出 7:1} 这样，天主不仅对邪恶的法郎立梅瑟为天主，也向圣司祭亚郎立梅瑟为天主。\par

4. 再说，请听我们用以称呼祂的天主子名号。他们说：\Quote{虽然天主没有儿子，你们却使那被钉的耶稣作了天主的长子。}然而，当祂藉梅瑟差遣到法郎那里去，对他说，\BibleQuote{\Emph{以色列是我的长子。我吩咐过你，让我的儿子去事奉我，你若不肯让他去，我要杀你的儿子，你的长子。}}\BibleRef{出 4:22} 并且也藉着先知\BibleRef{欧 11:1} 为这事作证，责备他们说：\BibleQuote{\Emph{我从埃及召了我的儿子来。我怎样呼唤他们，他们竟离开去祭献巴耳，向雕刻的偶像焚香。}} 依撒意亚论到他们说：\BibleQuote{\Emph{我养育儿女，将他们养大，他们竟悖逆我。}}\BibleRef{依 1:2} 经上又记载：\BibleQuote{\Emph{你们是主你们天主的子女。}}\BibleRef{申 14:1} 论到撒罗满，祂说：\BibleQuote{\Emph{他要作我的儿子，我要作他的父亲。}}\BibleRef{撒下 8:14} 同样，我们也称基督为天主子，因为藉着祂我们获得了对天主的认识；正如祂称以色列为\BibleQuote{\Emph{我的长子。}}，又如祂论到撒罗满说：\BibleQuote{\Emph{他要作我的儿子。}}我们也称祂为天主，正如祂以自己的名号给梅瑟命名。达味也论到他们说：—— \BibleQuote{\Emph{你们是神，都是至高者的儿子。}} 当他们不肯改正时，祂论到他们说：—— \BibleQuote{\Emph{你们要像世人一样死去，像任何一位王子一样倒毙。}}\par

5. 因为在世上，\OriginalTerm{天主性}{Divinity}的名号是作为最高的尊荣赐予的，天主喜悦谁，就将这名号加给他。但天主的名字有许多，且是尊崇的，正如祂将祂的名字授予梅瑟，对他说：—— \BibleQuote{\Emph{我是你祖先的天主，亚巴郎的天主，依撒格的天主和雅各伯的天主。这是我的名字，直到永远；这是我的纪念，直到万世。}}\BibleRef{出 3:6} 祂称自己的名为\OriginalTerm{艾希亚·阿舍尔·艾希亚（我是自有者）}{Ahiyah ashar Ahiyah}、\OriginalTerm{全能者}{El Shaddai}和\OriginalTerm{万军之主}{Adonai Sabaoth}。天主就是被这些名号称呼的。那伟大而尊崇的\OriginalTerm{天主性}{Godhead}名号，祂也不吝惜地加给祂的义人；正如，虽然祂是伟大的君王，却毫不吝惜地将伟大而尊崇的\OriginalTerm{王权}{Kingship}名号加给那些身为祂受造物的人。\par

6. 因为天主藉祂先知的口称外邦君王\Person{尼布甲尼撒}为\Quote{万王之王}。因为\Person{耶肋米亚}说：——\Quote{凡不肯将颈项放在我仆人\Person{尼布甲尼撒}——万王之王——之轭下的民族和邦国，我必以饥荒、刀剑和瘟疫惩罚那民族。}虽然祂是大君王，祂并不吝惜将\OriginalTerm{王权}{Kingship}的名号赐予人。同样，虽然祂是伟大的天主，祂却并不吝惜将\OriginalTerm{天主性}{Godhead}的名号赐予血肉之子。虽然万有的父性都属于祂，祂也称人为父。因为祂对会众说：——\Quote{你的子孙要接替你父亲的位置。}虽然\OriginalTerm{权柄}{authority}是属于祂的，祂却将权柄赐给人，使人彼此辖制。\OriginalTerm{钦崇}{worship}本是为尊崇祂，祂却允许世上的人彼此尊崇。因为即使一个人钦崇邪恶者、外邦人和\Quote{拒绝恩宠的人}，天主也不因此指责他。关于钦崇，祂命令祂的子民说：\BibleQuote{你们不可敬拜日月和天上的万象；也不可想要敬拜地上任何受造之物。} \BibleRef{申 4:17} 看，我们良善造物主的恩宠和慈爱：祂不吝惜将天主性、钦崇、王权和权柄的名号赐予人；因为祂是世上万物的父，祂使人受尊荣、被高举、得荣耀，超越一切受造物。因为祂以祂圣洁的双手塑造了人；将祂的圣神吹入人内，并从古时就成了人的\OriginalTerm{居所}{dwelling-place}。祂住在人内，在人中间行走。因为祂藉先知说：\BibleQuote{我要在他们内居住，在他们中间行走。} \BibleRef{肋 26:12} 此外，\Person{耶肋米亚}先知也说：\BibleQuote{你们若改善你们的品行和作为，你们就是主的殿宇。} \BibleRef{耶 7:4} 古时\Person{达味}也说：\Quote{主啊，世世代代祢是我们的居所；高山尚未孕育，大地尚未生产，世界尚未造成之前，从永远到永远，祢就是天主。}\par

7. 你如何理解这事？因为一位先知说：——\Quote{主啊，祢是我们的\OriginalTerm{居所}{dwelling-place}。}另一位说：——\Quote{我要在他们内居住，在他们中间行走。}首先，祂作了我们的居所，然后祂住在我们内，在我们中间行走。因为对智慧人来说，这两件事既真实又简明。\Person{达味}说：——\Quote{主啊，世世代代祢是我们的居所；高山尚未孕育，大地尚未生产，世界尚未造成之前。}我亲爱的，你们知道，天上地下的一切受造物都是先受造的，而人则是最后。当天主决定创造世界及其一切美物时，祂首先在祂的意念中\OriginalTerm{孕育}{conception}并塑造了人；在\Person{亚当}在祂思想中被孕育之后，祂才孕育了受造物；正如祂所说：——\Quote{高山尚未孕育，大地尚未生产之前。}因为人在孕育上比受造物更古老、更在先，但在出生上，受造物却比\Person{亚当}更古老、更在先。\Person{亚当}被孕育并居住在祂的思想中；当\Person{亚当}在孕育中被持守在祂的意念中时，祂藉着口中的言语创造了所有的受造物。当祂完成并装饰了世界，其中一无所缺时，祂才从祂的思想中将\Person{亚当}带出，用祂的手塑造了人；\Person{亚当}看见了已完工的世界。祂将权柄赐给他，管理祂所造的一切，就像一个父亲想要为儿子举办婚宴，为他聘定妻子、建造房屋，预备并装饰儿子所需的一切；然后举办婚宴，将家中的权柄交给儿子。同样，在\Person{亚当}被孕育之后，祂将他带出，并将权柄赐给他，管理祂所造的一切。关于这事，先知说：——\Quote{主啊，世世代代祢是我们的\OriginalTerm{居处}{habitation}；高山尚未孕育，大地尚未生产，世界尚未造成之前。从永远到永远，祢是主。}为了使人不认为有另一位天主，或在先或在后，他说：——\Quote{从永远直到永远，}正如\Person{依撒意亚}说：——\Quote{我是首先的，我是末后的。}当天主从祂思想中将\Person{亚当}带出之后，祂塑造了他，将祂的圣神吹入他内，赐给他分辨的知识，使他能分辨善恶，并能知道是他被天主所造。既然人认识了他的造物主，天主就在他的思想中被形成、被孕育，他就成了他的造物主天主的殿宇，正如经上所载：\Quote{你们是天主的殿宇。}同样，祂亲自说：——\Quote{我要在他们内居住，在他们中间行走。}但那些不认识其造物主的\Person{亚当}的子孙，天主就不在他们内形成，不住在他们内，也不被孕育在他们的思想中；他们在祂面前只算为走兽和其余的受造物。\par

8. 现在，通过这些事，顽固的人将会信服：我们称\Person{基督}为天主之子，这并不是什么奇怪的事。因为，看，他（天主）孕育了众人，并从他的思念中生出他们。他们也必不得不承认，\OriginalTerm{天主性}{Godhead}之名也属於他（\Person{基督}），因为他（天主）也使义人也共享天主的名号。至于我们朝拜\Person{耶稣}——我们藉由他认识了天主——这一点，让他们羞愧吧，因为他们自己也对不洁的外邦人中的异教人俯伏朝拜并尊崇，只要那些人拥有权势；而（对此）却毫无责备。天主已将这种朝拜的尊荣赐给\Person{亚当}的子孙，使他们以此相互尊崇——尤其是他们中间那些卓越而配受尊敬的人。因为，如果他们朝拜并以朝拜之名尊崇异教人——这些人沉溺于异教的邪恶，甚至否认天主之名——却没有将他们当作创造主来朝拜，如同只朝拜他们一样，因而没有犯罪；那么，我们朝拜并尊崇\Person{耶稣}，岂不更加合宜？他扭转了我们顽固的心，使我们脱离一切虚妄错谬的朝拜，并教导我们朝拜、事奉并服役独一的天主，我们的父和我们的创造主。也教导我们明白：世上的君王妄自称神，盗用大天主之名；他们是无信者，且强迫人行无信之事。人们在他们面前俯伏朝拜，事奉并尊崇他们，如同对待雕刻的偶像与神像一样，然而法律从未谴责这些，不视为罪过。正如\Person{达尼尔}也曾朝拜\Place{巴比伦}王\Person{拿步高}——一个无信者且强迫人背叛信仰的人——却未受谴责。\Person{若瑟}也向\Person{法郎}朝拜，圣经并没有记载他因此犯罪。至于我们，我们确实知道\Person{耶稣}是天主，是天主之子；藉着他，我们认识了他的父，并且我们众人都已远离其他一切朝拜。因此，我们无法报答他为我们所承受的一切。但让我们以朝拜向他表达尊崇，以回报他为我们所受的苦难。\par

9. 此外，我们还必须证明，这位\Person{耶稣}早在古时就被先知们预先应许了，且被称为天主之子。\Person{达味}说：\BibleQuote{你是我的儿子，我今日生了你。}他又说：\BibleQuote{在神圣的荣耀中，自母胎、从亘古，我已生了你，如同婴孩。}而\Person{依撒意亚}说：\BibleQuote{因有一婴孩为我们诞生了，有一个儿子赐给了我们；他的王权负载他的肩上；他的名字要称为神奇谋士、永世的大能天主、和平之王。他的王权与和平没有穷尽。}\BibleRef{依 9:6}所以，请告诉我，\Place{以色列}的贤士，这位诞生者是谁？他的名字被称为\Emph{婴孩}和\Emph{儿子}和\Emph{神奇谋士}、\Emph{永世的大能天主}，以及\Emph{和平之王}，他的王权与\Emph{和平}（他说）\Emph{没有穷尽}呢？因为，如果我们称\Person{基督}为天主之子，这是\Person{达味}教导我们的；我们称他为天主，这是我们从\Person{依撒意亚}学来的。\Emph{他的王权负载他的肩上}；因为他背负自己的十字架，走出了\Place{耶路撒冷}。他以婴孩诞生，\Person{依撒意亚}又说道：\BibleQuote{看，童贞女将怀孕生子，他的名字要称为厄玛奴耳，意思是，天主与我们同在。}\par

10. 如果你说基督还没有来临，我也可以容许你这好辩的言论。因为经上记载，当祂来临时，\Emph{外邦人将期待祂}。看！我，一个外邦人，已听说基督必要来临。当祂尚未来临时，我已经预先信了祂；并且借着祂，我敬拜以色列的天主。当祂来临时，祂难道会责备我，因我在祂来临之前就预先信了祂吗？但你，愚昧的人，先知不容许你说基督还没有来临；因为\Person{达尼尔}反驳你，说：\BibleQuote{\Emph{过了六十二个星期，\OriginalTerm{默西亚}{Messiah}必要被处死。在祂来临时，圣城将被毁灭，其结局将如洪水；直到预定之事成就，这城将一直荒凉。}}\BibleRef{达 9:26} 你期待并希望，在基督来临时，以色列将从各地聚集，\Place{耶路撒冷}将被建立并有人居住。但\Person{达尼尔}作证说，当基督来临并被处死后，\Place{耶路撒冷}将被毁灭，且将持续荒凉，直到所定之事永远成就。关于基督的苦难，\Person{达味}说：\BibleQuote{\Emph{他们穿透了我的手脚，我的骨骼都在呼喊。他们注视我、察看我，瓜分我的衣服，为我的长衣拈阄。}} \Person{依撒意亚}说：\BibleQuote{\Emph{请看！我的仆人必要被认识、被启示、被举扬，以致众人对祂惊愕不已。这人的容貌被毁，甚于众人，祂的形状甚于人子。}}\BibleRef{依 52:13} 他又说：\BibleQuote{\Emph{祂要洁净许多民族，君王要在祂面前惊异。}}\BibleRef{依 52:15} 他在那段话中说：\BibleQuote{\Emph{他在祂面前生长如嫩芽，又如出自干地中的根苗。}}\BibleRef{依 53:2} 在那段话的结尾他说：\BibleQuote{\Emph{他为了我们的罪恶被刺死；为了我们的过犯受辱；我们平安的惩罚已落在祂身上，因祂的创伤，我们得痊愈。}}\BibleRef{依 53:5} 人凭借什么创伤得了痊愈？\Person{达味}并没有被处死，因为他安享高寿，死后葬在\Place{白冷}。如果他们说是论\Person{撒乌耳}的，因为\Person{撒乌耳}在\Place{基肋波亚}山上与培肋舍特人交战时阵亡，并且说他们\Emph{穿透了他的手脚}，因为他的尸体被钉在\Place{贝特商}的城墙上；然而这并不适用于\Person{撒乌耳}。当\Person{撒乌耳}的肢体被穿透时，他的骨骼已无知觉，因为他已经死了。\Person{撒乌耳}死后，他们才将他的尸体和他儿子们的尸体悬挂在\Place{贝特商}的城墙上。但\Person{达味}说\BibleQuote{\Emph{他们穿透了我的手脚，我的骨骼都在呼喊。}}时，他在下一节中说：\BibleQuote{\Emph{天主，求你速来援助我；从刀剑下救出我的灵魂。}} 如今，基督从刀剑下被救出，从\OriginalTerm{阴府}{Sheol}中上来，恢复生命，第三日复活，这样，\Emph{天主速来援助了祂}。但\Person{撒乌耳}呼求主，主没有应允他；他透过先知求问，也没有得到答复。他改装自己，求问女巫，从那里得知实情。他在培肋舍特人面前战败，见战事已失利，就用自己的刀剑自刎。再者，在那段话中\Person{达味}说：\BibleQuote{\Emph{我要向我的弟兄宣扬你的名，在会众中赞颂你。}} 这些事怎能应用于\Person{撒乌耳}呢？\Person{达味}又说：\BibleQuote{\Emph{你决不让你的圣者见到腐朽。}} 但这一切都完全适用于基督。当祂来到他们中间时，他们不接纳祂；反而以假见证邪恶地审判了祂。祂被人悬挂在木架上，他们用钉子穿透祂的手脚，将祂钉牢；祂的骨骼都在呼喊。在那一天，发生了一件大奇事，即在正午时分，光明变为黑暗，正如\Person{匝加利亚}预言说：\BibleQuote{\Emph{那日子为主所知。那不是白昼，也不是黑夜；到了晚上，必有光明。}}\BibleRef{匝 14:7} 那么，以奇事为特征的那日子，既非白昼也非黑夜，到了晚上却有光明，是什么日子呢？显然就是他们钉死祂的那一天，因为在那一天的正午，黑暗降临，到了晚上却有光明。他又说：\BibleQuote{\Emph{那日必有寒冷和冰霜。}} ——正如你们所知，在他们钉死祂的那一天，天气寒冷，他们生了一堆火取暖，那时\Person{西满}来同他们站在一起。他又说：\BibleQuote{\Emph{刀剑兴起，攻击牧人，攻击我同伴；}\Emph{击打牧人，羊群四散；我必转回我的手攻击那牧者。}} 此外，\Person{达味}论到祂的受难说：\BibleQuote{\Emph{他们以苦胆给我当食物，我渴了，他们给我醋喝。}} ——他在那段话中又说：\BibleQuote{\Emph{他们迫害了祢所击打的人，加增了被杀者的痛苦。}} 因为他们加给他许多苦难，许多经上没有记载的事，咒骂和侮辱，这些经文无法尽述，因为他们的辱骂是可恨的。然而，\BibleQuote{\Emph{主乐于羞辱祂，使祂受苦。}}\BibleRef{依 53:10}，\BibleQuote{\Emph{他为了我们的过犯被刺死}}，\BibleQuote{\Emph{为了我们的罪恶受辱}}，并且\BibleQuote{\Emph{替我们成了罪。}}\BibleRef{格后 5:21}。\par

11. 我们朝拜那些仁慈，并在祂圣父的尊威前屈膝下拜，因为是祂将我们的敬拜转向了祂自己。我们称祂为天主，正如\Person{梅瑟}被称为天主；又称祂为\OriginalTerm{首生者}{Firstborn}和子，正如\Person{以色列}被称为\OriginalTerm{首生者}{Firstborn}和子；又称祂为耶稣（\Person{若苏厄}），正如\Person{农的儿子若苏厄}曾如此被称一样；又称祂为司祭，相似\Person{亚郎}；为君王，相似\Person{达味}；为伟大先知，相似众先知；为牧者，相似那些牧养引导\Person{以色列}的牧者。祂也如此称呼子女们，如祂所说：\BibleQuote{外邦的子民要听从我。}祂使我们与祂自己成为兄弟，祂曾说：\BibleQuote{我要向我的弟兄宣扬祢的名。}我们也成了祂的朋友，正如祂对门徒们说的：\BibleQuote{我称你们为朋友，}\BibleRef{若 15:15} 正如祂的圣父称\Person{亚巴郎}为\BibleQuote{我的朋友}。\BibleRef{依 41:8}祂对我们说：\BibleQuote{我是善牧、门、道路、葡萄树、撒种者、新郎、珍珠、灯、光、君王、天主、救主、救赎者。}祂被冠以许多名号。\par

12. 我给你们写了这简短的论证，我亲爱的，为叫你们能以此反驳犹太人，论到他们说天主没有儿子，以及论到我们称祂为天主、天主子、君王和\OriginalTerm{一切受造物的首生者}{Firstborn of all creatures}。\BibleRef{哥 1:15}\par

\SourceRecord{Translated by John Gwynn.；Nicene and Post-Nicene Fathers, Second Series；Vol. 13.；Edited by Philip Schaff and Henry Wace.；Buffalo, NY: Christian Literature Publishing Co.,；1890.；<http://www.newadvent.org/fathers/370117.htm>.}

% structure: p0001 source=h1 target=section reason=page-title

\section{第21篇 论迫害}

1. 我曾听到一种辱骂，这使我大为烦恼。那些不洁的人（即异教徒）说，这个由万民中聚集起来的民族，没有天主。那些不虔敬的人也这样说：\Quote{如果他们有天主，为什么祂不替自己的民族报仇呢？} 黑暗更浓重地笼罩了我，因为犹太人也辱骂我们，并在我们民族的子女面前自高自大。有一天，一个在犹太人中被看作智者的人问我，说：—— \Person{耶稣}，那位被称为你们师傅的，为你们写道：\BibleQuote{你们若有信德像一粒芥子，你们就要对这座山说：\InnerQuote{移开！}它就会从你们面前移开；甚至（你们也会说）：\InnerQuote{被举起，投入海中！}它也会服从你们。} 这样看来，在你们全体人民中，竟没有一个智者，他的祈祷蒙垂听，他祈求天主使你们的迫害者停止迫害你们。因为在那段经文里，确是为你们记载着：\BibleQuote{你们没有不能做的事。}\par

2. 当我看到他亵渎并大肆攻击这道（即基督宗教）时，我的心意纷乱，我知道他不会接受我对他引用的那些话的解释。于是我也就法律和先知书中的话问他，对他说：—— 你相信即使在你们分散时，天主仍与你们同在吗？他向我承认说：\Quote{天主与我们同在，因为天主曾对以色列说过：—— \BibleQuote{即使在他们敌人的境内，我仍没有离弃他们，也没有废弃我与他们所立的盟约。}}\BibleRef{肋 26:44} 我回答他说：—— \Quote{我听到你这样说，这很好，即天主与你们同在。我就要用你的话来反驳你。因为我说过，先知曾对以色列说过，这是出自天主的口的话：—— \BibleQuote{如果你经过海洋，我必与你同在，河流不会淹没你；如果你在火中行走，你必不会被烧伤，火焰也不会烧到你；因为上主你的天主与你同在。}}\BibleRef{依 43:2} 因此，在你们全体民众中，没有一个人是义人、善人和智者，能够经过海洋而存活，不被淹没；或经过河流而不被淹没；或能在火中行走，看看自己是否不会被灼伤，火焰是否会烧着他。如果你给我一个解释，我也不会被你说服，正如你也不接受我对你发问的那些话的解释一样。\par

3. 此外，我就记在\WorkTitle{厄则克耳}上的另一句话问他；即他对\Place{耶路撒冷}所说的：—— \BibleQuote{\Place{索多玛}和她的女儿们将被重建，如同古时一样；你和你的女儿们也将如同古时一样。}\BibleRef{则 16:55} 于是他给我解释这句话，开始辩护，对我说：\Quote{关于天主借先知对\Place{耶路撒冷}所说的，\BibleQuote{\Place{索多玛}和她的女儿们将被重建，如同古时一样；你和你的女儿们也将如同古时一样，}这句话的意思是，\Place{索多玛}和她的女儿们将像古时一样处在自己的地方，并受\Place{以色列}管辖；\Place{耶路撒冷}和她的女儿们将像古时一样处于王权的荣耀中。} 当我听到他这种辩护时，在我眼中甚为可鄙，我便对他说：—— \Quote{既然先知的话是在愤怒中说的，那么整段都是愤怒的，还是部分是愤怒的，部分是和善的呢？} 他回答说：—— \Quote{一段愤怒的经文完全是愤怒，其中没有平安。} 我便对他说：—— 既然你教导我说那段愤怒的经文中没有平安，那就不要争辩，也不要亵渎地听，我要就这句话指教你。因为从头到尾，整段话都是在愤怒中说的。因为他曾对\Place{耶路撒冷}说：—— \BibleQuote{我指着我的生命起誓——上主天主的断语——你的妹妹\Place{索多玛}和她的女儿们，并没有行过你和你的女儿们所行过的事。}\BibleRef{则 16:48} 又对她（\Place{耶路撒冷}）说：—— \BibleQuote{你应当羞愧，承担你的羞辱，因为你在你的罪孽上超过了你的姐妹，她们比你还显得正直。}\BibleRef{则 16:52} 既然他说\Place{索多玛}和她的女儿们比\Place{耶路撒冷}和她的女儿们还显得正直，\Place{耶路撒冷}在罪孽上超过了\Place{索多玛}，那么当\Place{以色列}被聚集时，它的所在地理应在\Place{索多玛}和\Place{哈摩辣}。因为 \BibleQuote{他们的葡萄树是来自\Place{索多玛}的葡萄树，是来自\Place{哈摩辣}田园的栽植；他们的葡萄是苦的，他们结的穗子是苦的。}\BibleRef{申 32:32} \Person{依撒意亚}也称他们为 \BibleQuote{\Place{索多玛}的首领，\Place{哈摩辣}的百姓。}\BibleRef{依 1:10} 因为如果\Place{以色列}被聚集，他们理应与\Place{索多玛}的首领和\Place{哈摩辣}的百姓同住在\Place{索多玛}和\Place{哈摩辣}；吃\Place{索多玛}的葡萄树和\Place{哈摩辣}的栽植中摘取的苦葡萄，收集苦穗子；吃毒蛇的蛋，穿蜘蛛网的衣服，\BibleRef{依 59:5} 与葡萄园的野葡萄一同被使用，\BibleRef{依 5:2} 变为被弃的银渣。\BibleRef{耶 6:30} 而比\Place{耶路撒冷}更正直的\Place{索多玛}和她的女儿们，将被重建如同古时一样。至于罪孽超过\Place{索多玛}的\Place{耶路撒冷}，将继续留在她的罪恶中，\BibleQuote{要留在荒凉中，直到命定的终结。}\BibleRef{达 9:27}\par

\Person{厄则克耳}说：——\BibleQuote{这是\Place{索多玛}和她的女儿的罪恶：她们没有扶持穷人和贫乏者的手；当我看见她们这些事，我倾覆了她们。} \BibleRef{则 16:49} 你们要思考并观察：从\Place{索多玛}倾覆直到\Place{耶路撒冷}建成，共有八百九十六年。从\Person{亚巴郎}藉天使蒙\Person{天主}告知\BibleQuote{明年此时我必回到你这里，你的妻子\Person{撒辣}要有一个儿子} \BibleRef{创 18:14}的时候，直至\Person{雅各伯}进入\Place{埃及}，共有一百九十一年；\Person{雅各伯}的子孙在\Place{埃及}住了二百二十五年。这样，从\Person{依撒格}受孕及\Place{索多玛}倾覆时起，共有四百一十六年；从\Place{以色列}出离\Place{埃及}，直到\Person{撒罗满}建成\Place{耶路撒冷}的大厦及圣殿，共有四百八十年。因此，从\Person{依撒格}受孕和\Place{索多玛}倾覆，直到\Place{耶路撒冷}的大规模建造，共有八百九十六年。从\Place{耶路撒冷}的大规模建造直到\Place{耶路撒冷}的毁灭，共有四百二十五年。从\Place{索多玛}倾覆之时直到\Place{耶路撒冷}变为荒废，所有年数的总和是一千三百二十一年。这就是\Place{索多玛}及其女儿在\Place{耶路撒冷}之先荒废的所有年数。那比\Place{耶路撒冷}更为正义的城，至今仍无人居住。所以，从\Place{索多玛}倾覆直到\Person{马其顿的斐理伯}之子\Person{亚历山大}的帝国第六百五十五年，年数的总和是二千二百七十六年。从巴比伦人使\Place{耶路撒冷}变为荒废直到如今，共有九百五十五年。\Place{耶路撒冷}在巴比伦人使其变为荒废之后，曾有人居住，就是在\Person{达尼尔}所证言的那七十个星期期间。然后，它在最后一次毁灭中被罗马人变为荒废，并且永不再有人居住，因为\BibleQuote{它必将荒凉，直到所定之事的结局。} \BibleRef{达 9:27} 这样，\Place{耶路撒冷}前次和后次的荒凉年数，总共四百六十五年，从中减去\Place{巴比伦}的七十年，就是三百九十五年。\par

我写这全部论证给你们，是因为犹太人自傲，（说）\Quote{已向我们订立盟约，我们必将被聚集。} 因为如果\Place{索多玛}的罪恶还不如\Place{耶路撒冷}的大，尚至今无人居住，而我们这样说，它将永远不得恢复，那么罪恶比\Place{索多玛}及其女儿更大的\Place{耶路撒冷}，又怎能恢复呢？至于\Place{索多玛}，\Person{天主}已经二千二百七十六年没有怜悯她；那么我们能说祂会怜悯\Place{耶路撒冷}吗？因为按照上文所写的计算，从她沦为荒废之日直到如今，只有三百九十五年。但关于他所说的，\BibleQuote{索多玛和她的女儿必如从前一样得业} \BibleRef{则 16:55}，以及论及\Place{耶路撒冷}他所说的，\BibleQuote{你和你的女儿必成为如从前一样}，这就是这段话的力量：它们将永远不会再有人居住；因为\Person{主}对祂所发怒的那地也如此诅咒说：——\BibleQuote{这地必不耕种，不生产，不长一草一木，却要像\Place{索多玛}和\Place{哈摩辣}一样，就是\Person{主}向之发怒、不予平息的那地。} \BibleRef{申 29:23} 因此，我的听者，你要确信：\Place{索多玛}和她的女儿将永无人居住；她们将如从前一样，就是说，如同那时她们尚未有人居住，如同\Person{主}向他们发怒、未予平息之时那样。\Place{耶路撒冷}和她的女儿将如从前一样，（即）如同先前阿摩黎人的山那荒凉之时，当时\Person{亚巴郎}在其上建筑祭台，将儿子\Person{依撒格}缚于其上；又如同\Person{达味}从耶步斯人\Person{敖尔难}购买禾场，并在那里建筑祭台时的荒凉一样。因为你们要思考观察：\Person{亚巴郎}献其子的这座山，就是\Place{耶步斯}山，即\Place{耶路撒冷}。\Person{达味}从\Person{敖尔难}购买的禾场，就是建筑圣殿之处。因此，\Place{耶路撒冷}必如从前一样荒凉。你要思考：当\Person{厄则克耳}预言这段话时，\Place{耶路撒冷}仍处于她的伟大中，其中的人正背叛\Place{巴比伦}王。\Person{先知}所说的，是以愤怒和谴责针对\Place{耶路撒冷}。\par

6. 我的听者，要思量且观察：假如天主给\Place{索多玛}和它的同伴留下了希望，祂就不会用烈火与硫磺倾覆它们——那是世界末日的标记——反而会把它们交于某个王国，使之受罚。正如经上记载：当\Person{耶肋米亚}使列国和王国喝了愤怒之杯后，论到每一座城说，在它们喝过这杯之后，\BibleQuote{\Emph{我必转回} \Emph{\Place{厄蓝}、\Place{提洛}、\Place{漆冬}、\Place{阿孟}子孙、\Place{摩阿布}}和\Emph{\Place{厄东}}的俘虏。} 论到这些王国各自，他说：\BibleQuote{\Emph{在末后的日子，我必转回她的俘虏。}} 如今我们见到，\Place{提洛}在\Emph{漂流七十年}之后有人居住且繁荣，又在她得了她卖淫的酬报，\Emph{与一切王国行淫}之后。她\Emph{拿着琴，弹得美妙，增添了众多音乐。} \Place{厄蓝}地区也有人居住且繁荣。至于\Place{巴比伦}，\Person{耶肋米亚}说：\BibleQuote{\Emph{巴比伦必倾覆，不得再起。}} \BibleRef{耶 51:64} 看！直到今日它仍处于荒凉，并将永远如此。关于\Place{耶路撒冷}，他也说：\BibleQuote{\Emph{\Place{以色列}的童女要跌倒，不得再起；她被丢弃在地，无人将她扶起。}} \BibleRef{亚 5:1} 因为\Person{耶肋米亚}关于\Place{巴比伦}所说的预言若是真实的，那么关于\Place{耶路撒冷}的预言也是真实可信的。\Person{依撒意亚}对\Place{耶路撒冷}说：\BibleQuote{\Emph{我必不再向你发怒，也不再责备你。}} \BibleRef{依 54:9} 确实，祂必不再向她发怒，也不永远责备她；因为荒凉中的事物，祂不责备，她也不会惹祂发怒。\par

7. 至于那些指责我们的人，他们\Quote{你们受迫害，却未得拯救}，他们自己应感到羞愧，因为他们每每受迫害，甚至在得救之前多年都在受迫害。他们在\Place{埃及}被迫服事达二百二十五年。在\Person{巴辣克}和\Person{德波辣}的时代，\Place{米德杨人}使\Place{以色列}服役。\BibleRef{民 4:2} 在\Person{厄胡得}的时代，\Place{摩阿布人}统治他们；\BibleRef{民 3:12} 在\Person{依弗大}的时代，\Place{阿孟人}；\BibleRef{民 11:5} 在\Person{三松}的时代，\Place{培肋舍特人}，\BibleRef{民 13:1} 又在\Person{厄里}和\Person{撒慕尔}先知的时代；\BibleRef{撒上 4:1} 在\Person{阿哈布}的时代，\Place{厄东人}；在\Person{希则克雅}的时代，\Place{亚述人}。\BibleRef{列上 20:11} \Place{巴比伦}王将他们从本土拔除，并驱散他们；他多次试探和迫害他们之后，他们仍不改正，正如祂对他们说的：\BibleQuote{\Emph{我徒然打击你们的子女，因为他们不接受惩戒。}} \BibleRef{耶 2:30} 祂又说：\BibleQuote{\Emph{我清除了先知，并藉我口之言将他们杀戮。}} \BibleRef{欧 6:5} 祂对\Place{耶路撒冷}说：\BibleQuote{\Emph{耶路撒冷啊，当藉苦难和鞭笞受教导，免得你的生命离开你。}} \BibleRef{耶 6:7} 但是他们离弃了祂，敬拜偶像，正如\Person{耶肋米亚}论到他们说：\BibleQuote{\Emph{你们到遥远的岛屿去，派人前往\Place{刻达尔}，仔细察考，看看有没有这样的事：列国是否更换了自己的神明——那并非真神的神明。但我的百姓却以无益之物换取了我的光荣。诸天哪，你们要因此惊愕，要颤抖，要极其惊恐，这是主说的；因为我的百姓做了两件恶事：他们离弃了我这活水的泉源，各自挖掘了蓄水池，是破裂不能存水的蓄水池。}} \BibleRef{耶 2:10} 因为破裂的蓄水池就是对雕像和偶像的畏惧。祂呼唤诸天惊愕，因为他们敬拜天上的万象。而诸天将受罚，正如所记：\BibleQuote{\Emph{它们将被卷起如书卷，其万象都要凋零。}} \BibleRef{依 34:4}\par

8. 我所写给你们的一切论述，我所爱的，从一开始，都是因为犹太人指责了我们百姓的子女；但现在，就我所能领悟的，我将教导你们关于受迫害者：他们已获得大赏报，而迫害者终归于耻笑和蔑视。\par

9. \Person{雅各伯}遭受迫害，而\Person{厄撒乌}是迫害者。\Person{雅各伯}领受了祝福和长子名分，\Person{厄撒乌}却被剥夺了这两样。\Person{若瑟}遭受迫害，他的兄弟们是迫害者；\Person{若瑟}被高举，迫害他的人在他面前下拜，如此，他的梦和他的异象都得以应验。遭受迫害的\Person{若瑟}是受迫害的\Person{耶稣}的预像。他的父亲给\Person{若瑟}穿上一件彩衣；而祂的圣父给\Person{耶稣}穿上从童贞女（所取）的身体。他的父亲爱\Person{若瑟}胜过爱他的兄弟们，而\Person{耶稣}是祂圣父的钟爱之子。\Person{若瑟}看见异象，做了梦；\Person{耶稣}则应验了异象和众先知。\Person{若瑟}与他的兄弟们同是牧者；而\Person{耶稣}是牧者的首领。当他的父亲打发\Person{若瑟}去看望他的兄弟们时，他们见他来了，就密谋杀害他；当圣父打发\Person{耶稣}去看望祂的兄弟时，他们说：\BibleQuote{这是继承人；来，我们杀掉他！}\BibleRef{玛 21:38} 他的兄弟们将\Person{若瑟}投入坑中；祂的兄弟们则将\Person{耶稣}带入阴府。\Person{若瑟}从坑里上来，\Person{耶稣}从阴府中复活。\Person{若瑟}从坑里上来后，就拥有了管辖他兄弟们的权柄；\Person{耶稣}从阴府中复活后，圣父赐给祂一个伟大的、超绝的名号，\BibleRef{斐 2:9} 为使祂的兄弟们服事祂，使祂的仇敌伏在祂的脚下。\Person{若瑟}被兄弟们认出后，他们都惊惶、畏惧，并因他的尊大而惊奇；将来，当\Person{耶稣}在末日来临时，当祂在威荣中显现时，祂的兄弟们将在祂面前惊惶、畏惧、丧胆，因为他们曾将祂钉死在十字架上。再者，\Person{若瑟}因\Person{犹大}的计谋被卖到\Place{埃及}；而\Person{耶稣}则藉\Person{犹达斯依斯加略}的手被交付给犹太人。他们卖\Person{若瑟}的时候，他对他兄弟们一言不答；\Person{耶稣}也不言语，对审判祂的判官也不做任何答辩。他的主人不义地将\Person{若瑟}交付监牢；祂的同乡不义地判了\Person{耶稣}的罪。\Person{若瑟}交出了两件衣裳，一件交在他兄弟们的手中，一件交在他主人妻子的手中；\Person{耶稣}交出了自己的衣裳，被兵士们分开。\Person{若瑟}三十岁时，侍立在\Person{法郎}面前，成为\Place{埃及}的主宰；\Person{耶稣}年约三十岁时，来到\Place{约旦河}接受圣洗，领受圣神，出去宣讲。\Person{若瑟}用粮食养育了\Place{埃及}；\Person{耶稣}用生命之粮养育了全世界。\Person{若瑟}娶了邪恶而不洁的司祭的女儿为妻；\Person{耶稣}则迎娶了由不洁的外邦人中（所选取）的教会作自己的净配。\Person{若瑟}在\Place{埃及}逝世并安葬；\Person{耶稣}在\Place{耶路撒冷}逝世并安葬。\Person{若瑟}的骨骸由他的兄弟们从\Place{埃及}带回；\Person{耶稣}则由圣父从阴府中复活，并将祂的身体连同祂一起提升到天堂，完全无损。\par

10. \Person{梅瑟}也遭受迫害，如同\Person{耶稣}遭受迫害一样。\Person{梅瑟}出生时，他们将他藏起来，免得他被迫害者杀害。\Person{耶稣}出生时，他们带着祂逃往\Place{埃及}，免得迫害祂的\Person{黑落德}杀害祂。在\Person{梅瑟}出生的日子，婴儿常被投溺在河中；在\Person{耶稣}出生时，\Place{白冷}及其周围的婴孩则被杀害。天主对\Person{梅瑟}说：\BibleQuote{那些图谋杀害你生命的人已经死了}\BibleRef{出 4:10}；天使在\Place{埃及}对\Person{若瑟}说：\BibleQuote{起来，带着孩子，往以色列地去，因为那些企图杀害孩子性命的人已经死了}\BibleRef{玛 2:20}。\Person{梅瑟}带领自己的百姓脱离\Person{法郎}的奴役；\Person{耶稣}则解救万民脱离\Person{撒殚}的奴役。\Person{梅瑟}在\Person{法郎}的宫中长大；\Person{耶稣}在\Person{若瑟}带着祂逃往\Place{埃及}时，在\Place{埃及}长大。当\Person{梅瑟}在河面上漂流时，\Person{米黎盎}站在河岸上；在天使\Person{加俾额尔}向\Person{玛利亚}预报之后，\Person{玛利亚}生下了\Person{耶稣}。当\Person{梅瑟}祭献羔羊时，\Place{埃及}的长子都被杀了；当他们将\Person{耶稣}这真羔羊钉在十字架上时，杀害祂的子民因杀害祂而灭亡。\Person{梅瑟}为他的百姓降下玛纳；\Person{耶稣}则将祂的圣体赐给万民。\Person{梅瑟}用木头使苦水变甜；\Person{耶稣}则藉祂的十字架，即祂被钉的木架，使我们的苦变为甘甜。\Person{梅瑟}为他的百姓颁布法律；\Person{耶稣}则将祂的盟约赐给万民。\Person{梅瑟}伸出手臂战胜了\Person{阿玛肋克}；\Person{耶稣}则以祂十字架的标记战胜了\Person{撒殚}。\Person{梅瑟}从磐石中为他的百姓引出水来；\Person{耶稣}则派遣\Person{西满刻法}（\OriginalTerm{磐石}{rock}）到万民中传扬祂的教义。\Person{梅瑟}揭去脸上的帕子与天主交谈；\Person{耶稣}则揭去万民脸上的帕子，使他们能聆听并接受祂的教义。\Person{梅瑟}将手覆在他所派遣的使者（宗徒）身上，他们便领受了司祭职；\Person{耶稣}将手覆在祂的宗徒身上，他们便领受了圣神。\Person{梅瑟}登上山并在那里去世；\Person{耶稣}升了天，坐在祂圣父的右边。\par

11. \Person{农的儿子若苏厄}也曾受迫害，就如我们的救赎主\Person{耶稣}受迫害一样。\Person{农的儿子若苏厄}被不洁的外邦人迫害；而我们的救赎主\Person{耶稣}被愚昧的民众迫害。\Person{农的儿子若苏厄}从迫害者手中夺回产业，交给自己的子民；而我们的救赎主\Person{耶稣}从祂的迫害者手中夺回产业，交给外邦民族。\Person{农的儿子若苏厄}使太阳在天空停住，并报复那些迫害他的民族；而我们的救赎主\Person{耶稣}使太阳在正午落下，为使那钉死祂的迫害民众蒙羞。\Person{农的儿子若苏厄}将产业分给自己的子民；而我们的救赎主\Person{耶稣}应许将生命之地赐给外邦。\Person{农的儿子若苏厄}保存了妓女\Person{辣哈布}的性命；而我们的救赎主\Person{耶稣}将虽被偶像崇拜的淫乱所玷污的教会聚集起来，并赐予生命。\Person{农的儿子若苏厄}在第七天推倒并拆毁\Place{耶里哥}的城墙；而我们的救赎主\Person{耶稣}，在祂的第七天，在安息的安息日，这个世界将要崩溃与毁灭。\Person{农的儿子若苏厄}用石头砸死了\Person{阿干}，因为他偷取了当灭之物；而我们的救赎主\Person{耶稣}将\Person{犹达斯}从他的门徒与朋友中分离出来，因为他偷取了穷人的钱。\Person{农的儿子若苏厄}临死时，在他的子民中立下了一项见证；而我们的救赎主\Person{耶稣}，在祂被接升天时，在祂的宗徒中立下了一项见证。\par

12. \Person{依弗大}也曾受迫害，就如\Person{耶稣}受迫害一样。\Person{依弗大}被自己的弟兄赶出父家；而\Person{耶稣}被自己的弟兄赶逐、举高并钉在十字架上。\Person{依弗大}虽受迫害，却兴起作了他子民的领袖；\Person{耶稣}虽受迫害，却兴起成了万民的君王。\Person{依弗大}许下誓言，并将自己的首生女儿献为祭；而\Person{耶稣}被作为祭品，为所有外邦人奉献给圣父。\par

13. 同样，\Person{达味}遭受迫害，如\Person{耶稣}遭受迫害一样。\Person{达味}由\Person{撒慕尔}傅油为王，代替犯罪的\Person{撒乌耳}；而\Person{耶稣}由\Person{若翰}傅油为大司祭，代替那执守法律的司祭们。\Person{达味}在受傅后遭受迫害；\Person{耶稣}也在受傅后遭受迫害。\Person{达味}起初只统治一个支派，然后统治全以色列；而\Person{耶稣}从起初统治少数信祂的人，最终祂将统治全世界。\Person{撒慕尔}傅油\Person{达味}时，达味年三十岁；\Person{耶稣}约三十岁时，由\Person{若翰}举行了覆手礼。\Person{达味}娶了国王的两个女儿；\Person{耶稣}娶了两位国王的两个女儿，即选民会众和外邦会众。\Person{达味}以善报恶于仇敌\Person{撒乌耳}；\Person{耶稣}教训说：\Emph{要为你们的仇敌祈祷}。\BibleRef{路 6:28} \Person{达味}是合乎天主心意的人；\BibleRef{撒上 13:14} 而\Person{耶稣}是天主子。\Person{达味}由迫害者\Person{撒乌耳}手中得了王位；\Person{耶稣}由迫害者以色列手中得了王位。\Person{达味}为仇敌\Person{撒乌耳}逝世哀哭作挽歌；\Person{耶稣}为\Place{耶路撒冷}——祂的迫害者，即将被毁——而哀哭。\Person{达味}将王位传给\Person{撒罗满}，归到自己祖先那里去了；\Person{耶稣}将钥匙交给\Person{西满}，然后升天，回到派遣祂的那位那里。因\Person{达味}的缘故，罪过得赦于他的后裔；因\Person{耶稣}的缘故，罪过得赦于外邦。\par

14. 同样，\Person{厄里亚}遭受迫害，如\Person{耶稣}遭受迫害一样。女凶手\Person{依则贝耳}迫害\Person{厄里亚}；而迫害人的杀人会众迫害\Person{耶稣}。\Person{厄里亚}因以色列的罪，使天闭塞不下雨；\Person{耶稣}因民众的罪，借着自己的来临，使圣神离开先知们。\Person{厄里亚}除灭了\Person{巴耳}的先知们；\Person{耶稣}践踏了\Person{撒殚}及其军旅。\Person{厄里亚}复活了寡妇的儿子；\Person{耶稣}复活了寡妇的儿子，以及\Person{拉匝禄}和会堂长的女儿。\Person{厄里亚}用少许饼养活寡妇；\Person{耶稣}用少许饼饱饫了数千人。\Person{厄里亚}乘旋风升天去了；我们的救赎主\Person{耶稣}升了天，坐在圣父的右边。\Person{厄里叟}领受了\Person{厄里亚}的神恩；\Person{耶稣}向宗徒们的面上嘘了气。\par

15. 同样，\Person{厄里叟}遭受迫害，如\Person{耶稣}遭受迫害一样。\Person{厄里叟}被\Person{阿哈布}的儿子，即那凶手的儿子迫害；\Person{耶稣}被凶杀的民众迫害。\Person{厄里叟}讲预言，\Place{撒玛黎雅}便有了丰裕；\Person{耶稣}说：\Emph{吃我的肉，喝我的血的人，必得永生。} \Person{厄里叟}用少许饼饱饫了一百人；\Person{耶稣}用五个饼饱饫了四千男人，还有妇女和儿童。\Person{厄里叟}使水变成油；\Person{耶稣}使水变成酒。\Person{厄里叟}救寡妇脱离债主；\Person{耶稣}救负债的外邦脱离罪债。\Person{厄里叟}使铁浮起，使木头沉下；\Person{耶稣}将我们里面沉下的提上来，将轻浮的沉下去。一个死人（被放在）\Person{厄里叟}的骨骸上，便活了；所有死在罪恶中的外邦人，被放在\Person{耶稣}的骨骸上，也都活了。\par

16. \Person{希则克雅}也曾被迫害，如同\Person{耶稣}被迫害一样。\Person{希则克雅}被迫害，并被他敌人\Person{散乃黑黎布}辱骂；\Person{耶稣}也被愚昧的百姓辱骂。\Person{希则克雅}祈祷，并战胜了他的仇敌；而藉着\Person{耶稣}的十字架，我们的\OriginalTerm{仇敌}{Adversary}被战胜了。\Person{希则克雅}是全\Place{以色列}的君王；\Person{耶稣}是万邦的君王。因为\Person{希则克雅}患病，太阳向后倒退；因为\Person{耶稣}受苦，太阳失去了光辉。\Person{希则克雅}的敌人成为死尸；而\Person{耶稣}，祂的仇敌将被摔在祂的脚下。\Person{希则克雅}出自\Person{达味}家族；而\Person{耶稣}，按肉身说，是\Person{达味}的儿子。\Person{希则克雅}说：—— \BibleQuote{\Emph{在我的日子，必有和平与真理}}；\BibleRef{列下 20:19}，\Person{耶稣}对祂的门徒说：—— \BibleQuote{\Emph{我将我的平安留给你们}}。\BibleRef{若 14:27} \Person{希则克雅}祈祷，他的病得了医治；\Person{耶稣}祈祷，从\OriginalTerm{死者之境}{abode of the dead}中复活了。\Person{希则克雅}从病中起来后，增加了年岁；\Person{耶稣}复活后，获得了极大的光荣。\Person{希则克雅}在生命延长之后，死亡辖制了他；但\Person{耶稣}，在祂复活之后，死亡将永远不再辖制祂。\par

17. \Person{约史雅}也曾被迫害，如同\Person{耶稣}被迫害一样。\Person{约史雅}被迫害，\Person{跛子法郎}杀了他；\BibleRef{列下 23:29}，\Person{耶稣}被迫害，那些因罪而瘸腿的百姓杀了祂。\Person{约史雅}洗净了\Place{以色列}地的不洁；\Person{耶稣}洗净并除去了普世的不洁。\Person{约史雅}尊崇并光荣了他天主的圣名；\Person{耶稣}却说：—— \BibleQuote{\Emph{我已光荣了（祂的名），还要再光荣}}。\BibleRef{若 12:28} \Person{约史雅}因\Place{以色列}的邪恶\Emph{撕裂了自己的衣服}；\BibleRef{列下 22:11} \Person{耶稣}因百姓的邪恶\Emph{撕裂了圣殿的帐幔}。\BibleRef{玛 27:51} \Person{约史雅}说：—— \BibleQuote{\Emph{将要临于这百姓的义怒何等大}}；\Person{耶稣}却说：—— \BibleQuote{\Emph{将要临于这百姓的义怒，他们必倒在刀剑之下}}。\BibleRef{路 21:23} \Person{约史雅}从圣殿中逐出不洁；\Person{耶稣}从祂父的殿中逐出不洁的商贩。为了\Person{约史雅}，\Place{以色列}的女儿们哀哭恸哭，如\Person{耶肋米亚}所说：—— \BibleQuote{\Emph{以色列的女儿啊，为约史雅哭泣}}；而为\Person{耶稣}，\Place{以色列}的女儿们哭泣哀恸，如\Person{匝加利亚}所说：—— \BibleQuote{\Emph{大地要哀恸，各族一起哀恸。}}\par

18. \Person{达尼尔}也曾被迫害，如同\Person{耶稣}被迫害一样。\Person{达尼尔}被\Place{加色丁}人，即异民的会众迫害；\Person{耶稣}也被\Place{犹太人}，即恶人的会众迫害。\Place{加色丁}人控告\Person{达尼尔}；\Place{犹太人}在总督前控告\Person{耶稣}。他们把\Person{达尼尔}投入狮子圈，但他得蒙解救，毫无受伤地从其中出来；他们将\Person{耶稣}投入\OriginalTerm{死者之境}{abode of the dead}的深渊，但祂上升了，死亡不能辖制祂。关于\Person{达尼尔}，他们以为他落入圈中就不复出来；关于\Person{耶稣}，他们说：\BibleQuote{\Emph{祂既然倒了，就必不再起来}}。针对\Person{达尼尔}，吞噬毁灭的狮子的口被封闭；针对\Person{耶稣}，吞噬毁灭的死亡之口被封闭。他们封了\Person{达尼尔}的圈，并严加看守；他们也严加看守\Person{耶稣}的坟墓，正如他们所说：\BibleQuote{\Emph{派卫兵看守坟墓}}。\BibleRef{玛 27:64} 当\Person{达尼尔}上来时，控告他的人蒙受了羞辱；当\Person{耶稣}复活时，所有钉死祂的人都蒙受了羞辱。审判\Person{达尼尔}的君王因控告他的\Place{加色丁}人的恶行而极其忧闷；\BibleRef{达 6:14}，审判\Person{耶稣}的\Person{比拉多}极其忧闷，因为他知道\Place{犹太人}是由于嫉妒才控告了祂。\BibleRef{玛 28:18} 因\Person{达尼尔}的祈祷，他百姓的俘虏从\Place{巴比伦}上来；\Person{耶稣}因祈祷，使万民的俘虏归回。\Person{达尼尔}解释了\Person{拿步高}的异象和梦；\Person{耶稣}阐释并解释\WorkTitle{法律}和\WorkTitle{先知}的预像。当\Person{达尼尔}解释了\Person{贝尔沙匝}的异象，他获得了国度的三分之一权柄；当\Person{耶稣}完成了预像和先知，祂的父将天上地下的一切权柄交给了祂。\Person{达尼尔}看见了奇迹并说出隐秘；\Person{耶稣}启示隐秘并应验了经上所记。\Person{达尼尔}为他的百姓被带去作质；\Person{耶稣}的身体为万民作了质。因\Person{达尼尔}，君王平息了对\Place{加色丁}人的义怒，使他们未被杀；因\Person{耶稣}，祂的父平息了对万民的义怒，使他们未因己罪被杀死亡。\Person{达尼尔}向君王请求，君王使他兄弟有权掌管\Place{巴比伦}省的事务；\BibleRef{达 2:49} \Person{耶稣}向天主请求，祂将权柄赐给了祂的兄弟，即门徒们，去掌管\Person{撒殚}及其军旅。\Person{达尼尔}论及\Place{耶路撒冷}说，直到所定的事成就，她必荒凉；\Person{耶稣}论及\Place{耶路撒冷}说：\BibleQuote{\Emph{在她内决不留一块石头在另一块石头上，因为她不认识她蒙受的眷顾的日子}}。\BibleRef{路 19:44} \Person{达尼尔}预见为他百姓所剩余的星期；\Person{耶稣}来临并予以应验。\par

19. \Person{哈纳尼雅}和他的兄弟们也曾受迫害，如同耶稣受迫害一样。\Person{哈纳尼雅}和他的兄弟们被\Person{拿步高}迫害；而耶稣被犹太人这百姓迫害。\Person{哈纳尼雅}和他的兄弟们被投入火窑中，烈火却如甘露般清凉，临于义人身上。耶稣也降入黑暗之处，冲破其门户，带出其中的囚徒。\Person{哈纳尼雅}和他的兄弟们从火窑中上来，火焰烧灭了控告他们的人；而耶稣复活，从黑暗之中上来，祂的控告者和钉死祂的人，终将在火焰中被焚烧。当\Person{哈纳尼雅}和他的兄弟们从火窑中上来时，\Person{拿步高}王战栗且惊愕；当耶稣从死者居所中复活时，钉死祂的那百姓惊恐战栗。\Person{哈纳尼雅}和他的兄弟们不敬拜\Place{巴比伦}王的肖像；耶稣也约束万民，不敬拜无生命的偶像。因\Person{哈纳尼雅}和他的兄弟们，\BibleQuote{各民族与各语言都颂扬拯救他们脱离烈火的天主}：\BibleRef{达 3:28} 同样，因着耶稣，万民与万语必将颂扬那拯救祂圣子、使祂未见腐朽的天主。在\Person{哈纳尼雅}和他兄弟们的衣服上，烈火无能为力；同样，在信从耶稣的义人的身体上，末日之时，烈火也必将无能为力。\par

20. \Person{摩尔德开}也曾受迫害，如同耶稣受迫害一样。\Person{摩尔德开}受恶人\Person{哈曼}迫害；耶稣受悖逆的百姓迫害。\Person{摩尔德开}藉着祈祷拯救他的百姓脱离\Person{哈曼}之手；耶稣藉着祈祷拯救祂的百姓脱离撒殚之手。\Person{摩尔德开}被救脱离迫害者之手；耶稣被救脱离迫害者之手。因为\Person{摩尔德开}穿上苦衣坐下，他拯救了\Person{艾斯德尔}和他的百姓脱离刀剑；因为耶稣穿上肉身并受到光照，祂拯救了教会及其子女脱离死亡。因着\Person{摩尔德开}，\Person{艾斯德尔}得君王欢心，得以进来取代不愿遵行王命的\Person{瓦市提}而坐；同样，因着耶稣，教会得\Person{天主}欢心，得以进入君王面前，取代不遵行祂旨意的会众。\Person{摩尔德开}劝诫\Person{艾斯德尔}，应与她的侍女们禁食，好使她与她的百姓脱离\Person{哈曼}之手；耶稣也劝诫教会及其子女（禁食），好使她与她的子女脱离义怒。\Person{摩尔德开}获得了迫害者\Person{哈曼}的尊荣；耶稣从祂的圣父那里获得了大荣耀，取代那些属于愚昧百姓的迫害者。\Person{摩尔德开}践踏在迫害者\Person{哈曼}的颈项上；至于耶稣，祂的仇敌将被置于祂的脚下。在\Person{摩尔德开}面前，\Person{哈曼}宣告说：\BibleQuote{凡君王愿意显扬的人，就当这样对待他}；\BibleRef{艾 6:11} 至于耶稣，宣讲祂的人从迫害祂的百姓中出来，他们说：—— \BibleQuote{这人真是天主子。} \BibleRef{玛 27:54} \Person{摩尔德开}的血要向\Person{哈曼}及其众子追讨；而\BibleQuote{耶稣的血}，迫害祂的人\BibleQuote{归在我们和我们的子孙身上}。\BibleRef{玛 27:25}\par

21. 我所写给你们的这些纪念，我亲爱的，论及受迫害的耶稣，以及受迫害的义人，乃是为使今日为那位受迫害的耶稣的缘故而受迫害的人能得到安慰，因为祂已为我们写下，并亲自安慰了我们；祂曾说：—— \BibleQuote{他们若迫害了我，也要迫害你们；正因你们不属于世界，就如我不属于世界一样，所以他们要迫害你们。} 祂先前也为我们写道：—— \BibleQuote{你们的父母、兄弟、亲族要交出你们；你们要因我的名被众人憎恨。} \BibleRef{路 21:16} 祂又教导我们：—— \BibleQuote{当人押送你们到会堂，到官长及有权势的人前，以及到掌管世界的君王面前时，不要预先思虑怎样申辩，或申辩什么；因为我必赐给你们口才和智慧，是你们的仇敌所不能抵抗及辩驳的。因为说话的不是你们，而是你们父的圣神在你们内说话。} 这圣神曾藉\Person{雅各伯}的口，向迫害他的\Person{厄撒乌}说话；也是那智慧的圣神，曾藉受迫害的\Person{若瑟}的口，在法郎面前说话；也是那藉着\Person{梅瑟}的口，在\Place{埃及}地所行一切奇事中说话的圣神；也是当\Person{梅瑟}按手在\Person{农}的儿子\Person{若苏厄}身上时，赐给他的智慧之神，使那些迫害他的民族在他面前全然消灭；也是那藉受迫害的\Person{达味}的口吟咏圣咏的圣神，达味以此吟唱圣咏，安抚迫害者\Person{撒乌耳}脱离恶神；也是那充满\Person{厄里亚}，并藉着他斥责迫害者\Person{依则贝耳}和\Person{阿哈布}的圣神；也是那在\Person{厄里叟}内说话，并向迫害他的君王预言并昭示未来一切事的圣神；也是那在\Person{米加雅}口中炽热，当他斥责迫害者\Person{阿哈布}说：—— \BibleQuote{如果你能平安回来，主就没有藉我发言。} \BibleRef{列上 22:28} 的圣神；也是那在\Place{巴比伦}地保全\Person{达尼尔}及其兄弟的圣神；也是那拯救被掳之地的\Person{摩尔德开}和\Person{艾斯德尔}的圣神。\par

22. 我亲爱的，请听这些\OriginalTerm{殉道者}{martyr}、\OriginalTerm{宣信者}{confessor}和受迫害者的名字。亚伯尔被杀害，他的血从地下喊冤。雅各伯受迫害，逃亡而成为流亡者。若瑟受迫害，被卖并被投入坑中。梅瑟受迫害，逃往米德杨。农的儿子若苏厄受迫害，投入战争。依弗大、三松、基德红和巴辣克，这些人也受迫害。这些人正是真福宗徒所说的：\BibleQuote{\Emph{时间不够我叙述他们的胜利。}}\BibleRef{希 11:32}达味也受撒乌耳的迫害，他行走在\BibleQuote{\Emph{山岭、洞窟和地穴中}}。\BibleRef{希 11:38}撒慕尔也受迫害，为撒乌耳而哀哭。此外，希则克雅受迫害，被苦难束缚。厄里亚受迫害，行走在旷野。厄里叟受迫害，成为流亡者；米加雅受迫害，被投入监牢。耶肋米亚受迫害，被投入污泥坑中。达尼尔受迫害，被投入狮子坑中。阿纳尼雅和他的兄弟们也受迫害，被投入火窑。摩尔德开、艾斯德尔和他们同族的子民，受哈曼的迫害。犹达玛加伯和他的弟兄受迫害，他们也忍受羞辱。那有福的妇人的七个儿子，受严酷鞭打的折磨\BibleRef{加下 7:1}，是\OriginalTerm{宣信者}{confessor}和真正的\OriginalTerm{殉道者}{martyr}；厄肋阿匝尔虽年事已高、年迈老成，却立下崇高榜样，宣认信仰，成为真正的殉道者。\BibleRef{加下 6:18}\par

23. 耶稣的\OriginalTerm{殉道}{martyrdom}伟大而卓越。祂在苦难和\OriginalTerm{宣认}{confession}上超越了所有之前和之后的人。在祂之后，有忠信的殉道者斯德望，被犹太人用石头砸死。西满（伯多禄）和保禄也是完美的殉道者。雅各伯和若望行走在他们师傅基督的足迹上。此后还有其他宗徒在各处宣认信仰，证明是真正的殉道者。至于我们在西方的弟兄们，在戴克里先的时代，整个天主的教会在他们全境遭受了极大的苦难和迫害。教会倾覆并被连根拔除，许多\OriginalTerm{宣信者}{confessor}和\OriginalTerm{殉道者}{martyr}宣认了信仰。在他们受迫害之后，主仁慈地垂顾了他们。在我们的时代，这些事也因我们的罪过而临于我们；但这也是为应验经上所载，正如我们的救赎主所说：\BibleQuote{\Emph{这些事必然发生。}}宗徒也说：\BibleQuote{\Emph{我们也处于这宣认的云彩之下。}}；\BibleRef{希 11:1}，这正是我们的光荣，许多人在其中宣认并被杀害。\par

\SourceRecord{Translated by John Gwynn.；Nicene and Post-Nicene Fathers, Second Series；Vol. 13.；Edited by Philip Schaff and Henry Wace.；Buffalo, NY: Christian Literature Publishing Co.,；1890.；<http://www.newadvent.org/fathers/370121.htm>.}

% structure: p0001 source=h1 target=section reason=page-title

\section{第二十二篇论证（论死亡与末世）}

1. 正直、公义、良善和智慧的人并不畏惧死亡，也不因死亡而战栗，因为他们面前有伟大的希望。他们时时刻刻惦念着死亡、自己的离世，以及亚当子孙将受审判的末日。他们知道，由于亚当违背了诫命，死亡便借着审判的判决掌了权；正如宗徒所说：—— \Emph{从亚当到梅瑟，死亡一直为王，连那些没有犯过与亚当同样罪过的人也在它的权下；这样，死亡就临到了亚当所有的子孙}，正如它临到了亚当一样。那么，死亡怎样从亚当到梅瑟为王呢？显然，当天主为亚当颁布诫命时，曾警告他说：—— \BibleQuote{你吃这棵知善恶树上的果子那一天，必定要死。}\BibleRef{创 2:17}所以他违背了诫命，吃了那树上的果子以后，死亡就统治了他和他所有的后裔。甚至对那些没有犯罪的人，死亡也因亚当违背诫命而统治了他们。\par

2. 他为什么说：—— \Emph{从亚当到梅瑟，死亡为王} 呢？谁这样缺乏知识，竟以为死亡仅仅从亚当到梅瑟统治呢？然而，从这话他应当明白：他说了—— \Emph{它临到了众人}。所以，从梅瑟直到世界末日，它临到了众人。然而，梅瑟宣讲了死亡的国度已被废除。因为当亚当违背了诫命，死亡的判决便临到他的后裔，死亡指望牢牢地束缚死所有人的子孙，并永远作他们的君王。但当梅瑟来临时，他宣讲了复活，死亡便知道它的国度将被废除。因为梅瑟说：—— \BibleQuote{勒乌本应生存，不至死亡，人数也不减少。}\BibleRef{申 33:6} 当圣者从荆棘丛中呼唤梅瑟时，这样对他说：—— \BibleQuote{我是亚巴郎的天主，依撒格的天主，雅各伯的天主。}\BibleRef{出 3:6} 死亡一听见这话，就震动、恐惧、惊惶、不安，并知道它并不能永远作亚当子孙的君王。从它听见天主对梅瑟说：—— \BibleQuote{我是亚巴郎的天主，依撒格的天主，雅各伯的天主} ——那时起，死亡就双手拍击，因为它得知天主是死者与生者的君王，且注定亚当的子孙要从它的黑暗中出来，连同身体复活。还要注意，我们的救赎者耶稣，当祂向撒杜塞人重复这话，而他们正与祂争论有关死人复活时，也这样说过：—— \BibleQuote{天主不是死人的天主，而是活人的天主：所有的人为祂都是活着的。}\BibleRef{路 20:38}\par

3. 天主为让死亡知道，它的权柄并非永远控制普世所有的后裔，就因着哈诺克蒙天主悦纳，将他接去归于自己，使他成为不死者。祂又将厄里亚提升到天上，死亡便不能辖制他。亚纳说：—— \BibleQuote{上主使人死，也使人活；使人降入阴府，也要将人由阴府提出。}\BibleRef{撒上 2:6} 而且，梅瑟犹如出自天主的口说：—— \BibleQuote{我使人死，也使人活。}\BibleRef{申 32:39} 再者，先知依撒意亚也曾说：—— \BibleQuote{您的亡者将再生，他们的尸体将要起立；睡在尘埃中的人们都要苏醒歌咏。}\BibleRef{依 26:19} 死亡听见这一切，不胜惊愕，就坐下哀悼。\par

4. 当那死亡毁灭者耶稣来到，穿上了出自亚当种类的身体，并在祂的身体内被钉十字架，尝了死亡；当（死亡）由此发觉祂已降到它那里时，它便从座位上震动，看见耶稣就惊惶不安；它关上自己的门户，不愿接待祂。然后祂冲破它的门户，进入它内，开始洗劫它一切的财产。当死者在黑暗中看见光明时，便从死亡的束缚中抬起头来，向外观看，看见了默西亚君王的光彩。于是，死亡的黑暗权势都坐下哀悼，因为它从自己的权位上被贬黜了。死亡尝了那致它死命的药剂，它的双手垂了下去，得知死者将复活，并脱离它的统治。当祂借着洗劫死亡的财产而磨难了它时，它便悲泣、痛哭、哀号说：\Quote{从我的领域出去，不要再进来。这个活着进入我领域的是谁呢？} 正当死亡惊恐地喊叫时（因为它看见自己的黑暗开始消散，一些睡了的义人也起来，要与祂一同上升），那时祂使死亡知道，当祂在时期的圆满来到时，要把它权下所有的囚犯都带出来，他们要出去见光明。于是，当耶稣在死者中完成了祂的职事后，死亡把祂送出领域，不容祂留在那里。且要把祂像其他死者一样吞灭，它不以为乐。它对这圣者毫无权柄，祂也不被交给腐朽。\par

5. 当它急切地送祂出去，祂从其领域出来后，便给它留下生命的应许，作为毒素；好使它的权势渐渐消灭。就如人把毒素放在为维持生命而给的食物中，当他发觉自己从食物里中了毒，就把混合了毒素的食物从腹中吐出来；但那药物的毒性仍留在他的肢体里，以致身体的构造渐渐消散败坏。这样，死去的耶稣便使死亡归于无有；因为借着祂，生命得以为王；借着祂，死亡被废除，有话来对它说：—— \BibleQuote{死亡！你的胜利在哪里？}\BibleRef{格前 15:55}\par

6. 因此，你们这些亚当的子孙啊，所有曾受死亡辖制的人，要记念死亡，并记念生命；不要像你们的原祖那样违背诫命。哦，头戴冠冕的君王啊，要记念死亡，死亡将取去你们头上的冠冕，并且死亡要作你们的君王，直到你们为审判而复活的时候。哦，傲慢、自大、骄傲的人啊，要记念死亡，死亡必消灭你们的傲慢，分解肢体，拆散关节，使身体和容貌归于腐朽。高傲的人必被死亡降卑，强横严酷的人必埋没在死亡的黑暗里。他必除去一切骄傲，他们必朽坏，化为尘土，直到审判。哦，富足的人啊，要记念死亡；因为时候将到，你们就近死亡时，你们的财富和产业便毫无用处。他不会在你们面前摆设珍馐，也不会为你们预备丰盛筵席。贪食者惯享甘肥的身体必在那里腐朽。他们必终止宴乐，不再记念。在那里，蛆虫必吞食他们的身体，他们美丽衣裳之上必披上黑暗。他们不思念世界的终结，死亡必在他们降到他那里时使他们惊惶。所以，他们必坐在压迫和死荫里，不再记念这个世界，直到末时来到，他们为审判而复活。哦，贪婪、勒索、抢夺同伴的人啊，要记念死亡，不要增添你们的罪；因为在那个地方，罪人不再悔改；抢夺别人财物的人，必不能保住自己的财物，却要去到那无人使用财富的地方。他必归于虚无，离弃他的尊荣，但他的罪孽要留存到审判的日子。\par

7. 哦，倚靠这个世界的人啊，你们眼中要轻看这个世界；因为你们在这世界是寄居的、是客旅，你们不知道哪一天你们就要被取去离开它。因为死亡会突然来临，把心爱的儿女从父母那里夺去分开，把父母从宝贝孩子那里夺去。死亡为自己带去珍贵的独生子女，他们的父母失去他们，被人藐视。死亡把自己珍贵的朋友分开带走，他们眷爱的人为他们哀哭悲恸。死亡带走并俘虏那些因美貌而被渴慕的人，好羞辱他们的形体，使他们朽坏。那些容貌荣耀的人，死亡把他们带去，他们成为尘土，直到审判。死亡从未婚夫那里带走已订婚的少女，把她们囚禁在他的洞房——他的黑暗之处。死亡带走并分开那些已订婚、已定名许配给童贞女的新郎，童贞女们为他们坐而哀恸。死亡为他自己带走并分开所有美好的青年，他们原以为直到年老也不会看见死亡。死亡为自己带走并聚集那些日子满足的可爱婴孩，他们的父母对他们还未曾厌倦。死亡为自己带走富足的人和奢侈之子；\BibleQuote{他们遗弃财物，如同抛弃海中的波浪。}\BibleRef{德 29:18}死亡为自己带走精巧的工匠，他们曾以其奇妙的作为提升世界。死亡为自己带走精明人和智慧人，他们成为愚拙，不分善恶。死亡为自己带走这世界富有禀赋的人，他们的禀赋被摧毁，永不能建立。死亡为自己带走强壮和大人物，他们的力量被降卑削弱，归于无有。那些自信力量不会削弱的人，在死亡之日，地位比他们低的人却收拾他们的遗体。那些相信死后会得到体面殡葬的人，却遭遇狗来吞食他们。那些相信会葬在自己出生之地的人，却不知道自己将在被掳之地受辱而终，归入坟墓。那些信赖自己产业，以为可以留给子女作产业的人，他们却不知自己将被敌人抢夺。死亡为自己带走勇士和战士，他们曾想蹂躏广大的世界。死亡带走那些以一切美物装饰自己的人，埋葬时却遭遇驴子的埋葬。死亡辖制未出生的胎儿，在他们出生前就俘虏他们。死亡带走那些以盛况受尊崇的人，当他们降到他那里，到黑暗无光的领域，他们就受轻贱。死亡不羞于面对头戴冠冕的君王。死亡不畏惧骄横和蹂躏各地的人。死亡不尊重尊贵人的情面，也不接受富人的贿赂。死亡不藐视穷人，他的灵魂不鄙弃一无所有的人。死亡不尊敬生活豪奢的人，在他那里善与恶无别。他不顾及长者胜过孩童的尊荣。他把明智的主人变为无知，那些曾急促并自忧，在那里与他一同获取财物的人，都被剥夺了收益。死亡带走奴隶及其主人；在那里，主人并不比仆人更受尊重。\BibleQuote{那里有小和大，听不见压迫者的声音。脱离了主人的奴隶}\BibleRef{约 3:18}在那里也不理会从前压迫他的人。死亡捆绑并俘虏看守囚犯的人，以及被囚禁的人。藉着死亡，囚犯得释放，不再害怕迫害他们的人。\par

8. 生活安逸的人害怕死亡；但受苦的人怀着希望，盼望快快被接去。所有富人都因死亡而战栗；但穷人渴望死亡，好从劳苦中得安息。死亡使有权势者想起他时就恐惧；但病人怀着希望仰望他，借着他可以忘记疼痛。再者，小孩子害怕死亡，因为死亡临到时他们就要离开快乐；但年高的老人祈求死亡，就是那些需要日用食粮的人。\par

9. 和平之子记念死亡；他们离弃并除去忿怒和仇恨。他们如同客旅寄居在这世上，并为前头的路程预备行装。他们思念上面的事，默想上面的事，而轻视眼下的事物。他们把财宝送到没有危险的地方，就是没有蛀虫、也没有盗贼的地方。他们住在世上如同寄居者，远方的儿女；盼望被差出这个世界，进入义人的城。他们在寄居之地刻苦己身；在他们被掳的房屋中，他们不被缠累或占据。每天每时，他们都面向上天，要去往他们列祖的安息之所。他们在这世上如同囚犯，又如同王的抵押品被看守着。至终他们在世上没有安息，也不寄望于这世界会存到永远。那些获取财产的人，不会因财产欢乐；那些生儿育女的人，死亡使他们充满忧愁。那些建造城市的人，不会留在其中；而那些为任何事奔忙劳苦的人，与愚昧人毫无分别。无知的人啊，凡信赖这世界的，无论他是谁！\par

10. 亲爱的，你们要记念，在心里比较思量：从前世代的人中，有谁被留在世上活到永远呢？死亡带走了前代的人，包括伟人、有权势者和智慧人。有谁获得了大财产，到他离去的时候能带走它们呢？从地上所聚集的，仍归回地的怀抱；人离开他的财产时，是赤身而去的。智慧人获取财物时，把其中一些预先送到前面去，正如约伯说：\Emph{我的见证人在天上}；又说：\Emph{我的弟兄和朋友都在天主那里}。主也吩咐那些获得财产的人，在天上\Emph{为自己结交朋友}，并且在那里\Emph{积蓄财宝}。\par

11. 明智的经师，你也要记念死亡，免得你心高气傲，以致忘记审判的判词。死亡不放过智慧人，也不看精明人的情面。死亡把明智的经师们带去，使他们忘记所学，直到众义人复活的时候来临。\par

12. 在那里他们将要忘记这个世界。那里他们没有缺乏；他们将以丰盛的爱彼此相爱。他们的身体不再沉重，他们要轻快地飞翔，\Emph{如同鸽子飞向它们的窗户}。\BibleRef{依 60:8} 在他们的思想中不再记念邪恶，他们的心中也不生任何不洁。在那里没有本性的欲望，因为他们已断绝对一切私欲的依赖。他们心中不再起忿怒或邪淫；他们要将所有生发罪过的事物完全除去。他们心中充满彼此的爱火；仇恨完全不会在他们里面扎根。他们无需建造房屋，因为要住在光中，住在圣人们的居所里。他们无需编织的衣裳，因为要以永恒的光为衣。他们无需食物，因为要斜倚在祂的餐桌上，永得滋养。那境界的空气令人愉悦而荣耀，其光辉照耀，美好而喜乐。那里栽植着美丽的树木，果实不落，叶子不枯。它们的枝条华美，芬芳怡人，其滋味永不会使灵魂厌倦。那境界辽阔无垠，但其中的居民看见远处如同近处。在那里产业不会分割，没有人会对同伴说：\Quote{这是我的，那是你的}。他们在那里不会受贪欲的捆绑，也不会在记忆方面迷失。在那里一个人不会格外敬爱邻舍，而是众人以同一方式丰盛地彼此相爱。在那里他们不娶妻子，也不生养儿女；男女之间也没有区别；人人都是他们在天的圣父的儿女；正如先知所说：\Emph{我们岂不是同有一位圣父吗？岂不是同有一位天主创造了我们吗？}\BibleRef{拉 2:10}\par

13. 至于我先前所说的——那里人不娶不嫁，也不分男女——我们的主和祂的宗徒们已经教导了我们。因为主说：\BibleQuote{那堪当得那世界和从死者中复活的人，不娶也不嫁；因为他们不能再死；他们像天上的天使一样，是天主的子女}。\BibleRef{路 20:35} 宗徒也说：\BibleQuote{不再分男或女、为奴或自由，因为你们众人在耶稣基督内已是一个}。\BibleRef{迦 3:28} 因为，关于\Person{厄娃}，为繁衍后代，天主从\Person{亚当}周身取出她，使她成为一切生者的母亲；但在那世界里没有女性；如同在天上也没有女性，没有生育，也没有贪欲的运用\OriginalTerm{贪欲}{concupiscence}。那地方毫无欠缺，只有圆满和成全。老人不再死去，青少年不再衰老。年轻人娶妻生子，就是因着对衰老和死亡的期待，好使父亲去世后，子女可以接替他们。但这一切只在此世有用，因为在那地方，没有匮乏，没有任何欠缺，没有贪欲，没有生育，没有终结，没有衰竭，没有死亡，没有终止，没有衰老。那里没有仇恨，没有忿怒，没有嫉妒，没有厌倦，没有劳苦，没有黑暗，没有黑夜，没有虚假。那地方毫无任何缺乏，却满盈着光明、生命、恩宠、圆满、满足、更新、爱德，以及一切已写下但尚未\Emph{封缄}的美好许诺。因为在那里的是\BibleQuote{眼睛未曾看见，耳朵未曾听见，人心也未曾想到的}，\BibleRef{格前 2:9} 那是无法言说、人不能说出的。宗徒又说：\BibleQuote{那就是天主为爱祂的人所预备的}。\BibleRef{格前 2:9} 纵然人讲论许多，终无法表达。眼睛未曾看见的，他们无法讲述；耳朵未曾听见的，也不宜以任何与耳闻目睹之物相比较的方式谈论。至于那未曾进入人心的，谁敢谈论，好像它与任何人心思量过的事物相似呢？但讲话者这样比拟并称呼那地方是合宜的：天主的居所，生命的所在，成全之处，光明之境，荣耀之地，天主的安息日，休息之日，义人的安息，正直人的喜乐，义人与圣者的住所和居处，我们希望的所在，我们信靠的稳固住所，我们宝藏之处，那地方将平息我们的疲倦，除去我们的苦难，抚慰我们的叹息。我们应当用这些来比拟，也这样称呼那地方。\par

14. 再者，死亡将君王与城镇的建造者带走，他们都曾在辉煌中自仗权势。死亡也不放过列国的王侯。死亡将那些永不满足、不说\Quote{够了}的贪婪者掳掠而去；他以比他们更大的贪婪吞没了他们。死亡将那些掠夺者带走，他们并不因自己所蒙受的仁慈而停止掠夺同伴。死亡将压迫者带走，他们因死亡而止息了不义。死亡将迫害者带走，被迫害者则得享安宁，直到他们也被带往他那里。死亡将那些吞噬同伴的人带走，受践踏、受压迫的人暂时得以休息，直到他们自己也都被带走往那里去。死亡将那些思虑繁多的人带走，他们所思量的一切都消散归于无有。人筹划许多事，死亡忽然临到，他们便被带走；从此再不记得他们所思想的。有人为多年作计划，却不知道他活不过明天。亚当的子孙中，有人自高自大，向同伴夸耀；死亡临到，他的夸耀便归于无有。富人计划增添财产，却不知道他连自己已拥有的也不能长久保留。死亡将一切世人都带到他那里，将他们紧紧捆绑在他的住处，直到审判之时。对于那些没有犯过罪的人，死亡也作了王，这是因亚当为自己的罪所承受的审判的定案。\par

15. 那\OriginalTerm{生命之主}{Life-giver}——毁灭死亡者，将要来临，将死亡的权能从义人和恶人身上除去，归于虚无。死人将以大呼声复活，死亡将被倒空，失去他所有的掳物。亚当的一切子孙都要为审判聚集，各人要去到为他预备的地方。义人的复活者要进入生命，罪人的复活者要被交付于死亡。遵守诫命的义人将要离去，在他们复活的那日，不接近审判；正如\Person{达味}所求的：\Emph{不要将你的仆人带入审判}；他们的主在那日也不会使他们惊惶。\par

16. 你们要记起，宗徒也曾说：\BibleQuote{我们将要审判天使}。\BibleRef{格前 6:3} 我们的主对祂的门徒说：\Emph{你们要坐在十二宝座上，审判以色列家十二支派}。\Person{厄则克耳}论及义人说，\BibleRef{则 23:24} 他们要审判\Person{敖曷拉}和\Person{敖曷里巴}。既然义人要审判恶人，祂已明确指出他们不受审判。至于宗徒们所说\BibleQuote{我们将要审判天使}，你们听着，我要指教你们。那要被宗徒们审判的\Quote{天使}就是那些违犯法律的司祭；正如先知所说：\BibleQuote{司祭的唇舌应保卫智识，人也应从他的口中寻求教训，因为他是万军上主的使者}。\BibleRef{拉 2:7} 那些被称为\Quote{天使}的司祭，就是人们应从其口中寻求教训者，当他们违犯法律时，在末日要被宗徒们和那些遵守法律的司祭们审判。\par

17. \BibleQuote{恶人在审判时不得兴起，罪人也不得在义人的会中。} 正如善工圆满的义人不会进入审判受审，同样，那些罪恶众多、过犯满盈的恶人，也无需进入审判；他们复活后，将直接返回\OriginalTerm{阴府}{Sheol}，就如\Person{达味}所说：\BibleQuote{恶人必将返回阴府，所有忘记天主的民族亦然。} \Person{依撒意亚}也说：\BibleQuote{万民都像桶中的一滴水，如天平上的微尘；海岛犹如一粒沙被抛弃，万国在祂眼中如同虚无。祂视他们如毁灭和刀剑。} 因此，要明白并确信：所有不认识他们的造物主天主的民族，在天主眼中都是虚无的，不会临近审判，一旦复活就将返回阴府。\par

18. 但世上其余被称为罪人的人，将站立审判前受谴责。那些有过轻微过失的，审判者将予以谴责，告知他们违犯了诫命。审判之后，祂将赐予他们生命的产业。要明白，我们的主在福音中已向我们揭示：各人要按照自己的行为获得赏报。那领了钱的，显出了所赚的增值。谁的\OriginalTerm{米纳}{pound}或\OriginalTerm{塔冷通}{talent}赚了十倍的，就获得圆满、毫无欠缺的生命。赚了五倍的，就获得十份之半。一人被赐予十倍权柄，一人五倍。现在细想：五的增值次于十的增值；而索取赏报的工人胜过那些默默接受赏报的。那些整日劳苦的，坦然领受赏报并索取，深信主会加添更多。而那些只工作了一小时的，默默接受，知道是因着恩宠才获得怜悯和生命。罪孽深重的罪人，将由审判之地判定有罪，并进入折磨。从那时起，审判将主宰他们。\par

19. 再听宗徒的话：\BibleQuote{各人要照自己的工夫领受自己的赏报。}\BibleRef{格前 3:8} 劳苦少的，按他的懈怠领受；奋勉多的，按他的勤奋得赏。\Person{约伯}也说：\BibleQuote{天主决不行恶，全能者断不颠倒是非。因为祂照人的行为赏报人，按人所行的报应人。}\BibleRef{约 34:10} 宗徒又说：\BibleQuote{星辰的光芒各不相同，死人的复活也是这样。}\BibleRef{格前 15:41} 因此要明白，即使人们进入永生，赏报仍有高低，荣耀仍有大小，酬报仍有差异。等级高于等级，光的景象亦更美好。太阳超越月亮，月亮大于与她同在的星辰。且看，月亮和星辰都在太阳的权下，它们的光辉被太阳的辉煌吞没。太阳与月亮星辰同具能力，为不废去已与白日分开的黑夜。太阳被造时，即被称为光体。再看，太阳、月亮和星辰都称为光体；但光体超越光体。太阳使月亮的光暗淡，月亮同样使星辰的光幽暗；而星辰的光辉彼此超越。\par

20. 还要从现世的事例中明白：那些辛苦劳作者，以及与人同工的雇工。有些人按日工资雇用同伴，领取他们劳苦的工价；有些人按月受雇，到了约定的时间，计算并领取该时的工资。日工资有别于月工资；年工资更超过月工资。\par

21. 再者，也可从现世的权柄中明白。有人因尽职讨国王欢心，从掌权者那里获得尊荣。一人从国王手中领受冠冕，成为某一地区的总督。另一人的权柄之下，国王安置了城镇；并且在服饰上也优胜于部下。有人领受礼物与赠品，荣耀各有区别。有一人，国王赐予他管理全库房的荣耀；另一人则按较低的身份侍奉国王，其职权仅限于供应日常饮食。\par

22. 论到惩罚，我也说并非人人大同。犯大恶者受大折磨；触犯较轻者受较轻折磨。有些人\BibleQuote{本国的子民，反要被驱逐到外边黑暗里；那里要有哀号和切齿。}\BibleRef{玛 8:12} 其他人则按其所应得被投入火中；因为经上并未说他们会切齿，也没有说那里有黑暗。还有些人将被投入另一处所，在那里\BibleQuote{他们的蛆虫不死，火也不灭；但一切有血肉的，必视他们为可憎之物。}\BibleRef{依 66:24} 还有些人的门将被关上，审判者会对他们说：\BibleQuote{我不认识你们。}\BibleRef{玛 25:12} 考虑一下，正如善行的赏报并非人人一样，恶行也是如此。人受审判并非只有一种方式，而是各人要依照自己的行为获得报应，因为审判者身披正义，不看人的情面。\par

23. 正如我向你们指示有关世界的事，即今世的君王和掌权者赐给下属的尊荣如何一个高过一个；我也曾向你们指示过这有关的事：正如君王有美好的礼物赐给那些他们赏识的人，同样他们也有监狱、锁链和脚镣，即各种不同的束缚。有人以重罪触怒君王，不经审问便被处死。另一人虽然触怒君王，却不至于死；他被加以束缚，直到受审判；并受责罚，君王便宽免了他的过犯。还有一人，是君王所看重的；他被置于监狱之外，享有自由，没有锁链，也没有束缚。被处以死刑的，与受束缚的，是有分别的；而刑罚的轻重，各依其罪过的轻重而定。但请来就我们的救主，祂曾说：\BibleQuote{我父的家里有许多住处。}\BibleRef{若 14:2}\par

24. 我亲爱的，那些理解力低下的人，对我写给你们的这些事辩论说：\Quote{义人将在何处领受美好的赏报？而恶人将在何处受痛苦，按照他们的行为接受惩罚？} 人啊，你这样思想，我要问你，也请你告诉我：为什么死亡称为死亡？为什么阴府\OriginalTerm{阴府}{Sheol}称为阴府？因为经上记载：当科辣黑和他的同党反抗梅瑟时，\BibleQuote{地裂开了口，吞下了他们，他们活活地坠入阴府。}\BibleRef{户 16:32} 因此，那在旷野中张开的，就是阴府的口。达味也说：恶人必归回阴府。我们说，科辣黑和他的同党被吞入的那个阴府，恶人也必被赶回那里去。因为天主有权柄，如果祂愿意，可以在天上赐予生命的产业；如果祂喜欢，也可以在地上赐予。耶稣我们的主曾说：\BibleQuote{神贫的人是有福的，因为天国是他们的。}\BibleRef{玛 5:3} 而对那个与祂同钉十字架、信从祂的人，祂起誓说：\BibleQuote{你今天就要与我一同在伊甸乐园里。}\BibleRef{路 23:43} 宗徒也说：\BibleQuote{当义人复活时，他们将被提升到空中迎接我们的救主。}\BibleRef{得前 4:17} 然而，我们这样说：我们的救主对我们所说的是真实的：\BibleQuote{天地要过去。}\BibleRef{玛 24:35} 宗徒又说：\BibleQuote{所见到的希望，不算是希望。}\BibleRef{罗 8:24} 先知也说：\BibleQuote{诸天要像烟云消散，大地要像衣服变旧；地上的居民也要这样死亡。}\BibleRef{依 51:6} 约伯论及那些长眠的人说：\BibleQuote{直到诸天消逝，他们不会醒来，也不会从他们的睡眠中惊醒。}\BibleRef{约 14:12} 从这些事，你应确信：这大地，就是亚当的子孙被播种于其中的大地，以及覆盖在人上的穹苍\OriginalTerm{穹苍}{firmament}——就是那为分开上层诸天与大地及今生而设置的穹苍——都要过去，变旧，被毁灭。而天主要为亚当的子孙造一件新事，使他们在天国\OriginalTerm{天国}{Kingdom of Heaven}里承受产业。如果祂将产业赐给他们在地上，那也称为天国。如果赐在天上，为祂也是容易的事。因为就如世上的君王，虽然各居自己的地方，但凡是他们权柄所及之处，都称为他们的王国。同样，太阳是安置在天上的光体，凡是它的光芒所及之处，它的权能足以覆盖，无论是在海上或陆上。你再观察：世上的君主也有宴会和享乐，他们不论往哪里去，无论进入何地何种境遇，他们的宴会总随着他们；而且在他们喜欢的地方，他们也能设立监狱。因为太阳在十二小时内环绕，从东方至西方；当它完成了行程，它的光就在夜间隐藏，黑夜并不因它的权能而受到搅扰。而在夜间时辰，太阳转绕着它快速的行程，旋转着开始奔行在它惯常的轨道上。至于与你同在的太阳，你这位聪明人啊，从你的幼年直到老年终了，你也不知道它在夜间如何运行，好绕回它的轨道之处。有必要去探究那些向你隐藏的事吗？\par

25. 我写下这些纪念文字，是为我们的弟兄与所亲爱的、\Emph{天主的教会}的儿女，使这些文字在各地到达他们手中，当他们读到其中内容时，也能在祈祷中记念我的微不足道，并知道我也是个罪人，且有欠缺；但这是我一开始就提出并写下的信德，记在这些（我所写的）章节中。\Emph{信德}是基础，而在\Emph{信德}之上（建立）与之相称的行为。在信德之后，（我写了）\Emph{爱德}有两条诫命。在爱德之后，我写了\Emph{守斋}，连同它的论证与行为。在守斋之后，我写了\Emph{祈祷}，论及其果实与行为。在祈祷之后，我写了\Emph{战争}，以及\Person{达尼尔}所写关于列国的一切。在战争之后，我写了给\Emph{隐修士}的劝勉。在隐修士之后，我写了\Emph{\OriginalTerm{悔改}{Repentance}}。在悔改之后，我写了\Emph{死者的复活}。在死者的复活之后，我写了\Emph{谦逊}。在谦逊之后，我写了\Emph{牧者}，即教师们。在牧者之后，我写了\Emph{割损礼}，犹太人以此自傲。在割损礼之后，我写了\Emph{逾越节}，以及十四日。在逾越节之后，我写了\Emph{安息日}，犹太人为此骄傲。在安息日之后，我写了一份\Emph{劝勉}，因我们时代发生的纷争。在劝勉之后，我写了\Emph{食物}，即犹太人所认为不洁的。在食物之后，我写了\Emph{外邦人}，他们已进入并成为继承人，取代了原先的子民。在外邦人之后，我写并证明了\Emph{天主有子}。在\Emph{天主子}之后，我写了驳斥犹太人，他们毁谤\Emph{童贞}。在关于\Emph{童贞}的辩护之后，我再次写了\Emph{驳犹太人}，他们说：\Quote{我们被预定要聚集在一起。}在那辩护之后，我写了\Emph{施舍给穷人}。在穷人之后，我写了一篇关于\Emph{受迫害者}的论证。在受迫害者之后，我在最后写了\Emph{死亡与末世}。我按照字母表的二十二个字母写了这二十二篇论说。前十篇我写于\Person{马其顿的斐理伯之子亚历山大}之国的第六百四十八年，正如在它们结尾处所记载的。而这十二篇我写于希腊人与罗马人的王国，即\Person{亚历山大}的王国之第六百五十五年，以及波斯王第三十五年。\par

26. 我所写的这些事，都是按照我所能领悟的。但如果有人读到这些讲论，发现其中有些话与自己的想法不合，他不应轻视它们；因为这些章节中所写的一切，不是按照一个人的想法，也不是为了说服某一位读者；而是按照整个教会的想法，为了说服所有信仰的人。如果他以信服的心来读、来听，那很好；如果不是，我应当这样说：我是为那些愿意信服的人写的，不是为嘲弄者写的。再者，如果有读者发现我们所说的某些话，以某种方式表达，而另一位智者以另一种方式表达，不要因此困扰；因为各人都是按照他所能领悟的，向他的听众说话。因此，我，写下这些事的人，即使有些话与别的讲者所说不同，我仍要说：那些智者说得很好，但我这样讲也似乎是合宜的。如果有人对任何事对我说明并指教，我必不争辩地接受他的指教。凡诵读圣经，无论是旧约或新约，以信服的心诵读的，必能学习，并能教导。但若他对自己不明白的事争论不休，他的心便不能接受教导。但若他发现某些话太难，不明白其中的深意，就让他这样说：\Quote{凡所写的都写得好，只是我尚未领悟其意义。} 如果他向那些探究教理、有智慧、有辨识力的人请教那些对他来说太困难的事，那么，当十个智者为同一件事以十种不同的方式向他讲述时，他要接受那使他喜欢的；若有不喜欢的地方，也不可轻视那些智者；因为天主的话如同珍珠，无论转向哪一面，都呈现美丽的景象。门徒啊，你要记得达味所说的话：\Quote{我从我所有的老师那里得到了学习。} 宗徒也说：\Quote{要阅读一切在圣神内的圣经。凡事要察验；善美的要持守，各种恶事要远离。} 因为即使一个人的年日，多如从亚当直到世界终穷的所有日子，且坐下默思圣经，他也不能完全领悟话语深奥的全部力量。人无法企及天主的智慧；正如我在第十篇讲论中所写的。然而，凡不从那大宝库中提取的讲者的话，都是可诅咒且应被轻视的。因为国王的肖像（在他钱币上的）无论到哪里都被接受；但含有杂质的（钱币）则被拒绝，不被接受。如果有人要说：\Quote{这些讲论是由某某人讲的；} 他当好好明白，刻意查究讲者是谁，并未受命给他。我，一个微不足道的人，也写了这些事，我是由亚当所出，由天主的手所造，是圣经的门徒。因为我们的主曾说：\BibleQuote{凡是求的，就必得到；找的，就必找到；敲的，就必给他开。} \BibleRef{玛 7:8} 先知也说：\BibleQuote{此后，我要将我的神倾注在一切有血肉的人身上，他们要说预言。} \BibleRef{岳 2:28} 因此，凡读到我以上所写的任何内容的，务要以信服的心去读，并要为他——作为基督身体的弟兄——祈祷；藉着天主的全教会的祈求，他的罪得以赦免。凡读的人要明白所记载的：\BibleQuote{学习真道的，应与施教的人分享自己的一切财物。} \BibleRef{迦 6:6} 又记载说：\BibleQuote{撒种的和收割的将一同喜欢。} \BibleRef{若 4:36} \BibleQuote{各人要按自己的劳苦领受自己的赏报。} \BibleRef{格前 3:8} \BibleQuote{没有隐藏的事，将来不被揭露的。} \BibleRef{玛 10:26}\par

\SourceRecord{Translated by John Gwynn.；Nicene and Post-Nicene Fathers, Second Series；Vol. 13.；Edited by Philip Schaff and Henry Wace.；Buffalo, NY: Christian Literature Publishing Co.,；1890.；<http://www.newadvent.org/fathers/370122.htm>.}
